% Preambulo del PLAN-CORRECCION (XeLaTeX). Pensado para imprimir en A4, a una cara o a doble cara.
\documentclass[10pt,a4paper]{article}
\usepackage[a4paper,top=2.3cm,bottom=2.4cm,left=2.2cm,right=2.2cm,headsep=0.7cm,footskip=1.0cm]{geometry}
\usepackage{fontspec}
\setmainfont{Palatino Linotype}[Ligatures=TeX,Numbers=Lining]
\setsansfont{Segoe UI}[Ligatures=TeX,Scale=0.92,BoldFont={Segoe UI Semibold}]
\setmonofont{Consolas}[Scale=0.88]
\newfontfamily\simbolos{Segoe UI Symbol}[Scale=0.9]
\newfontfamily\mates{Cambria Math}[Scale=0.95]
\usepackage[spanish,es-tabla,es-nodecimaldot]{babel}
\usepackage{xcolor}
\usepackage{array,booktabs,longtable,tabularx,multirow}
\usepackage{enumitem}
\usepackage{titlesec}
\usepackage{fancyhdr}
\usepackage{ragged2e}
\usepackage{needspace}
\usepackage[most]{tcolorbox}
\usepackage[hidelinks,bookmarksnumbered]{hyperref}
\hypersetup{pdftitle={Plan de corrección — Tesis EduFEM (main\_v4)},pdfauthor={Auditoría de redacción, 2026-09-20}}

% ---------- color: sobrio, imprime bien en escala de grises ----------
\definecolor{tinta}{HTML}{1B1F24}
\definecolor{gris}{HTML}{5A6470}
\definecolor{grisclaro}{HTML}{E9ECEF}
\definecolor{papel}{HTML}{F6F7F8}
\definecolor{critica}{HTML}{9B1C1C}
\definecolor{alta}{HTML}{B45309}
\definecolor{media}{HTML}{1D4E89}
\definecolor{baja}{HTML}{4B5563}
\definecolor{verde}{HTML}{1F6F43}
\color{tinta}

% ---------- parrafos ----------
\setlength{\parindent}{0pt}
\setlength{\parskip}{4.5pt plus 1pt minus 1pt}
\linespread{1.06}
\emergencystretch=2em
\tolerance=1200
\hyphenpenalty=600
\widowpenalty=10000
\clubpenalty=10000
\raggedbottom

% ---------- titulos ----------
\titleformat{\section}{\sffamily\bfseries\LARGE\color{tinta}}{\thesection}{0.7em}{}[\vspace{2pt}{\color{tinta}\titlerule[1.2pt]}]
\titleformat{\subsection}{\sffamily\bfseries\large\color{tinta}}{\thesubsection}{0.6em}{}
\titleformat{\subsubsection}{\sffamily\bfseries\normalsize\color{gris}}{\thesubsubsection}{0.6em}{}
\titlespacing*{\section}{0pt}{18pt plus 4pt}{10pt}
\titlespacing*{\subsection}{0pt}{14pt plus 3pt}{5pt}
\titlespacing*{\subsubsection}{0pt}{10pt plus 2pt}{3pt}

% ---------- cabecera y pie ----------
\pagestyle{fancy}
\fancyhf{}
\renewcommand{\headrulewidth}{0.4pt}
\renewcommand{\footrulewidth}{0pt}
\fancyhead[L]{\sffamily\footnotesize\color{gris}Plan de corrección · Tesis EduFEM (main\_v4.pdf)}
\fancyhead[R]{\sffamily\footnotesize\color{gris}\nouppercase{\leftmark}}
\fancyfoot[C]{\sffamily\footnotesize\color{gris}\thepage}
\renewcommand{\sectionmark}[1]{\markboth{#1}{}}

% ---------- listas ----------
\setlist{leftmargin=1.4em,itemsep=2pt,topsep=3pt,parsep=0pt}

% ---------- piezas ----------
\newcommand{\chip}[2]{{\sffamily\bfseries\footnotesize\color{#1}\MakeUppercase{#2}}}
\newcommand{\casilla}{{\simbolos\char"2610}\hspace{0.35em}}
\newcommand{\rotulo}[1]{{\sffamily\bfseries\footnotesize\color{gris}#1}\hspace{0.4em}}
\newcommand{\cita}[1]{{\itshape «#1»}}

% ficha de hallazgo: #1 color, #2 cabecera (ID, severidad, ubicacion, esfuerzo)
\newtcolorbox{ficha}[2]{enhanced,breakable,colback=white,colframe=#1,boxrule=0pt,leftrule=2.6pt,
  arc=0pt,outer arc=0pt,left=8pt,right=4pt,top=3pt,bottom=4pt,before skip=9pt,after skip=9pt,
  title={#2},coltitle=tinta,colbacktitle=papel,fonttitle=\sffamily\small,toptitle=3pt,bottomtitle=3pt,
  titlerule=0pt,pad at break=2mm,before upper={\setlength{\parskip}{3.5pt}\small}}

% par «dice / debe decir» de la seccion de calibracion
\newtcolorbox{dice}{enhanced,breakable,colback=papel,colframe=grisclaro,boxrule=0.4pt,arc=1.5pt,
  left=6pt,right=6pt,top=4pt,bottom=4pt,before skip=3pt,after skip=2pt,fontupper=\small}
\newtcolorbox{debedecir}{enhanced,breakable,colback=white,colframe=verde,boxrule=0.6pt,arc=1.5pt,
  left=6pt,right=6pt,top=4pt,bottom=4pt,before skip=2pt,after skip=8pt,fontupper=\small}
\newtcolorbox{recuadro}[1]{enhanced,breakable,colback=papel,colframe=tinta,boxrule=0.6pt,arc=1.5pt,
  left=8pt,right=8pt,top=6pt,bottom=6pt,before skip=8pt,after skip=10pt,title={#1},
  fonttitle=\sffamily\bfseries\small,coltitle=white,colbacktitle=tinta}

\newcolumntype{L}[1]{>{\RaggedRight\arraybackslash\small}p{#1}}
\newcolumntype{Y}{>{\RaggedRight\arraybackslash\small}X}

\begin{document}
\thispagestyle{empty}
{\sffamily\color{gris}\small Auditoría de redacción · 20 de septiembre de 2026}\par\vspace{4pt}
{\sffamily\bfseries\fontsize{26}{30}\selectfont Plan de corrección}\par\vspace{4pt}
{\sffamily\Large Tesis EduFEM — \texttt{tesis/main\_v4.pdf} (168 páginas)}\par\vspace{2pt}{\color{tinta}\rule{\linewidth}{1.2pt}}\vspace{10pt}
Este plan acompaña al diagnóstico de AUDITORIA-REDACCION.md y está pensado para trabajar con él al lado de la tesis impresa o del PDF. No propone reescribir el documento: propone 225 correcciones ubicadas, de las que 31 son críticas o altas. Cada una da la sección, el folio impreso y la página del PDF, la cita literal de lo que hoy dice el texto, por qué es un problema y qué hacer —qué eliminar, qué fusionar, cuál es la aparición canónica o el texto reformulado—.
\begin{recuadro}{Cómo usar este plan}
\begin{itemize}
\item Empieza por la sección 2, que fija el orden de las pasadas y las fórmulas que se deciden una sola vez. La sección 3 (calibración) va antes que los capítulos porque cambia frases que otras correcciones citan.
\item Cada corrección lleva una casilla para marcarla, su ID estable (el mismo que en AUDITORIA-REDACCION.md y en datos/\allowbreak{}hallazgos.json) y el tiempo estimado. «También CAL-04» en la cabecera de una ficha significa que otro eje señaló la misma frase: se corrige una vez.
\item Los textos propuestos son propuestas: respetan el molde del tribunal y no prometen efectos sobre estudiantes, pero la voz es la tuya. Donde la propuesta pide «verificar» es porque depende de lo que dice una fuente que la auditoría no abrió.
\item Las fichas traen la corrección completa de lo crítico y lo alto; lo medio y lo bajo va en tablas compactas al final de cada capítulo, con la propuesta resumida (la versión íntegra está en datos/\allowbreak{}hallazgos.json).
\item Al terminar, la sección 7 es una lista de control para comprobar, con búsquedas, que no quedó ninguna frase sin calibrar ni ningún criterio sin veredicto.
\end{itemize}\end{recuadro}
{\small\tableofcontents}
\clearpage
\section{Diagnóstico ejecutivo}
El documento está en buen estado y es defendible: la cadena del problema a las conclusiones se puede seguir, toda afirmación del Resumen tiene respaldo, las 25 referencias están citadas y en orden, y no hay remisiones rotas. La frontera delicada —qué se midió y qué no— está declarada muchas veces y bien. El problema no es que falten salvedades, sino dónde están: lejos de las afirmaciones fuertes, que son justo las que el tribunal lee primero.\par
Quedan 225 hallazgos: 4 críticos, 27 altos, 177 medios y 17 bajos. Un verificador independiente intentó refutar cada crítico o alto contra el documento: descartó 16 y ajustó 86, casi siempre a la baja.\par
\subsection*{Los cinco problemas de mayor impacto}
\begin{enumerate}[leftmargin=1.5em,itemsep=5pt]
\item \textbf{La finalidad sobre el estudiante se afirma justo donde no lleva salvedad.} Objetivo general («para mejorar el criterio»), primera conclusión («cumplió con el objetivo general»), pregunta del problema («bajo criterio» que nadie midió), quinto objetivo («potencial para formar criterio») y cierre de la Justificación no llevan salvedad contigua. Se fija una salvedad canónica y se pega a cada uno, o se reformulan como propiedad del artefacto. {\small\color{gris}\sffamily CAL-01, CAL-03, CAL-04, CAL-05, CAL-06, CAL-14 · 2 h}
\item \textbf{Los resultados no cierran, uno por uno, contra los criterios que el propio documento fijó.} 2.1.6 enuncia los criterios en una oración de 180 palabras; cinco de diecisiete tienen cifra pero no veredicto; la tasa de tensión no dice entre qué mallas se mide (la Tabla 3.1 muestra 1,21 frente a 1,4); 3.8 calla las componentes secundarias (2,9 \% y 4,6 \%) y las Conclusiones suben a «libre de errores de orden». Dos tablas lo arreglan. {\small\color{gris}\sffamily TRZ-07, TRZ-06, CAL-15, CAL-16, CAL-12 · 3 h}
\item \textbf{Hay tres cifras que el tribunal puede refutar con las tablas del propio documento.} «Hasta el 0,56 \% frente a SAP2000» en el Resumen, con 4,56 \% y 2,72 \% en la Tabla 3.2 (falta «en desplazamientos»); la razón Q4/\allowbreak{}Q9 de Cook «no depende de la referencia», y pasa de 13 a 7 veces al moverla de 23,96 a 23,97; 23,96 atribuido a «la práctica de la comunidad». {\small\color{gris}\sffamily TRZ-01, CAL-02, TRZ-04, CON-03 · 1,5 h}
\item \textbf{La validez de contenido es un cotejo del autor y a ratos se viste de validación externa.} El 18/\allowbreak{}18 y el 4/\allowbreak{}4 los construye y aplica el autor; se declara, pero en una subordinada. Filas cuyo «instrumento» es un resultado de la V\&V del autor; dos instrumentos del mismo autor como «principal argumento de validez»; «contenido universal del método» sin respaldo; «no existe un entorno» sin decir qué se buscó; población sin nombre de asignatura. {\small\color{gris}\sffamily CAL-10, CAL-13, CAL-08, CAL-07, TRZ-03, DEF-01, TRZ-05 · 4 h}
\item \textbf{El lector no especialista se pierde donde se mide la evidencia principal, y el documento se repite.} El Resumen da el resultado central en norma L2 y seminorma H1 sin decir qué miden; el Capítulo 1 entra por la forma variacional sin puente desde el análisis matricial; el índice se recita en prosa cinco veces, las limitaciones se escriben enteras tres, y hay 101 oraciones de 70 palabras o más. {\small\color{gris}\sffamily ACC-01, ACC-02, ACC-03, RED-02, RED-03, CLA-15 · 4 h + patrones}
\end{enumerate}
\begin{recuadro}{Esfuerzo total}Esfuerzo estimado: 14 h para todo lo crítico y alto (dos jornadas); 82 h para lo medio y 5 h para lo bajo si se hiciera todo, hallazgo por hallazgo y sin descontar que muchos se resuelven juntos.\end{recuadro}
\clearpage
\section{Orden de trabajo}
\begin{center}\small\begin{tabular}{@{}lrrrrrr@{}}\toprule
\sffamily\bfseries Eje & \sffamily\bfseries Crítica & \sffamily\bfseries Alta & \sffamily\bfseries Media & \sffamily\bfseries Baja & \sffamily\bfseries Total & \sffamily\bfseries Esfuerzo\\\midrule
Calibración & 3 & 13 & 53 & 2 & 71 & 23,2 h\\
Trazabilidad & 1 & 6 & 19 & 1 & 27 & 10,8 h\\
Riesgo de defensa & 0 & 2 & 18 & 0 & 20 & 10,3 h\\
Redundancia & 0 & 2 & 15 & 3 & 20 & 9,2 h\\
Claridad & 0 & 0 & 24 & 1 & 25 & 15,4 h\\
Accesibilidad & 0 & 3 & 10 & 0 & 13 & 8,8 h\\
Consistencia & 0 & 1 & 16 & 4 & 21 & 7,3 h\\
Aparato formal & 0 & 0 & 16 & 2 & 18 & 10,5 h\\
Referencias & 0 & 0 & 5 & 0 & 5 & 3,4 h\\
Tipografía y edición & 0 & 0 & 1 & 4 & 5 & 2,5 h\\
\midrule\textbf{Total} & \textbf{4} & \textbf{27} & \textbf{177} & \textbf{17} & \textbf{225} & \textbf{101,5 h}\\\bottomrule\end{tabular}\end{center}
Críticos y altos suman 14 horas; medios, 82 horas; bajos, 5 horas. Las estimaciones son por hallazgo y no descuentan que muchos se resuelven juntos (una decisión de nombre arregla decenas de apariciones), de modo que el total real de lo medio y lo bajo es sensiblemente menor. Una secuencia razonable: dos jornadas para las pasadas 1 a 7, que son las que pesan en la defensa; una semana a ritmo parcial para las pasadas 8 a 11.\par\medskip
El orden importa porque las correcciones se encadenan: la salvedad canónica que se fija en la primera pasada reaparece en una docena de frases, la tabla de criterios de 2.1.6 es la que luego recorre 3.8, y el Resumen se corrige al final porque refleja todo lo anterior. Dentro de cada pasada conviene seguir el orden de páginas.
{\small\begin{longtable}{@{}L{0.7cm}L{3.8cm}L{7.9cm}L{2.6cm}@{}}\toprule
\sffamily\bfseries N.º & \sffamily\bfseries Pasada & \sffamily\bfseries Qué se hace y por qué en este orden & \sffamily\bfseries Hallazgos · tiempo\\\midrule\endhead\bottomrule\endfoot
1 & \textbf{Decidir las fórmulas} & Antes de tocar el texto: la salvedad canónica, los tres registros de veredicto, el nombre único de la variable dependiente, del destinatario, del artefacto, de los casos y del validador, y el nombre, sigla y semestre de la asignatura que define la población (hoy una perífrasis, con un DATO PENDIENTE en el fuente). & TRZ-05, CON-06, CON-08, CON-09, CON-17, CON-18\newline 1 h\\\addlinespace[4pt]
2 & \textbf{Frases estructurales} & Pregunta del problema (y su copia en 2.1.3), objetivo general, quinto objetivo específico, cierre de la Justificación, definición de variables y Tabla 2.1 (validez de contenido y coherencia son propiedades del artefacto, no dimensiones de la variable dependiente; «manipular» anuncia un experimento que no hubo). & CAL-04, CAL-01, CAL-06, CAL-05, CAL-17, CAL-40, CAL-41\newline 2 h\\\addlinespace[4pt]
3 & \textbf{Criterios y veredictos} & Convertir la oración de 180 palabras de 2.1.6 en una tabla de criterios numerados con umbral y procedencia del umbral; precisar entre qué mallas se mide la tasa de tensión; cerrar 3.2–3.5 con «criterio, observado, se cumple»; sustituir la prosa de 3.8 por la tabla criterio–umbral–cifra–veredicto; declarar las componentes secundarias. & TRZ-06, TRZ-07, TRZ-10, TRZ-15, CAL-15\newline 3 h\\\addlinespace[4pt]
4 & \textbf{Cifras refutables} & «en desplazamientos» en el Resumen; la razón Q4/\allowbreak{}Q9 de Cook con su dependencia de la referencia; la procedencia de 23,96; la cota 0,3 \%/\allowbreak{}0,6 \% acotada a Timoshenko; «planificación» fuera de la frase de 3.7.2; incertidumbre 0,05 \% frente a umbral 0,1 \% en el mismo folio. & TRZ-01, CAL-02, TRZ-04, CAL-12, CON-03, CON-11, CON-15\newline 1,5 h\\\addlinespace[4pt]
5 & \textbf{Validez de contenido} & Separar en la Tabla 3.7 las filas cuyo instrumento es una función del software de las que remiten a un resultado del autor; rebajar la coincidencia de los dos instrumentos de «principal argumento de validez» a coherencia interna; retirar «contenido universal del método»; declarar la autoría de los instrumentos con frase propia y no en una subordinada; presentar los cuatro principios como síntesis propia; acotar «no existe un entorno» al conjunto revisado y decir qué se buscó; responder por qué valen 67 expertos mexicanos de ingeniería mecánica. & CAL-10, CAL-13, CAL-08, CAL-44, CAL-53, TRZ-03, CAL-07, TRZ-13, DEF-01, TRZ-09, TRZ-18\newline 4 h\\\addlinespace[4pt]
6 & \textbf{3.7–3.8, Conclusiones y Aportes} & El cierre de 3.8 («expone lo que el estudiante necesita»), la primera conclusión, «libre de errores de orden», «aporte pedagógico central», «atiende de forma directa al criterio», instrumentos «reutilizables», la última frase de la tesis; el párrafo que falta sobre el riesgo de que EduFEM se use como otra caja negra; el inventario de limitaciones escrito una sola vez. & CAL-14, CAL-03, CAL-16, CAL-11, CAL-60, CAL-61, CAL-65, CAL-67, DEF-06, RED-03, TRZ-25, TRZ-27\newline 3 h\\\addlinespace[4pt]
7 & \textbf{Resumen} & Al final, porque refleja todo lo anterior: L2 y H1 con una glosa de qué miden, «ocho módulos» frente a siete etapas más M0, quién estableció la validez de contenido, el modal que la Introducción sí lleva («puede aprender a operar sin formar criterio»), y la frase final con los verbos ya decididos. & ACC-01, TRZ-01, CAL-18, CAL-19, CAL-20\newline 1 h\\\addlinespace[4pt]
8 & \textbf{Accesibilidad de la teoría} & El puente desde el análisis matricial de pórticos al continuo al abrir 1.3; qué miden la norma L2, la seminorma H1 y la tasa de convergencia, en llano, en 1.13 y al abrir 3.2; glosas de media línea en la primera aparición de cada término; la Tabla 1.1 legible sin el texto. & ACC-02, ACC-03, ACC-04, ACC-07, ACC-09, ACC-11\newline 4 h\\\addlinespace[4pt]
9 & \textbf{Redundancia} & Dejar una sola recitación del índice; la clasificación metodológica, una vez (conservando en la Introducción la cláusula del nivel explicativo y su remisión a la prueba de campo); la salvedad de la prueba de campo, con oración completa solo donde es canónica y con remisión en el resto; la literatura de Lee, Bishay y Pérez-Santiago contada una vez con sus localizadores. & RED-02, TRZ-02, RED-04, RED-05, RED-09, RED-16\newline 4 h\\\addlinespace[4pt]
10 & \textbf{Patrones de estilo, nombres y aparato} & Por búsqueda: un nombre por concepto; oraciones de 70 palabras o más (reescribir las peores, no todas); verbos absolutos; impersonal que oculta al sujeto; figuras B.7–B.12 sin citar; las 18 ecuaciones numeradas que nadie cita; Anexos C y F citados con propósito; 22 citas que exigen localizador; años de cinco claves del .bib; Anexo E. & CLA-15, CAL-43, CLA-19, APF-04, TRZ-14, TRZ-12, REF-02, REF-04, CON-19\newline según alcance\\\addlinespace[4pt]
11 & \textbf{Recompilar y controlar} & Regenerar índices, comprobar que los folios 85 y 95 no queden casi vacíos, releer el Resumen y pasar la lista de control de la sección 7. & sección 7\newline 1 h\\\addlinespace[4pt]
\end{longtable}}
\subsection*{Fórmulas canónicas (decidir una vez, repetir igual)}
\begin{debedecir}\rotulo{Salvedad canónica (una sola, siempre igual)}«…ese efecto no se midió en este trabajo y queda como hipótesis pendiente de la prueba de campo diseñada en las recomendaciones.» Con oración completa en cuatro lugares —hipótesis, 2.1.6, 3.8 y cierre de Conclusiones— y, en el resto, pegada a la afirmación como inciso o sustituida por una remisión a esos lugares.\end{debedecir}
\begin{debedecir}\rotulo{Tres registros de veredicto}Medido por la V\&V: «se comprueba» o «se cumple el criterio fijado en la Subsección 2.1.6», acotado a los casos ensayados. Cotejo documental del autor: «el autor cotejó… contra el consenso publicado» o «cada destreza tiene un instrumento identificado por el autor». Literatura: «se fundamenta en lo que la literatura reporta de otras herramientas; no hay evidencia sobre EduFEM». «Queda confirmada» vale para una cláusula solo junto a «frente a los criterios fijados a priori».\end{debedecir}
\begin{debedecir}\rotulo{Propiedad del artefacto en lugar de efecto}Donde el texto dice que EduFEM «permite formar», «hace comprender», «ancla conceptos» o «atiende al criterio», el sujeto pasa a ser el software y el verbo, lo que hace: «expone», «muestra», «permite inspeccionar», «permite reproducir a mano», «pone a la vista». El estudiante solo es sujeto de lo que puede hacer con el programa, no de lo que aprende.\end{debedecir}
\clearpage
\section{Calibración de afirmaciones}
Es el eje prioritario. El trabajo verificó un motor y cotejó contenido y coherencia de diseño; no midió criterio, aprendizaje ni comprensión de ningún estudiante. Esta sección lista, en el orden del documento, cada frase crítica o alta que atribuye a EduFEM un efecto sobre el estudiante, transfiere sin salvedad lo que la literatura reporta de otras herramientas, usa un verbo absoluto, reclama primacía sin acotar o generaliza desde tres casos. Para cada una: lo que dice hoy, por qué es un problema, el tipo de afirmación (medida, fundamentada en literatura o supuesta por diseño) y la orden de eliminarla o su versión calibrada. El hallazgo de fondo del verificador independiente ordena el trabajo: en este documento las salvedades abundan, pero están lejos de las afirmaciones fuertes; los dos lugares donde de verdad no hay ninguna son el objetivo general y el cierre del párrafo de literatura de la Justificación. Las frases que ya están bien calibradas y sirven de modelo —Limitaciones, folio 91; el cierre de 1.2, folio 20; la hipótesis «por etapas», folios 6 y 7— no aparecen aquí.
\subsection{Introducción}
\needspace{6\baselineskip}
\noindent\casilla\textbf{CAL-04}\quad\chip{alta}{alta}\quad Introducción · Planteamiento del problema · folio 2 · PDF 17\quad{\color{gris}\footnotesize también DEF-02}\hfill 25 min\par\vspace{1pt}
{\small\color{gris}Calibración · afirmación supuesta}\par
\begin{dice}\rotulo{Dice}«¿de qué manera el uso del software de análisis estructural como caja negra podrá afectar en el bajo criterio para interpretar la respuesta estructural obtenida por el MEF en los estudiantes de la Carrera de Ingeniería Civil...?»\end{dice}
{\small\rotulo{Por qué}La pregunta científica da por establecido que el criterio de esa población es «bajo» —premisa que el trabajo no midió en ningún estudiante— y formula además un efecto causal que esta etapa no contrasta. El régimen es incorrecto: afectar es transitivo. La misma formulación se repite textualmente en la matriz de consistencia, de modo que el defecto viaja al instrumento que el tribunal usa para verificar la coherencia interna.\par}
\begin{debedecir}\rotulo{Debe decir / qué hacer}«¿de qué manera la forma en que el software expone el canal de cálculo y la respuesta estructural —como caja negra o de manera transparente, interactiva y numéricamente verificable— podrá incidir en el criterio para interpretar la respuesta estructural obtenida por el MEF en los estudiantes de la Carrera de Ingeniería Civil de la Universidad Autónoma Tomás Frías en la gestión 2026?» Y a continuación, fijando la frontera: «El diagnóstico de que ese criterio resulta insuficiente procede de la literatura y de los antecedentes revisados en la sección 1.2, y no de una medición realizada en esa población; ese efecto no se midió en este trabajo y queda como hipótesis pendiente de la prueba de campo diseñada en las recomendaciones.» El localizador de la fuente que sostiene el diagnóstico debe ser el que ya se cita en el planteamiento del problema; no introducir una referencia nueva sin verificarla. Replicar la formulación, palabra por palabra, en la subsección 2.1.3 (folio 41, PDF 56), que la retoma textualmente.\end{debedecir}
{\small\rotulo{Misma corrección en}§2.1.3 Matriz de consistencia, f. 41.\par}
{\small\color{gris}\rotulo{DEF-02}Pregunta previsible del tribunal de metodología: «usted parte de que estos estudiantes tienen bajo criterio, ¿de dónde sale ese diagnóstico y lo midió aquí?». El problema da por establecida una condición de la población local que ningún dato propio sostiene —la evidencia que el documento ofrece es un consenso de expertos mexicanos y una encuesta profesional—, y el párrafo tampoco declara qué parte de la interrogante responde este trabajo y cuál queda abierta, justo donde el tribunal busca la frontera entre lo comprobado y lo pendiente.\par}
\vspace{6pt}
\needspace{6\baselineskip}
\noindent\casilla\textbf{CAL-05}\quad\chip{alta}{alta}\quad Introducción · Justificación · folio 4 · PDF 19\hfill 15 min\par\vspace{1pt}
{\small\color{gris}Calibración · afirmación supuesta}\par
\begin{dice}\rotulo{Dice}«Exponer el mapeo isoparamétrico, el Jacobiano, la matriz de deformación y la cuadratura de Gauss sobre el elemento real que el estudiante construyó, y mostrar la respuesta que producen, permite anclar conceptos abstractos en una geometría tangible y juzgar el resultado con conocimiento de cómo se obtuvo.»\end{dice}
{\small\rotulo{Por qué}La frase cierra el párrafo de literatura trasladando a EduFEM, en indicativo y sin salvedad, el efecto cognitivo que las tres fuentes citadas (Lee, Bishay, Rutten) reportan de otras herramientas. «Anclar conceptos» y «juzgar el resultado con conocimiento» son efectos sobre el estudiante que este trabajo no midió, enunciados como propiedad ya verificada del artefacto.\par}
\begin{debedecir}\rotulo{Debe decir / qué hacer}«Sobre esa base, EduFEM expone el mapeo isoparamétrico, el Jacobiano, la matriz de deformación y la cuadratura de Gauss sobre el elemento real que el estudiante construyó, junto con la respuesta que producen: es la decisión de diseño con que este trabajo traslada esos resultados a la elasticidad plana. Que ese anclaje mejore la interpretación de los resultados es la conjetura que orienta el diseño: ese efecto no se midió en este trabajo y queda como hipótesis pendiente de la prueba de campo diseñada en las recomendaciones.»\end{debedecir}
\vspace{6pt}
\needspace{6\baselineskip}
\noindent\casilla\textbf{CAL-01}\quad\chip{critica}{crítica}\quad Introducción · Objetivos (objetivo general) · folio 5 · PDF 20\quad{\color{gris}\footnotesize también DEF-03}\hfill 20 min\par\vspace{1pt}
{\small\color{gris}Calibración · afirmación supuesta}\par
\begin{dice}\rotulo{Dice}«... para mejorar el criterio de los estudiantes de ingeniería civil en la interpretación de la respuesta estructural obtenida por el MEF.»\end{dice}
{\small\rotulo{Por qué}El objetivo general cierra con un efecto sobre estudiantes que el trabajo no midió y, a diferencia de la hipótesis, no lleva ninguna salvedad ni remisión a la prueba de campo. Al quedar así, todo «objetivo cumplido» posterior —en 3.8 y en Conclusiones— se lee como afirmación de efecto logrado. Es el punto que el tribunal leerá primero y del que colgará la pregunta de si se midió algo en estudiantes.\par}
\begin{debedecir}\rotulo{Debe decir / qué hacer}«... que integre la formulación matemática, la construcción y visualización interactiva del modelo y la verificación numérica, y que exponga de manera transparente, interactiva y numéricamente verificable el canal de cálculo y la respuesta estructural sobre el modelo del propio estudiante, como medio didáctico orientado a la formación del criterio de los estudiantes de ingeniería civil para interpretar esa respuesta.» Inmediatamente después de la lista de objetivos específicos: «La finalidad formativa orienta el diseño y no es un resultado de esta etapa: ese efecto no se midió en este trabajo y queda como hipótesis pendiente de la prueba de campo diseñada en las recomendaciones.» Si el molde del tribunal obliga a conservar la fórmula «para mejorar el criterio», dejarla y añadir a continuación esa misma salvedad, con las mismas palabras.\end{debedecir}
{\small\color{gris}\rotulo{DEF-03}Pregunta de apertura casi segura: «¿a cuántos estudiantes midió?». Los dos enunciados que el tribunal lee primero y en voz alta prometen un efecto y un potencial sobre el estudiante que el trabajo no midió y que el resto del documento acota con cuidado; el contraste entre la promesa del objetivo y la salvedad de 3.8 es lo que convierte la pregunta en una contradicción visible.\par}
\vspace{6pt}
\needspace{6\baselineskip}
\noindent\casilla\textbf{CAL-06}\quad\chip{alta}{alta}\quad Introducción · Objetivos (objetivo específico 5) · folio 6 · PDF 21\hfill 15 min\par\vspace{1pt}
{\small\color{gris}Calibración · afirmación supuesta}\par
\begin{dice}\rotulo{Dice}«... para establecer su potencial para formar criterio.»\end{dice}
{\small\rotulo{Por qué}Una tabla de especificaciones y una matriz de principios establecen cobertura de contenido y coherencia de diseño, no un «potencial para formar criterio»: el objetivo promete una conclusión que el instrumento empleado no puede sostener, que el capítulo 3 dará por cumplida y que las propias Limitaciones desmienten («validez de contenido no es efectividad»). El eslabón objetivo-conclusión-limitación queda contradictorio.\par}
\begin{debedecir}\rotulo{Debe decir / qué hacer}Objetivo específico 5: «Validar el contenido del software frente a las destrezas de interpretación y verificación de la respuesta que los expertos exigen al egresado, mediante una tabla de especificaciones, y su coherencia con los principios de aprendizaje que fundamentan el diseño, para establecer la cobertura de contenido del artefacto y la coherencia de su diseño, condición previa al efecto formativo: ese efecto no se midió en este trabajo y queda como hipótesis pendiente de la prueba de campo diseñada en las recomendaciones.» Suprimir «para establecer su potencial para formar criterio». Comprobar que la fila correspondiente de la Tabla 2.2 (matriz de consistencia) y la conclusión del folio 89 empleen el mismo enunciado, sin la finalidad formativa.\end{debedecir}
{\small\rotulo{Misma corrección en}Limitaciones, f. 91.\par}
\vspace{6pt}
\subsection{Capítulo 1 · 1.1–1.2 Estado del arte y fundamentos pedagógicos}
\needspace{6\baselineskip}
\noindent\casilla\textbf{CAL-07}\quad\chip{alta}{alta}\quad §1.1 Antecedentes y estado del arte (síntesis, cierre del OE1) · folio 17 · PDF 32\quad{\color{gris}\footnotesize también DEF-04}\hfill 30 min\par\vspace{1pt}
{\small\color{gris}Calibración · afirmación fundamentada}\par
\begin{dice}\rotulo{Dice}«no existe un entorno que exponga de forma transparente, interactiva y numéricamente verificable ese canal completo en elasticidad plana sobre el modelo concreto del alumno»\end{dice}
{\small\rotulo{Por qué}El diagnóstico que sostiene el problema científico y cierra el primer objetivo específico se enuncia como inexistencia universal, cuando lo que una revisión acotada puede sostener es que no se halló tal entorno entre los recursos que identificó y examinó. Agrava el punto que la sección nombra seis categorías de recurso y nueve productos pero no declara en ninguna parte el alcance de la búsqueda: fuentes consultadas, términos y periodo.\par}
\begin{debedecir}\rotulo{Debe decir / qué hacer}Se mantiene la reformulación propuesta -«... y, entre los recursos que esta revisión identificó y pudo examinar (Tabla 1.1), no se halló un entorno que exponga de forma transparente, interactiva y numéricamente verificable ese canal completo en elasticidad plana sobre el modelo concreto del alumno.»- con tres precisiones: (a) la forma canónica modelo está en el folio 4 (PDF 19), no en el folio 6; (b) el texto remite a que «el problema científico de la introducción» formula esa misma brecha, de modo que la acotación debe aplicarse también allí, previa verificación de la redacción exacta del problema; (c) la exigencia de declarar bases de datos, términos de búsqueda y periodo solo procede si el autor puede documentar esa búsqueda; si no la registró, basta con declarar una sola vez, al abrir 1.1, cuántos recursos se revisaron y con qué criterio de inclusión, sin inventar un protocolo de revisión sistemática que no se ejecutó.\end{debedecir}
{\small\rotulo{Misma corrección en}Introducción · Justificación (forma ya acotada, sirve de modelo), f. 6.\par}
{\small\color{gris}\rotulo{DEF-04}Pregunta previsible: «usted afirma que no existe un entorno con estas características; ¿qué buscó, dónde y con qué criterio?». La afirmación pasa de «ninguno de los revisados» a «no existe» sin que el documento declare en ningún punto la estrategia de la revisión —bases consultadas, términos, ventana temporal, criterios de inclusión—, de modo que basta que un miembro del tribunal nombre un recurso no considerado para derribar la premisa de oportunidad del trabajo.\par}
\vspace{6pt}
\subsection{Capítulo 2 · 2.1 Diseño metodológico}
\needspace{6\baselineskip}
\noindent\casilla\textbf{CAL-08}\quad\chip{alta}{alta}\quad §2.1.5 Instrumentos · folio 45 · PDF 60\hfill 20 min\par\vspace{1pt}
{\small\color{gris}Calibración · afirmación supuesta}\par
\begin{dice}\rotulo{Dice}«... la limitación de ese criterio —expertos mexicanos de ingeniería mecánica, que el propio estudio declara no generalizables a otros contextos— se acota porque los ítems que se toman son contenido universal del método y no preferencias de un país»\end{dice}
{\small\rotulo{Por qué}La limitación que la única fuente externa declara de sí misma se neutraliza con una afirmación del propio autor —que esos ítems son contenido universal del método— presentada como hecho y sin respaldo, justo en el punto donde el trabajo es más vulnerable por carecer de panel propio. Es un supuesto razonable, pero enunciado como dato invierte la carga de la prueba delante del tribunal.\par}
\begin{debedecir}\rotulo{Debe decir / qué hacer}«El criterio externo es el consenso publicado de expertos, que hace las veces de juicio de expertos sin panel propio. Ese consenso declara no generalizables sus resultados a otros contextos [3, p. 1171]\allowbreak{} -lo recogieron expertos mexicanos de ingeniería mecánica-, y este trabajo lo adopta bajo un supuesto que no verifica con un panel propio: que los ítems seleccionados corresponden a la formulación del método y a la interpretación de sus resultados, contenido común a cualquier curso introductorio de MEF, y no a preferencias de una industria o de un país. Contrastar ese supuesto con docentes de la carrera queda pendiente, junto con la prueba de campo que diseñan las recomendaciones.»\end{debedecir}
\vspace{6pt}
\subsection{Capítulo 2 · 2.2 El software EduFEM}
\needspace{6\baselineskip}
\noindent\casilla\textbf{CAL-09}\quad\chip{alta}{alta}\quad §2.2.6 Preproceso y lienzo · folio 56 · PDF 71\hfill 20 min\par\vspace{1pt}
{\small\color{gris}Calibración · afirmación supuesta}\par
\begin{dice}\rotulo{Dice}«la orientación de cada cuadrilátero se fuerza a sentido antihorario calculando el área con signo, lo que garantiza un determinante Jacobiano positivo»\end{dice}
{\small\rotulo{Por qué}La afirmación es falsa como enunciado general: en un cuadrilátero no convexo el determinante jacobiano puede ser negativo en algún punto de Gauss aunque los vértices estén ordenados en sentido antihorario. Además contradice a la propia subsección, que dos párrafos antes describe un validador que rechaza elementos cuyo determinante no supera el umbral mínimo y enumera el jacobiano negativo entre las advertencias. Es el único punto de la zona donde el texto afirma algo técnicamente incorrecto.\par}
\begin{debedecir}\rotulo{Debe decir / qué hacer}«Como en el modo de dibujo, la orientación de cada cuadrilátero se fuerza a sentido antihorario calculando el área con signo. Esa orientación es condición necesaria, pero no suficiente: en un cuadrilátero no convexo el determinante Jacobiano puede ser negativo en algún punto de Gauss aun con los vértices bien ordenados, de modo que la comprobación definitiva corre por cuenta del validador de la Subsección 2.2.5, que evalúa det J en los puntos de Gauss y rechaza el elemento si no supera el umbral mínimo.»\end{debedecir}
\vspace{6pt}
\subsection{Capítulo 3 · 3.1–3.5 Verificación y validación numérica}
\needspace{6\baselineskip}
\noindent\casilla\textbf{CAL-02}\quad\chip{critica}{crítica}\quad §3.5 Membrana de Cook · folio 72 · PDF 87\quad{\color{gris}\footnotesize también CON-02}\hfill 30 min\par\vspace{1pt}
{\small\color{gris}Calibración · afirmación medida}\par
\begin{dice}\rotulo{Dice}«el error del Q9 es trece veces menor que el del Q4; esa razón no depende de la incertidumbre de la referencia, que solo afecta a la tercera cifra del error del Q9»\end{dice}
{\small\rotulo{Por qué}La razón se calcula a 2178 GDL con un error del Q9 de -0,044 \%, que la nota de la misma sección declara no significativo por estar por debajo del 0,1 \%. Con la referencia de 23,97 que esa nota contempla, los errores pasan a -0,634 \% (Q4) y -0,088 \% (Q9) y la razón baja a unas siete veces; a 578 GDL, en cambio, ambos errores superan la incertidumbre y la razón queda entre doce y quince. La afirmación «no depende de la incertidumbre de la referencia» contradice, por tanto, los datos de la Tabla 3.5 del propio documento, y la cifra de trece se repite en 3.7.1.\par}
\begin{debedecir}\rotulo{Debe decir / qué hacer}«A igualdad de grados de libertad, el error del Q9 es más de diez veces menor que el del Q4: con 578 GDL (Q4 con N = 16 y Q9 con N = 8) la razón está entre doce y quince veces según se tome como referencia 23,96 o 23,97; con 2178 GDL el error del Q9 (-0,044 \%) cae ya por debajo del 0,1 \% y la razón deja de ser significativa.» Eliminar «esa razón no depende de la incertidumbre de la referencia...». En 3.7.1 cambiar «trece veces menos error con el Q9» por «más de diez veces menos error»; el «un orden de magnitud» del Resumen se sostiene y puede quedar.\end{debedecir}
{\small\rotulo{Misma corrección en}§3.5, nota de la referencia de Cook, f. 71; §3.7.1 Alcances, f. 80; Resumen, f..\par}
{\small\color{gris}\rotulo{CON-02}La afirmación es aritméticamente falsa y contradice la nota al pie de la página anterior. Con los valores de la propia Tabla 3.4 y u\_\allowbreak{}ref = 23,96 los errores son -0,594 \% (Q4, N=32) y -0,044 \% (Q9, N=16), razón 12,9; con el u\_\allowbreak{}ref = 23,97 que la misma nota respalda por extrapolación de Richardson pasan a -0,634 \% y -0,088 \%, razón 7,2: el error del Q9 cambia en su segunda cifra significativa, no en la tercera, y la razón entre errores depende justamente de lo que la frase declara indiferente. El mismo «trece veces» se repite sin salvedad en 3.7.1 (folio 80), de modo que la corrección debe propagarse. Nota de ubicación: la tabla de convergencia de Cook es la Tabla 3.4; la Tabla 3.5 es la síntesis comparativa Q4/\allowbreak{}Q9.\par}
\vspace{6pt}
\subsection{Capítulo 3 · 3.6–3.8 Validez de contenido, interpretación e hipótesis}
\needspace{6\baselineskip}
\noindent\casilla\textbf{CAL-10}\quad\chip{alta}{alta}\quad §3.6.2, Tabla 3.7 (tabla de especificaciones) · folio 76 · PDF 91\hfill 60 min\par\vspace{1pt}
{\small\color{gris}Calibración · afirmación mixta}\par
\begin{dice}\rotulo{Dice}«FEA30 prueba de convergencia refinando hasta una diferencia dada | Secuencias MMS y Cook; extrapolación de Richardson | Sección 3.5»\end{dice}
{\small\rotulo{Por qué}En varias filas (FEA30 y FEA34 entre ellas) la columna «instrumento de EduFEM» no es una función de la aplicación sino un resultado de la verificación que hizo el autor, y la evidencia remite solo al Capítulo 3, contra la regla de marcado de la Subsección 2.1.5, que exige que el instrumento exista en el software distribuido y el Capítulo 2 lo describa. El MMS y la extrapolación de Richardson no están entre lo que la aplicación ofrece, de modo que el «dieciocho de dieciocho» que sostiene la cláusula (c), el Resumen y las conclusiones queda sobrecalibrado. El texto ya declara parte del límite («secuencias que el estudiante debe construir, sin un generador automático»), pero no distingue interfaz de guion.\par}
\begin{debedecir}\rotulo{Debe decir / qué hacer}Tarea de verificación, no de reescritura a ciegas. En FEA30 y FEA34, dejar en la columna «instrumento de EduFEM» únicamente funciones que el software distribuido ofrezca y que el Capítulo 2 describa (mallado y refinamiento sucesivo, conversión entre tipos de elemento, lectura de desplazamientos, tensiones y reacciones en el post-proceso y en la memoria de cálculo), y trasladar «secuencias MMS y Cook; extrapolación de Richardson» y «contraste con la solución analítica y con SAP2000» a la columna de evidencia o a una nota que los identifique como procedimientos ejecutados por el autor con los guiones publicados junto al código. Añadir en la columna de evidencia el localizador del Capítulo 2 que describe cada instrumento, como en el resto de la tabla. Reportar el recuento desglosado: cuántas entradas quedan cubiertas por funciones de la interfaz y cuántas por procedimientos del repositorio. Si alguna fila queda solo con el guion, declararla cubierta «por el repositorio, no por la interfaz» o dejar la celda vacía, como exige la Subsección 2.1.5. Conservar la nota ya existente sobre la ausencia de generador automático de secuencias de malla.\end{debedecir}
{\small\rotulo{Misma corrección en}§3.6.2, fila FEA34, f. 77; §2.1.5 (regla de marcado), f. 45.\par}
\vspace{6pt}
\needspace{6\baselineskip}
\noindent\casilla\textbf{CAL-11}\quad\chip{alta}{alta}\quad §3.7.1 Alcances · folio 80 · PDF 95\hfill 15 min\par\vspace{1pt}
{\small\color{gris}Calibración · afirmación supuesta}\par
\begin{dice}\rotulo{Dice}«es el argumento pedagógico central, porque permite al estudiante constatar de primera mano la ventaja del elemento de mayor orden frente al bloqueo, un fenómeno que de otro modo permanece abstracto.»\end{dice}
{\small\rotulo{Por qué}Se atribuye a la comparación un efecto sobre el estudiante —que constata de primera mano y que sin ella el fenómeno le resultaría abstracto— que ninguna observación de este trabajo respalda. Es la formulación más explícita del deslizamiento en todo el capítulo de interpretación.\par}
\begin{debedecir}\rotulo{Debe decir / qué hacer}«...y la comparación a igualdad de grados de libertad —más de diez veces menos error con el Q9— es el resultado de mayor interés didáctico declarado en el diseño, porque pone la ventaja del elemento de mayor orden frente al bloqueo en una cifra reproducible sobre un mismo modelo, en lugar de enunciarla como consecuencia teórica.» Suprimir «permite al estudiante constatar de primera mano» y «un fenómeno que de otro modo permanece abstracto». La cifra se coordina con CALb-09.\end{debedecir}
\vspace{6pt}
\needspace{6\baselineskip}
\noindent\casilla\textbf{CAL-12}\quad\chip{alta}{alta}\quad §3.7.1 Alcances · folio 80 · PDF 95\hfill 20 min\par\vspace{1pt}
{\small\color{gris}Calibración · afirmación medida}\par
\begin{dice}\rotulo{Dice}«Para los problemas de referencia estudiados, el error en las magnitudes primarias —tensión normal y flecha— se mantiene por debajo del 0,3 \% frente a la solución analítica y del 0,6 \% frente al modelo de SAP2000»\end{dice}
{\small\rotulo{Por qué}La cota se generaliza en plural a «los problemas de referencia estudiados» cuando solo la viga de Timoshenko tiene solución analítica cerrada y modelo equivalente en SAP2000: la membrana de Cook no tiene solución analítica y su Q4 conserva un error de -0,594 \% con 2178 GDL, y el MMS no produce errores porcentuales frente a una referencia de ese tipo. La misma generalización se repite en Aportes.\par}
\begin{debedecir}\rotulo{Debe decir / qué hacer}«En la viga de Timoshenko, el único caso de esta batería con solución analítica cerrada y con modelo equivalente en SAP2000, el error en las magnitudes primarias —tensión normal y flecha— se mantiene por debajo del 0,3 \% frente a la solución analítica y del 0,6 \% frente al modelo de SAP2000.» Repetir la misma acotación en Aportes.\end{debedecir}
{\small\rotulo{Misma corrección en}Aportes del trabajo, f. 90.\par}
\vspace{6pt}
\needspace{6\baselineskip}
\noindent\casilla\textbf{CAL-13}\quad\chip{alta}{alta}\quad §3.7.2 Validez de contenido y coherencia · folio 81 · PDF 96\hfill 45 min\par\vspace{1pt}
{\small\color{gris}Calibración · afirmación supuesta}\par
\begin{dice}\rotulo{Dice}«Esa coincidencia es el principal argumento de validez de los dos instrumentos, porque se construyeron por vías distintas —el contenido que los expertos exigen y los principios que la literatura documenta— y no se calibraron entre sí.»\end{dice}
{\small\rotulo{Por qué}El argumento central de validez descansa en la convergencia de dos instrumentos construidos por el mismo autor, sobre el mismo artefacto y con el mismo conocimiento previo. Además, la independencia no se cumple para la matriz de principios: los cuatro principios gobiernan las decisiones de diseño desde la Sección 1.2, de modo que la matriz solo podía arrojar ese resultado y comprueba coherencia interna, no concordancia independiente. El límite se declara en 3.1 y en 2.1.5, pero no donde se enuncia el argumento de validez, que es donde el tribunal lo leerá.\par}
\begin{debedecir}\rotulo{Debe decir / qué hacer}Sustituir la frase por: «Los dos instrumentos no tienen el mismo estatuto y conviene distinguirlos. La tabla de especificaciones coteja el software con un criterio de contenido publicado y ajeno al proyecto, y es la que aporta el argumento de validez de contenido. La matriz de principios recorre, en cambio, los mismos principios que guiaron el diseño desde la Sección 1.2: comprueba la coherencia interna entre fundamento y decisión —que ninguna decisión quede sin principio y ningún principio sin encarnación—, no una concordancia independiente. Que ambas señalen el mismo conjunto de decisiones es consistente con esa lectura, no una corroboración cruzada: las construyó el mismo autor sobre el mismo artefacto, tal como declara la Subsección 2.1.5. Su comprobación externa queda para la prueba de campo de las recomendaciones.» Corregir además, en el folio 80, «La matriz de principios confirma desde la literatura lo que la tabla muestra desde el contenido» por «La matriz de principios coincide, desde la literatura, con lo que la tabla muestra desde el contenido».\end{debedecir}
{\small\rotulo{Misma corrección en}§3.7.2, f. 80; §3.1 (limite ya declarado), f. 62.\par}
\vspace{6pt}
\needspace{6\baselineskip}
\noindent\casilla\textbf{CAL-14}\quad\chip{alta}{alta}\quad §3.8 Validación de la hipótesis · folio 84 · PDF 99\hfill 20 min\par\vspace{1pt}
{\small\color{gris}Calibración · afirmación supuesta}\par
\begin{dice}\rotulo{Dice}«cubre las destrezas de interpretación y verificación que los expertos exigen, expone lo que el estudiante necesita para formar criterio sobre la respuesta estructural.»\end{dice}
{\small\rotulo{Por qué}La frase que cierra la respuesta al problema científico afirma una suficiencia sobre el aprendizaje —que lo expuesto es lo que el estudiante necesita para formar criterio— que ni se midió ni se estableció en esos términos. La salvedad de la oración siguiente cubre solo el efecto («que ese criterio mejore de hecho»), no la suficiencia de lo expuesto, de modo que la frase citable, en el lugar más visible del capítulo, queda sin acotar; el propio 3.7.2 declara que ninguno de los dos instrumentos establece el efecto.\par}
\begin{debedecir}\rotulo{Debe decir / qué hacer}Sustituir el cierre y la oración que le sigue por un solo período: «...y cubre las destrezas de interpretación y verificación que el consenso publicado de expertos declara necesarias para el egresado, expone el canal de cálculo completo, el procedimiento que produce la respuesta y los medios para verificarla sobre el modelo del propio usuario. Si esa exposición forma el criterio del estudiante —el efecto que el problema pregunta en última instancia— es la parte de la hipótesis que esta etapa no contrasta: se fundamenta por analogía con las herramientas de la misma familia que sí lo midieron u observaron (Sección 1.2) y la contrastará la prueba de campo diseñada en las recomendaciones.»\end{debedecir}
\vspace{6pt}
\needspace{6\baselineskip}
\noindent\casilla\textbf{CAL-15}\quad\chip{alta}{alta}\quad §3.8 Validación de la hipótesis (cláusula b) · folio 84 · PDF 99\hfill 20 min\par\vspace{1pt}
{\small\color{gris}Calibración · afirmación medida}\par
\begin{dice}\rotulo{Dice}«de hasta el 0,21 \% en sigma\_\allowbreak{}x y el 0,56 \% en desplazamientos frente a SAP2000 (Tabla 3.2, Tabla 3.3), dentro de los umbrales del 1 \% y del 3 \% fijados en la Subsección 2.1.6»\end{dice}
{\small\rotulo{Por qué}El veredicto de la cláusula (b) cita solo las magnitudes primarias y no advierte que el criterio de 2.1.6 excluye expresamente las componentes secundarias de tensión, que alcanzan el 2,9 \% frente a la solución analítica y el 4,6 \% en sigma\_\allowbreak{}y y el 2,7 \% en tau\_\allowbreak{}xy frente a SAP2000 (folio 79). Quien lea únicamente 3.8 concluirá que todos los errores medidos quedan bajo el 3 \%. No hay incumplimiento del criterio, sí omisión en el veredicto.\par}
\begin{debedecir}\rotulo{Debe decir / qué hacer}Verificar primero el alcance del criterio en la Subsección 2.1.6. Si los umbrales del 1 \% y del 3 \% están fijados sobre las magnitudes primarias, añadir tras «fijados en la Subsección 2.1.6»: «—umbrales fijados sobre las magnitudes primarias; las componentes secundarias de tensión, de magnitud entre seis y treinta y cinco veces menor que sigma\_\allowbreak{}x y recuperadas con un orden de convergencia inferior, alcanzan el 2,9 \% frente a la solución analítica y el 4,6 \% frente a SAP2000, según detalla la Subsección 3.7.1—». Si 2.1.6 no distingue entre magnitudes primarias y secundarias, corregir allí el criterio para que lo haga y decir en 3.8 que el veredicto de la cláusula (b) se emite sobre las magnitudes primarias. En cualquiera de los dos casos, que el lector de 3.8 no pueda concluir que todos los errores medidos quedan bajo el 3 \%. La armonización de «queda confirmada» con «queda comprobada» del Resumen se trata como corrección de coherencia aparte.\end{debedecir}
{\small\rotulo{Misma corrección en}§3.7.1 Alcances, f. 79.\par}
\vspace{6pt}
\subsection{Conclusiones y recomendaciones}
\needspace{6\baselineskip}
\noindent\casilla\textbf{CAL-03}\quad\chip{critica}{crítica}\quad Conclusiones (apertura) · folio 86 · PDF 101\hfill 25 min\par\vspace{1pt}
{\small\color{gris}Calibración · afirmación mixta}\par
\begin{dice}\rotulo{Dice}«El presente trabajo cumplió con el objetivo general planteado: se desarrolló EduFEM, un software educativo de escritorio escrito en el lenguaje de programación Python sobre bibliotecas numéricas de código abierto»\end{dice}
{\small\rotulo{Por qué}La primera frase de las Conclusiones declara cumplido el objetivo general y lo parafrasea sin su cláusula finalista —«para mejorar el criterio de los estudiantes de ingeniería civil en la interpretación de la respuesta estructural»—, con lo que un efecto no medido queda dado por logrado en el lugar más visible del documento, sin salvedad y sin remisión a la prueba de campo.\par}
\begin{debedecir}\rotulo{Debe decir / qué hacer}«El presente trabajo cumplió el objetivo general en su componente de artefacto: se desarrolló EduFEM, un software educativo de escritorio escrito en el lenguaje de programación Python sobre bibliotecas numéricas de código abierto [sigue igual hasta «la verificación numérica de los resultados»]\allowbreak{}. La finalidad con que ese objetivo se enunció —mejorar el criterio del estudiante para interpretar la respuesta estructural— no se midió en esta etapa y no se da por lograda: es la hipótesis pendiente que deberá contrastar la prueba de campo diseñada en las recomendaciones.»\end{debedecir}
{\small\rotulo{Misma corrección en}Introducción, Objetivos (objetivo general), f. 5.\par}
\vspace{6pt}
\needspace{6\baselineskip}
\noindent\casilla\textbf{CAL-16}\quad\chip{alta}{alta}\quad Conclusiones (verificación y validación del motor) · folio 88 · PDF 103\quad{\color{gris}\footnotesize también DEF-26}\hfill 15 min\par\vspace{1pt}
{\small\color{gris}Calibración · afirmación medida}\par
\begin{dice}\rotulo{Dice}«esas tasas son evidencia rigurosa de que el motor está libre de errores de orden.»\end{dice}
{\small\rotulo{Por qué}Las Conclusiones convierten en ausencia demostrada de errores lo que el cuerpo formula con alcance acotado: 3.2 (folio 64) enumera los mecanismos que las cuatro configuraciones del MMS ejercitan. Además, el propio 3.7.1 narra un defecto real —la matriz de extrapolación Eq4— que las normas de desplazamiento no delataron, de modo que «evidencia rigurosa de que el motor está libre de errores de orden» contradice al mismo documento y es justamente el tipo de afirmación que un tribunal con formación numérica cuestionará.\par}
\begin{debedecir}\rotulo{Debe decir / qué hacer}Sustituir por: «...esas tasas comprueban que, en las cuatro configuraciones ensayadas, no hay errores de orden en el ensamblaje, la integración por cuadratura, el mapeo isoparamétrico, la matriz constitutiva de ambos estados planos ni la imposición de restricciones; la corrección del operador de extrapolación se comprueba aparte, con la reproducción exacta de los campos polinómicos de cada elemento. Fuera de esos mecanismos y de esas configuraciones, las tasas no acreditan nada.»\end{debedecir}
{\small\rotulo{Misma corrección en}§3.2 Verificación por soluciones manufacturadas, f. 64; §3.7.1 Alcances, f. 79.\par}
{\small\color{gris}\rotulo{DEF-26}Pregunta del miembro de métodos numéricos: «usted cuenta que encontró un error en la matriz de extrapolación que su batería de convergencia no detectó y luego concluye que el motor está libre de errores de orden; ¿en qué quedamos?». La afirmación absoluta de las Conclusiones contradice el relato del propio trabajo dos páginas antes y permite poner en duda toda la verificación; el relato del defecto es un punto fuerte que el verbo absoluto convierte en munición.\par}
\vspace{6pt}
\subsection{Calibración de severidad media y baja, y patrones de léxico}
{\small\begin{longtable}{@{}L{1.6cm}L{1.6cm}L{5.25cm}L{5.65cm}r@{}}
\toprule
\sffamily\bfseries ID & \sffamily\bfseries Dónde & \sffamily\bfseries Qué pasa & \sffamily\bfseries Qué hacer & \sffamily\bfseries min\\
\midrule\endhead
\bottomrule\endfoot
\mbox{\casilla CAL-17}\newline {\color{media}\sffamily\bfseries\scriptsize MEDIA} & f. 39-40\newline{\scriptsize Calibración} & La validez de contenido y la coherencia son propiedades del software, pero la tabla las cuelga como dimensiones de la variable dependiente —el criterio del estudiante—, de modo que un lector rápido concluye que la dependiente se midió en dos de sus tres dimensiones. Esa operacionalización es la que habilita después decir que la dependiente «se observa» y que la hipótesis se comprueba «en parte», cuando de la dependiente no se observó nada: es el punto por donde el tribunal puede impugnar todo el contraste. & Cuerpo (folio 39): cambiar solo la atribución, sin duplicar la salvedad que ya está: «La variable dependiente es el criterio del estudiante para interpretar y juzgar la respuesta estructural obtenida por el MEF: desplazamientos, tensiones, reacciones, calidad de la malla y fenómenos numéricos. Esta etapa no la observa, porque… & 45\\\addlinespace[3pt]
\mbox{\casilla CAL-18}\newline {\color{media}\sffamily\bfseries\scriptsize MEDIA} & f. i\newline{\scriptsize Calibración} & La pasiva impersonal más el verbo de establecimiento hacen leer que la validez la establecieron los 67 expertos, cuando el cotejo lo construyó y lo aplicó el propio autor sobre un consenso ya publicado, sin panel de jueces convocado para esta tesis ni índice de acuerdo; la única salvedad («las construye el autor») vive en 2.1.5, folio 46, y en Limitaciones, lejos de las cinco apariciones del verbo fuerte. Es la afirmación metodológica central del trabajo y está sobrecalibrada justo donde más lectores se detienen. & Resumen: «Para establecer la validez de contenido, el autor construyó una tabla de especificaciones que coteja los instrumentos del software con las destrezas de interpretación y verificación de resultados que un consenso publicado de 67 expertos exige al egresado: las quince destrezas y los tres conceptos del universo tienen… & 30\\\addlinespace[3pt]
\mbox{\casilla CAL-19}\newline {\color{media}\sffamily\bfseries\scriptsize MEDIA} & f. i\newline{\scriptsize Calibración} & El Resumen afirma en indicativo, como hecho observado, una consecuencia sobre el estudiante que la Introducción enuncia con modal y apoya en la literatura, no en una observación propia del autor. & «... de modo que, según advierte la literatura, el estudiante puede aprender a operar la herramienta sin formar criterio para interpretar y juzgar lo que obtiene.» & 5\\\addlinespace[3pt]
\mbox{\casilla CAL-20}\newline {\color{media}\sffamily\bfseries\scriptsize MEDIA} & f. ii\newline{\scriptsize Calibración} & La frase acota bien el alcance, pero «su efecto ... se fundamenta en la literatura» presupone que el efecto existe y que sólo falta medirlo, cuando lo fundamentado es la expectativa; y no dice que esa literatura trata de otras herramientas. Un lector rápido retiene «la hipótesis queda comprobada». & «La parte de diseño y contenido de la hipótesis se cumple. La parte de efecto —que esa exposición mejore el criterio del estudiante— no se midió en este trabajo: se apoya, por analogía, en lo que la literatura reporta de otras herramientas didácticas del método y queda como hipótesis pendiente de la prueba de campo diseñada en… & 10\\\addlinespace[3pt]
\mbox{\casilla CAL-21}\newline {\color{media}\sffamily\bfseries\scriptsize MEDIA} & f. 2\newline{\scriptsize Calibración} & Un estudio de consenso de opinión no «confirma» y no habla de «criterio»: reporta que la formación privilegia la operación sobre los conceptos. El constructo criterio sobre la respuesta es del autor —1.2 lo define como tal— y aquí se atribuye a la fuente. En el mismo párrafo, «Esta opacidad tiene una consecuencia pedagógica concreta» enuncia como hecho una inferencia. & «Un estudio de consenso entre expertos de la industria y de la academia reporta que la formación tiende a privilegiar las destrezas de operación sobre los conceptos fundamentales y reclama que el egresado sepa planificar, verificar y validar sus análisis, no solo ejecutarlos [3, pp. 1159, 1171]\allowbreak{}: es lo que este trabajo denomina… & 10\\\addlinespace[3pt]
\mbox{\casilla CAL-22}\newline {\color{media}\sffamily\bfseries\scriptsize MEDIA} & f. 3\newline{\scriptsize Calibración} & La Justificación presenta como resultados firmes lo que las secciones 1.2 y 2.2.1 acotan después —evaluación cualitativa sin grupo de control en la primera fuente; mejora entre cohortes atribuida sobre todo a la autoría de los alumnos en la segunda—, de modo que la literatura llega inflada al primer lugar donde el tribunal la lee y sostiene la pertinencia del trabajo con una fuerza que las propias secciones posteriores contradicen. & «La literatura sobre enseñanza del método reporta, con evidencia de alcance desigual que la sección 1.2 detalla, que la simulación interactiva y la manipulación directa de los objetos matemáticos del MEF conducen a una comprensión más intuitiva de sus procedimientos —evaluación cualitativa, sin grupo de control— [9, p. 157; 10,… & 15\\\addlinespace[3pt]
\mbox{\casilla CAL-23}\newline {\color{media}\sffamily\bfseries\scriptsize MEDIA} & f. 4\newline{\scriptsize Calibración} & «Cada operación» es absoluto: la memoria de cálculo desarrolla el procedimiento sobre el modelo del usuario, pero no toda operación de un modelo arbitrario es reproducible a mano. En el mismo párrafo se afirma sin evidencia que la organización en tres fases hace el flujo «pedagógicamente legible», que es de nuevo un efecto sobre el lector. & «... el estudiante puede rehacer a mano, o en una hoja de cálculo, las operaciones que la memoria desarrolla —la matriz de deformación, la cuadratura de Gauss, la matriz de rigidez elemental y su ensamblaje— y contrastarlas con las del programa, elemento por elemento.» Y en el plano práctico: «... lo que reproduce el flujo de… & 15\\\addlinespace[3pt]
\mbox{\casilla CAL-24}\newline {\color{media}\sffamily\bfseries\scriptsize MEDIA} & f. 5\newline{\scriptsize Calibración} & La finalidad del cuarto objetivo es absoluta cuando lo que tres casos de referencia permiten es acotar el error frente a ellos; la misma fórmula se repite en 2.1.1. Generalizar desde tres casos a la corrección de toda respuesta es el tipo de salto que el tribunal puede señalar sin ser especialista. & «... para acotar, frente a referencias de exactitud conocida, el error de la respuesta que el software expone.» Aplicar la misma redacción en 2.1.1. & 10\\\addlinespace[3pt]
\mbox{\casilla CAL-25}\newline {\color{media}\sffamily\bfseries\scriptsize MEDIA} & f. 7\newline{\scriptsize Calibración} & «Ese efecto» es la mejora del criterio para interpretar la respuesta, pero según la propia sección 1.2 las fuentes reportan otra cosa: comprensión más intuitiva en una evaluación cualitativa sin control, y mejora entre cohortes atribuida sobre todo a la autoría de los alumnos. «Herramientas de la misma familia» no está definido y la transferencia a EduFEM es una inferencia por analogía que el texto no declara como tal: la literatura queda avalando un efecto que ninguna de esas fuentes midió. & Adoptar como forma canónica la del folio 20 -«hacen plausible el efecto que la hipótesis postula... no constituyen evidencia de eficacia de EduFEM, que este trabajo no mide»- y corregir solo la atribución, en lugar de insertar un párrafo largo en cuatro lugares (duplicaría lo que el cierre de 1.2 ya dice). Folio 7: «... es la… & 30\\\addlinespace[3pt]
\mbox{\casilla CAL-26}\newline {\color{media}\sffamily\bfseries\scriptsize MEDIA} & f. 7\newline{\scriptsize Calibración} & Al tipificar la investigación se afirma que el artefacto «resuelve» un problema formativo, que es justamente lo que queda por probar; la misma fórmula se repite en 2.1.1. Es un solo verbo, pero aparece en el lugar donde se declara qué tipo de trabajo es y qué se propone demostrar. & «... la propuesta y construcción de un artefacto (EduFEM) que se ofrece como solución tentativa a un problema formativo concreto.» En 2.1.1: «Por su finalidad, el trabajo es de carácter propositivo: propone un artefacto dirigido a un problema formativo, cuya eficacia formativa no se midió en este trabajo y queda como hipótesis… & 10\\\addlinespace[3pt]
\mbox{\casilla CAL-27}\newline {\color{media}\sffamily\bfseries\scriptsize MEDIA} & f. 9\newline{\scriptsize Calibración} & Afirma una transferencia de aprendizaje del estudiante que no se midió; lo generalizable es la formulación matemática, no la comprensión. En el mismo apartado, «directamente generalizables» refuerza el absoluto y carece de localizador de página. & «El dominio plano actúa así como un puente: la formulación isoparamétrica, la cuadratura de Gauss y el ensamblaje que el software expone en dos dimensiones son los mismos que rigen el caso tridimensional [ref., p. xx]\allowbreak{}, planteado como línea de continuación del proyecto.» Sustituir «directamente generalizables» por… & 15\\\addlinespace[3pt]
\mbox{\casilla CAL-28}\newline {\color{media}\sffamily\bfseries\scriptsize MEDIA} & f. 12\newline{\scriptsize Calibración} & Justifica el uso de ediciones con más de diez años mediante una afirmación absoluta sobre el contenido de ediciones que no se consultaron y que ninguna fuente respalda. Es exactamente la pregunta que el tribunal formulará sobre la actualidad de la bibliografía, y tal como está la respuesta es indefendible: el autor no puede saber qué reproducen ediciones que no leyó. & «Las obras de referencia del método —Zienkiewicz y Taylor, Bathe, Hughes, Cook y colaboradores, Strang y Fix— se citan en las ediciones consultadas, anteriores en más de diez años a este trabajo, por su carácter de fuentes clásicas: la formulación isoparamétrica, la cuadratura de Gauss, el tratamiento del bloqueo y la teoría de… & 10\\\addlinespace[3pt]
\mbox{\casilla CAL-29}\newline {\color{media}\sffamily\bfseries\scriptsize MEDIA} & f. 14\newline{\scriptsize Calibración} & Superlativo sobre un universo que no se contó ni se delimitó; la revisión no declara cuántas publicaciones se examinaron ni con qué criterio, de modo que la primacía no es verificable. & «VisualFEA es el producto comercial con orientación didáctica declarada que con más frecuencia aparece en las publicaciones consultadas» o, si no se llevó un recuento, suprimir el superlativo. & 5\\\addlinespace[3pt]
\mbox{\casilla CAL-30}\newline {\color{media}\sffamily\bfseries\scriptsize MEDIA} & f. 15\newline{\scriptsize Calibración} & El «vacío» se define por la conjunción de nueve atributos que son los de EduFEM —incluidas filas como memoria de cálculo automática, fenómenos numéricos observables y V\&V publicada—, de modo que ningún recurso ajeno puede satisfacerla por construcción; además la afirmación no se acota al conjunto revisado. La tabla sirve para delimitar el aporte, no para medir una carencia del campo, y presentarla como carencia es circular. & Folio 15: conservar íntegra la enumeración de atributos y cambiar solo el encuadre: «Entre los recursos revisados, ninguno reúne de forma simultánea los atributos que este trabajo persigue: una interfaz íntegramente en español, el carácter libre y de código abierto, el dominio del continuo elástico plano, la exposición… & 40\\\addlinespace[3pt]
\mbox{\casilla CAL-31}\newline {\color{media}\sffamily\bfseries\scriptsize MEDIA} & f. 18\newline{\scriptsize Calibración} & Transfiere sin ninguna salvedad un tamaño de efecto medido en simulaciones de ciencias experimentales, en educación general, al objeto de esta tesis, equiparando por decisión del autor las matrices del MEF con los fenómenos físicos no observables de aquellos estudios. La analogía es legítima como orientación de diseño; enunciada como equivalencia, hace que el primer principio pedagógico llegue al tribunal como pronóstico respaldado. & «... entre los resultados que destaca está el gran efecto de aprendizaje de visualizar fenómenos que de otro modo son invisibles [12, p. 151]\allowbreak{}. Este trabajo asume, por analogía, que las matrices intermedias y los fenómenos numéricos ocupan dentro del canal de cálculo una posición equivalente a la de aquellos fenómenos no… & 20\\\addlinespace[3pt]
\mbox{\casilla CAL-32}\newline {\color{media}\sffamily\bfseries\scriptsize MEDIA} & f. 18\newline{\scriptsize Calibración} & «Mejoran la instrucción tradicional» es ambiguo —no se sabe si la complementan o la sustituyen— y «evidencia sólida» es una valoración que conviene atribuir con la formulación de la fuente, porque de esa frase cuelga el primer principio pedagógico. La auditoría no abrió la fuente: el cotejo queda a cargo del autor. & «... la revisión de una década de investigación cuasiexperimental sobre simulaciones por computadora en la enseñanza de las ciencias concluye que su uso como complemento de la instrucción tradicional produce mejoras de aprendizaje, en especial cuando refuerza o sustituye actividades de laboratorio, y que el efecto depende de… & 15\\\addlinespace[3pt]
\mbox{\casilla CAL-33}\newline {\color{media}\sffamily\bfseries\scriptsize MEDIA} & f. 18\newline{\scriptsize Calibración} & El mismo resultado de una fuente sobre otra herramienta se describe de tres maneras de alcance creciente -procesamiento de las ecuaciones, procedimientos del MEF, conceptos del análisis estructural-, y las dos últimas amplían lo que la versión más cauta atribuye a la fuente. El lector que coteje la cita no sabrá cuál es el alcance real del hallazgo ajeno que el trabajo invoca como apoyo. & Tomar como canónica la versión de 1.2 -«una comprensión más intuitiva del procesamiento de las ecuaciones»-, con su salvedad de evaluación cualitativa sin grupo de control, y repetirla sin variantes en la Justificación y en 2.2.1, comprobando en la página 168 de la fuente cuál es el alcance literal. & 15\\\addlinespace[3pt]
\mbox{\casilla CAL-34}\newline {\color{media}\sffamily\bfseries\scriptsize MEDIA} & f. 19\newline{\scriptsize Calibración} & Presenta como predicción teórica un efecto sobre el estudiante que la tesis no midió, convirtiendo un marco formulado para conductas de aprendizaje en general —no para el MEF— en aval anticipado del resultado de EduFEM. Además da por supuesto que el estudiante adoptará el modo constructivo al usar el programa, que es justamente lo que no se observó. La salvedad de cierre de 1.2 llega después y no cubre esta frase. & «Juzgar una respuesta -decidir si la malla es admisible, si las tensiones cierran, si el elemento bloquea- puede leerse, en los términos de ese marco, como una conducta constructiva. De ahí la conjetura de diseño de este trabajo: las capas superpuestas de EduFEM, el conmutador entre la fórmula y su valor numérico y la… & 20\\\addlinespace[3pt]
\mbox{\casilla CAL-35}\newline {\color{media}\sffamily\bfseries\scriptsize MEDIA} & f. 28\newline{\scriptsize Calibración} & La subordinada final atribuye a una decisión de diseño un efecto en el estudiante —que observe y capte el bloqueo— que el trabajo no midió. Es la única frase del capítulo 1 que enuncia un efecto sobre el estudiante como finalidad ya lograda. & «y conserva el comportamiento del Q4 sin corregirlo, de modo que el bloqueo quede visible en los resultados que el programa presenta en lugar de quedar compensado dentro del cálculo. Si esa exposición se traduce en un mejor criterio del estudiante no se midió en este trabajo y queda como hipótesis pendiente de la prueba de… & 15\\\addlinespace[3pt]
\mbox{\casilla CAL-36}\newline {\color{media}\sffamily\bfseries\scriptsize MEDIA} & f. 32\newline{\scriptsize Calibración} & Una prueba de regresión comprueba los casos que ejercita, no «cualquier» campo, y «exactamente» no distingue la identidad algebraica de la igualdad numérica que la prueba puede observar. & «una prueba de regresión comprueba que E reproduce los campos bilineales (o bicuadráticos) de prueba con un error del orden de la precisión de máquina, que es la verificación numérica de la identidad algebraica de la que E se deduce». & 10\\\addlinespace[3pt]
\mbox{\casilla CAL-37}\newline {\color{media}\sffamily\bfseries\scriptsize MEDIA} & f. 33\newline{\scriptsize Calibración} & La subordinada da por hecho que el alumno deriva la fórmula a mano, actividad que el trabajo no observó, cuando lo comprobable es que el módulo presenta esa forma de la métrica. & «... que es la forma de la métrica que el módulo presenta paso a paso, calculable a mano con el ángulo de cada esquina». & 10\\\addlinespace[3pt]
\mbox{\casilla CAL-38}\newline {\color{media}\sffamily\bfseries\scriptsize MEDIA} & f. 34\newline{\scriptsize Calibración} & La afirmación se emite en sentido pleno y la salvedad que la restringe llega dos oraciones después, de modo que la frase, citada suelta, atribuye al trabajo una validación experimental que no hizo. & «Este trabajo aplica ambas estrategias en la acepción que se precisa a continuación: verificación en sentido estricto y validación entendida como contraste con referencias externas de exactitud conocida, no como campaña experimental.» Así la salvedad deja de ser una corrección posterior y forma parte de la afirmación. & 15\\\addlinespace[3pt]
\mbox{\casilla CAL-39}\newline {\color{media}\sffamily\bfseries\scriptsize MEDIA} & f. 35\newline{\scriptsize Calibración} & Verbo absoluto: observar la tasa teórica de convergencia es la evidencia más fuerte disponible de que no hay errores que degraden el orden, pero no confirma la corrección del código, y el propio apartado admite dos oraciones después que las normas de desplazamiento no ejercitan la recuperación de tensiones. & «Una tasa observada inferior a la teórica delata un defecto de implementación; alcanzarla es la evidencia más exigente de que las rutinas ejercitadas por esa medida están correctamente implementadas, sin que ello descarte errores en las rutas que la medida no recorre.» & 15\\\addlinespace[3pt]
\mbox{\casilla CAL-40}\newline {\color{media}\sffamily\bfseries\scriptsize MEDIA} & f. 38\newline{\scriptsize Calibración} & De las variables del problema sólo la independiente puede observarse en el artefacto; la frase incluye también a la dependiente, aunque el texto deja a continuación el efecto a la prueba de campo. Es la formulación derivada del defecto estructural de CALa-06. & «La unidad de análisis de esta etapa es el artefacto: en él se observan la variable independiente y las condiciones de contenido y coherencia, con los instrumentos de la sección 2.1.5; la variable dependiente y el efecto que la hipótesis postula sobre el estudiante no se midieron en este trabajo y quedan como hipótesis… & 10\\\addlinespace[3pt]
\mbox{\casilla CAL-41}\newline {\color{media}\sffamily\bfseries\scriptsize MEDIA} & f. 39\newline{\scriptsize Calibración} & El verbo «manipular» anuncia una manipulación experimental que no ocurrió: la variable independiente se realiza en un solo nivel —la exposición transparente— y el nivel opuesto, el uso como caja negra, no se implementa ni se contrasta contra nada en este trabajo. Revisada toda la unidad, ningún apartado evalúa esa condición. El tribunal, al leer «se manipula», esperará un diseño con dos condiciones y no lo encontrará. & «La variable interviniente es EduFEM, la solución tentativa: el artefacto en el que la variable independiente toma el valor exposición transparente, interactiva y numéricamente verificable. Esta etapa no realiza el valor opuesto -el uso como caja negra-: el contraste con él es documental y se apoya en el análisis comparativo de… & 25\\\addlinespace[3pt]
\mbox{\casilla CAL-42}\newline {\color{media}\sffamily\bfseries\scriptsize MEDIA} & f. 41\newline{\scriptsize Calibración} & La salvedad distingue lo comprobable de lo no comprobable pero no dice dónde queda la parte no comprobable, de modo que el enunciado de la hipótesis en la matriz de consistencia sigue leyéndose como si la mejora del criterio se contrastara aquí. La remisión a la prueba de campo aparece recién en 2.1.4 y 2.1.6. & «[...]\allowbreak{} permitirá mejorar ese criterio. De esa hipótesis, esta etapa sólo contrasta las tres cláusulas relativas al artefacto —exposición, corrección y validez de contenido—; la cláusula de efecto sobre el criterio del estudiante no se midió en este trabajo y queda como hipótesis pendiente de la prueba de campo diseñada en las… & 15\\\addlinespace[3pt]
\mbox{\casilla CAL-43}\newline {\color{media}\sffamily\bfseries\scriptsize MEDIA} & f. 45\newline{\scriptsize Calibración} & Patrón en toda la zona: la familia garantiz- aparece 11 veces y confirm- 13 en el documento (recuento por script sobre el texto completo); una parte sustancial califica lo que sólo se sostiene por diseño, por la existencia de un antecedente o por tres casos de referencia. En el caso principal, de compartir funciones se sigue que los números coinciden, no que haya «coherencia» entre la explicación y el cálculo; y en 1.2, «la decisión que lo encarna y la evidencia que lo demuestra» convierte un registro de decisiones de diseño en demostración. No todas las apariciones son defecto: «garantizada por construcción» en el MMS es legítima. & Reservar «se cumple» y «garantiza» para los criterios numéricos con umbral fijado a priori y para lo que se sigue por construcción (diciéndolo: «por construcción»); en los demás casos usar «respalda», «es coherente con», «reporta» o «de modo que». Redacción canónica para el caso principal: «El sistema se ensambla y resuelve con… & 45\\\addlinespace[3pt]
\mbox{\casilla CAL-44}\newline {\color{media}\sffamily\bfseries\scriptsize MEDIA} & f. 46\newline{\scriptsize Calibración} & La limitación capital del trabajo —los dos instrumentos con que se contrasta la cláusula (c) los construye y los aplica el propio autor, sin jueces independientes— se despacha en una subordinada final de ocho palabras y sin enunciar su consecuencia. Es la primera objeción previsible del tribunal y la que hay que desactivar de frente; además es la salvedad que falta junto a las afirmaciones fuertes de CALa-01. & Reemplazar la subordinada final -no añadir párrafo nuevo, para no repetir la oración anterior-: «... de modo que cualquier lector puede rehacer cada veredicto con el software distribuido y el manual. Ambos instrumentos, sin embargo, los construye y los aplica el autor, sin un panel de jueces independiente: sus veredictos no son… & 25\\\addlinespace[3pt]
\mbox{\casilla CAL-70}\newline {\color{baja}\sffamily\bfseries\scriptsize BAJA} & f. 47\newline{\scriptsize Calibración} & En cuatro salvedades, por lo demás correctas, el artículo definido o el posesivo presuponen que el efecto existe y que sólo falta medirlo. La propia tesis tiene ya la fórmula justa en el cierre de 1.2: «el efecto que la hipótesis postula». & Emplear en las cuatro salvedades la fórmula del cierre de 1.2, «el efecto que la hipótesis postula», o «su eventual efecto»; en 3.8: «y con el efecto que la hipótesis postula pendiente de medición». & 10\\\addlinespace[3pt]
\mbox{\casilla CAL-45}\newline {\color{media}\sffamily\bfseries\scriptsize MEDIA} & f. 48\newline{\scriptsize Calibración} & El texto reconoce con honestidad que el efecto reportado por la segunda fuente se atribuye a que el alumno construye sus propias herramientas de cálculo, pero acto seguido invoca esa evidencia para justificar un diseño en el que el estudiante no programa nada, sin decir qué parte del hallazgo se considera transferible y por qué. La inferencia queda abierta justo donde el tribunal puede cerrarla en contra. & Añadir tras la salvedad: «EduFEM no pide al estudiante que escriba código, de modo que de ese segundo hallazgo sólo se retiene la parte que puede trasladarse a un programa de uso: que el procedimiento quede a la vista paso a paso y pueda reproducirse a mano. Si ese traslado conserva el efecto reportado no se midió en este… & 20\\\addlinespace[3pt]
\mbox{\casilla CAL-46}\newline {\color{media}\sffamily\bfseries\scriptsize MEDIA} & f. 55\newline{\scriptsize Calibración} & Atribuye al artefacto un efecto sobre la percepción y la experiencia del usuario que no se midió, con un verbo absoluto («evita») en lugar de describir lo que el programa hace. & «Con esta política, en mallas densas el lienzo dibuja sólo la silueta del dominio y la selección activa, mientras que en los modelos pequeños habituales en clase —que rara vez exceden unas decenas de elementos— el dibujo se mantiene completo en todo momento.» & 15\\\addlinespace[3pt]
\mbox{\casilla CAL-47}\newline {\color{media}\sffamily\bfseries\scriptsize MEDIA} & f. 57\newline{\scriptsize Calibración} & La fórmula «la interactividad como motor del aprendizaje» da por establecido un papel causal que la propia tesis califica diez folios antes como evaluación cualitativa sin grupo de control, y al hacerlo desplaza hacia EduFEM un efecto que este trabajo no midió. El resto de la frase, en cambio, separa bien lo que es decisión del autor. & «El diseño de estos módulos parte de lo que Lee reporta sobre la interactividad en sus simuladores del MEF [pp. 158, 166]\allowbreak{} —una evaluación cualitativa, sin grupo de control (sección 1.2)— y lo lleva a una filosofía minimalista que es decisión de este trabajo, no un resultado tomado de la literatura.» & 15\\\addlinespace[3pt]
\mbox{\casilla CAL-48}\newline {\color{media}\sffamily\bfseries\scriptsize MEDIA} & f. 58\newline{\scriptsize Calibración} & «Irresoluble a mano» es un absoluto impreciso —el integrando admite primitiva en principio, aunque impracticable para un cuadrilátero distorsionado— y choca con el argumento que la tesis sostiene en la Justificación y en la memoria de cálculo, según el cual el procedimiento que el programa muestra puede reproducirse a mano. & «contrapone la integral exacta, sin primitiva practicable para un elemento de geometría arbitraria, con su aproximación por cuadratura», añadiendo media línea sobre por qué importa: es esa impracticabilidad la que obliga a la cuadratura de Gauss y la que el módulo hace visible. & 15\\\addlinespace[3pt]
\mbox{\casilla CAL-49}\newline {\color{media}\sffamily\bfseries\scriptsize MEDIA} & f. 59\newline{\scriptsize Calibración} & La decisión de usar la paleta jet se justifica con dos afirmaciones que el texto da por sabidas y no respalda: que esa paleta es la práctica habitual de los programas comerciales —sin nombrar ninguno ni citar fuente— y que el estudiante encontrará después esa herramienta profesional, trayectoria que este trabajo no verificó. El pasaje no lleva ninguna llamada de cita. & «la paleta con la que programas de uso extendido en el medio profesional, como SAP2000 y ANSYS, presentan sus mapas de tensiones [ref., p. x]\allowbreak{}» y, para la segunda: «se privilegió la continuidad visual con las herramientas que el estudiante encontrará previsiblemente en la práctica profesional, supuesto que este trabajo no… & 25\\\addlinespace[3pt]
\mbox{\casilla CAL-50}\newline {\color{media}\sffamily\bfseries\scriptsize MEDIA} & f. 62\newline{\scriptsize Calibración} & «Garantiza», «valida» y «confirma» se emplean indistintamente para lo que se sigue de una demostración o de la construcción y para lo que una prueba puntual, con una tolerancia concreta, solo hace verosímil; así el verbo pierde valor donde sí corresponde (determinante jacobiano positivo, término fuente deducido de la solución exacta). Cuatro apariciones verificadas en el Capítulo 3 y el Anexo D, además de las que se reportan por separado por estar en secciones de veredicto (CALb-05 y CALb-13). La repetibilidad, en particular, depende de mantener el entorno —versiones de NumPy y SciPy, precisión de punto flotante—, cosa que el propio Anexo D reconoce en otro pasaje. & Fijar un verbo por tipo de evidencia y usarlo siempre igual, reservando «garantiza» para lo que se sigue de una demostración o de la construcción (determinante jacobiano positivo, término fuente deducido de la solución exacta, reproducción exacta de campos polinómicos) y «valida» para el contraste con referencias externas, en… & 40\\\addlinespace[3pt]
\mbox{\casilla CAL-51}\newline {\color{media}\sffamily\bfseries\scriptsize MEDIA} & f. 66\newline{\scriptsize Calibración} & La equivalencia con «los paquetes profesionales» se afirma sin nombrar ninguno, sin fuente y sin haber inspeccionado su implementación, que además es cerrada. & Sustituir el inciso por una formulación verificable: «adoptado, que separa el identificador público del nodo del índice ordinal en las matrices del sistema, tal como recomienda la literatura de implementación [cita con página]\allowbreak{}». Si no hay fuente a mano, suprimir el inciso: la prueba se sostiene sola. El verbo se corrige en… & 20\\\addlinespace[3pt]
\mbox{\casilla CAL-52}\newline {\color{media}\sffamily\bfseries\scriptsize MEDIA} & f. 73\newline{\scriptsize Calibración} & La variable dependiente es el criterio del estudiante y no puede observarse «en el artefacto»: lo que se observa es la cobertura del contenido que ese criterio exigiría. La confusión aparece en la frase que abre la sección donde se reportan las cláusulas (a) y (c), es decir, donde el tribunal fija qué se midió. & «Primero los indicadores de la variable independiente —que expone el software y cómo—, y después los dos instrumentos documentales con que se comprueba, sobre el artefacto, que su contenido y su diseño corresponden a lo que la variable dependiente exige; el criterio del estudiante mismo se mide en la prueba de campo diseñada en… & 10\\\addlinespace[3pt]
\mbox{\casilla CAL-53}\newline {\color{media}\sffamily\bfseries\scriptsize MEDIA} & f. 77\newline{\scriptsize Calibración} & El resultado de dieciocho sobre dieciocho —que sostiene la cláusula (c), el Resumen y las conclusiones— lo produjo el propio autor cotejando su propio artefacto, sin jueces independientes, y sobre un universo que no es el consenso completo sino el subconjunto de ítems que cae dentro del alcance declarado de EduFEM (la nota de la misma tabla enumera los excluidos: análisis tridimensional, placas y cáscaras, térmico, barras y vigas, CAD, simetría y contacto y los tipos de análisis más allá del estático lineal). El texto atribuye al autor solo la traducción de los códigos, nunca el juicio de correspondencia, de modo que un cien por ciento autoevaluado se presenta con el lenguaje de un criterio externo cumplido. & Sustituir la frase por: «Las quince destrezas y los tres conceptos tienen instrumento identificable en el software distribuido, con lo que queda cotejado el criterio de validez de contenido fijado en la Subsección 2.1.6. Conviene recordar aquí su condición, ya declarada en la Subsección 2.1.5: el cotejo entre cada ítem del… & 35\\\addlinespace[3pt]
\mbox{\casilla CAL-54}\newline {\color{media}\sffamily\bfseries\scriptsize MEDIA} & f. 79\newline{\scriptsize Calibración} & La frase que abre la interpretación de los resultados generaliza desde tres casos sin acotar el dominio; la salvedad correspondiente —que la validación se apoya en tres casos cuidadosamente seleccionados— aparece varias páginas después, al final de 3.7.4, de modo que el lector no especialista se lleva la idea de una fiabilidad general del motor. & «La batería de verificación y validación —un problema manufacturado, una solución analítica cerrada y un caso de referencia histórico— respalda la fiabilidad del núcleo numérico dentro del dominio que esos tres casos cubren: elasticidad lineal plana, elementos Q4 y Q9, y los tamaños de malla y tipos de carga ensayados.» & 15\\\addlinespace[3pt]
\mbox{\casilla CAL-55}\newline {\color{media}\sffamily\bfseries\scriptsize MEDIA} & f. 79\newline{\scriptsize Calibración} & «Márgenes estrechos» aplana una diferencia de casi veinte veces en la componente cortante (0,16 \% de SAP2000 frente a 2,9 \% de EduFEM) que la oración inmediatamente anterior acaba de declarar. & «Ambos modelos reproducen la solución analítica con errores inferiores al 3 \% en todas las componentes; en la cortante la diferencia entre ellos es apreciable (0,16 \% frente a 2,9 \%) y es atribuible a la distinta formulación —continuo de tensión plana frente a cáscara—, por lo que no permite afirmar la superioridad de una… & 15\\\addlinespace[3pt]
\mbox{\casilla CAL-56}\newline {\color{media}\sffamily\bfseries\scriptsize MEDIA} & f. 80\newline{\scriptsize Calibración} & Que dieciocho de los 49 ítems del consenso «definen el criterio» es una selección y una definición operativa del autor, no algo que la fuente afirme; en 2.1.4 y en 3.7.2 se enuncia como hecho. En 1.2 sí se declara como definición propia, de modo que basta trasladar esa fórmula a las otras dos. & 3.7.2: «...cubre las dieciocho destrezas y conceptos que este trabajo adopta como definición operativa del criterio sobre la respuesta». 2.1.4: «De sus 49 ítems este trabajo toma, como definición operativa del criterio sobre la respuesta, los que caen dentro del alcance de la elasticidad plana: ...». & 10\\\addlinespace[3pt]
\mbox{\casilla CAL-57}\newline {\color{media}\sffamily\bfseries\scriptsize MEDIA} & f. 80\newline{\scriptsize Calibración} & El valor pedagógico de decisiones de diseño y de resultados numéricos se afirma como hecho, y el software se vuelve sujeto de verbos de enseñanza («ejercita», «explican», «profundiza en la comprensión»), cuando quien ejercita o comprende sería el estudiante y eso no se observó. Ocho apariciones ubicadas, además de las que se reportan por separado por su ubicación (CALb-14, CALb-18, CALb-21, CALb-29, CALb-31). Dos de ellas están fuera de esta zona, en la Introducción (folios 4 y 9), y corresponden al grupo de esa unidad. & Sustituir el valor afirmado por la propiedad o la intención, y poner al estudiante como sujeto posible y al software como medio: «es pedagógicamente significativo» -\textgreater{} «tiene una lectura de modelado: muestra que el modelo de elasticidad bidimensional captura la contribución del cortante a la flecha»; «las que EduFEM ejercita» -\textgreater{}… & 30\\\addlinespace[3pt]
\mbox{\casilla CAL-58}\newline {\color{media}\sffamily\bfseries\scriptsize MEDIA} & f. 81\newline{\scriptsize Calibración} & El pasaje está bien construido en lo esencial —declara la analogía, remite a la prueba de campo y lleva localizadores de página—, pero «sí lo midieron» atribuye a esas fuentes una medición del juicio sobre la respuesta que, según los lectores de 1.2, corresponde a una evaluación cualitativa sin control y a una mejora entre cohortes atribuida sobre todo a la autoría de los alumnos; y «se fundamenta» presenta la expectativa como ya fundada sobre EduFEM. Tres apariciones con la misma redacción. & Fórmula única para las tres apariciones: «Que esa exposición mejore el criterio del estudiante para interpretar la respuesta estructural es la parte de la hipótesis que este trabajo no contrasta. La literatura reporta ese efecto en herramientas de la misma familia [11, pp. 4-5, 13; 9, p. 168]\allowbreak{}, lo que sustenta la expectativa… & 20\\\addlinespace[3pt]
\mbox{\casilla CAL-59}\newline {\color{media}\sffamily\bfseries\scriptsize MEDIA} & f. 81\newline{\scriptsize Calibración} & El reclamo de primacía está acotado a «los antecedentes revisados», pero el texto no dice cuántos son ni con qué criterio se reunieron, de modo que en la lectura queda con más fuerza de la que el alcance declarado permite. & «...aportan a esa tabla su última fila: una verificación y validación publicada junto con la herramienta, que ninguno de los antecedentes reunidos en la Sección 1.x —n herramientas educativas de MEF localizadas con los criterios allí declarados— documenta.» Completar sección, número y criterio de búsqueda. & 15\\\addlinespace[3pt]
\mbox{\casilla CAL-71}\newline {\color{baja}\sffamily\bfseries\scriptsize BAJA} & f. 86\newline{\scriptsize Calibración} & El adjetivo «genuina» refuerza una premisa que el trabajo toma de la literatura y no verificó en su propia población, sin marcar esa procedencia. & «...una dificultad de aprendizaje documentada en la literatura sobre la enseñanza del método y ligada a su carácter abstracto». & 5\\\addlinespace[3pt]
\mbox{\casilla CAL-60}\newline {\color{media}\sffamily\bfseries\scriptsize MEDIA} & f. 88\newline{\scriptsize Calibración} & Califica una decisión de diseño como «aporte pedagógico central» sin que ningún resultado del trabajo mida efecto pedagógico alguno; la fórmula se repite, con otras palabras, en 3.7.1 y en Aportes. & «...conectando la abstracción matemática con la geometría concreta del problema: es la decisión de diseño sobre la que se apoya el aporte pedagógico que este trabajo propone, y cuyo efecto sobre el criterio del estudiante queda pendiente de la prueba de campo especificada en las recomendaciones.» & 15\\\addlinespace[3pt]
\mbox{\casilla CAL-61}\newline {\color{media}\sffamily\bfseries\scriptsize MEDIA} & f. 89\newline{\scriptsize Calibración} & «Confirman ... en los términos en que fueron enunciados» incluye literalmente las cláusulas finalistas del quinto objetivo específico («para establecer su potencial para formar criterio») y del objetivo general («para mejorar el criterio»), y las da por satisfechas; el «y, con ello, el objetivo general» extiende a esa finalidad una comprobación que solo alcanzó al artefacto y a su contenido. La salvedad que sigue cubre la hipótesis, no los objetivos. & Retocar solo la oración existente, sin tocar la siguiente: «En conjunto, los resultados muestran que se alcanzaron los cinco objetivos específicos en lo que cada uno tiene de comprobable en esta etapa —los dos primeros como diagnóstico y fundamento, no como resultados medidos; el tercero con módulos hasta el ensamblaje y las… & 25\\\addlinespace[3pt]
\mbox{\casilla CAL-62}\newline {\color{media}\sffamily\bfseries\scriptsize MEDIA} & f. 89\newline{\scriptsize Calibración} & El impersonal y la agencia atribuida a los propios instrumentos («la tabla estableció», «la matriz estableció», «los instrumentos señalan») ocultan que fue el autor quien construyó las tablas y emitió el juicio, justo el límite que Limitaciones reconoce cuatro páginas después. Cuatro apariciones en el mismo capítulo. & Explicitar el sujeto y usar el verbo fijado para la evidencia cotejada por el autor: «el autor cotejó, ítem por ítem y con el criterio documentado en la Subsección 2.1.5, que las quince destrezas ... tienen instrumento identificable en EduFEM»; «al confrontar el diseño con los cuatro principios de aprendizaje, cada uno resultó… & 30\\\addlinespace[3pt]
\mbox{\casilla CAL-63}\newline {\color{media}\sffamily\bfseries\scriptsize MEDIA} & f. 90\newline{\scriptsize Calibración} & «Habitualmente disociadas» generaliza a todos los recursos educativos sobre el MEF una observación que procede del conjunto acotado revisado en el Capítulo 1, y «riguroso» adjudica al motor una cualidad que la batería de V\&V no mide como tal. Es una afirmación de novedad sin acotar, justo donde el tribunal buscará el aporte. & «...una herramienta de software libre, funcional y verificada, que articula tres capas que los recursos educativos sobre el MEF revisados como antecedentes presentan por separado: un motor de cálculo verificado por soluciones manufacturadas y contrastado con dos referencias externas, una interfaz de modelado interactiva y un… & 10\\\addlinespace[3pt]
\mbox{\casilla CAL-64}\newline {\color{media}\sffamily\bfseries\scriptsize MEDIA} & f. 90\newline{\scriptsize Calibración} & «Completa y verificada ... para los dos estados planos» no distingue lo que la batería cubre de lo que no: la deformación plana solo tiene verificación por MMS y ninguna referencia externa, la viga de Timoshenko se resolvió únicamente con Q9 y de la membrana de Cook solo se contrasta el desplazamiento, de modo que las tensiones del Q4 no se compararon con ninguna referencia externa. La propia recomendación del folio 94 reconoce el hueco de la deformación plana, que Limitaciones no declara. & Cambio mínimo en Aportes: «el motor implementa la formulación isoparamétrica Q4 y Q9 para los dos estados planos de la elasticidad, verificada por soluciones manufacturadas en ambos estados y contrastada con referencias externas en tensión plana». Cambio principal, en Limitaciones: añadir que el contraste con referencias… & 30\\\addlinespace[3pt]
\mbox{\casilla CAL-65}\newline {\color{media}\sffamily\bfseries\scriptsize MEDIA} & f. 91\newline{\scriptsize Calibración} & «Atiende de forma directa al criterio» atribuye a la herramienta un efecto sobre la variable dependiente, que no se midió; el enunciado aparece en Aportes, donde el tribunal lee lo que el trabajo reclama haber logrado. & «...EduFEM la expone con sus modos crudo y suavizado, sus reacciones contrastadas con la carga y su círculo de Mohr: es la parte del diseño orientada de forma directa a las destrezas de interpretación y verificación de la respuesta que recoge la tabla de especificaciones.» & 10\\\addlinespace[3pt]
\mbox{\casilla CAL-66}\newline {\color{media}\sffamily\bfseries\scriptsize MEDIA} & f. 91\newline{\scriptsize Calibración} & Hasta «verificar la respuesta» son propiedades del artefacto; «cerrando el lazo entre la teoría, la exploración y el cálculo manual» afirma una integración cognitiva que nadie observó. & Suprimir la cláusula final y cerrar en: «...y usar para verificar la respuesta por sus propios medios.» Si se quiere conservar la idea, enunciarla como propiedad del documento: «...un registro que enlaza la teoría, la exploración y el cálculo manual sobre un mismo modelo». & 10\\\addlinespace[3pt]
\mbox{\casilla CAL-67}\newline {\color{media}\sffamily\bfseries\scriptsize MEDIA} & f. 91\newline{\scriptsize Calibración} & Se declaran reutilizables dos instrumentos que se aplicaron a un único recurso y por un único evaluador, sin estudio de acuerdo entre jueces ni ensayo sobre un segundo recurso que respalde esa propiedad. & «Desde el punto de vista metodológico, el trabajo propone dos instrumentos para evaluar recursos de enseñanza del MEF sin panel propio de expertos: la tabla de especificaciones construida sobre los ítems del consenso publicado de Pérez-Santiago y Campos, y la matriz de principios. Ambos se aplicaron aquí a un único recurso y… & 20\\\addlinespace[3pt]
\mbox{\casilla CAL-68}\newline {\color{media}\sffamily\bfseries\scriptsize MEDIA} & f. 92\newline{\scriptsize Calibración} & En una recomendación el efecto solo cabe como hipótesis a contrastar; «profundizando en la comprensión» lo enuncia como consecuencia segura de un módulo que aún no existe. & «...permitirían un módulo educativo dedicado que contraste, sobre un mismo modelo, el comportamiento bloqueado y el corregido: el fenómeno que hoy solo se ilustra con el caso de referencia de Cook pasaría a poder manipularse, y su efecto sobre la comprensión del estudiante podría contrastarse en la prueba de campo.» & 15\\\addlinespace[3pt]
\mbox{\casilla CAL-69}\newline {\color{media}\sffamily\bfseries\scriptsize MEDIA} & f. 94\newline{\scriptsize Calibración} & La redacción presupone que la utilidad didáctica ya está fundamentada y que a la prueba solo le queda documentarla, cuando lo que la prueba debe decidir es si existe; la última frase de la tesis repite el supuesto al hablar de «ampliar» una fundamentación empírica que hoy no hay. & «Los resultados de esa prueba, favorables o no, permitirían afinar el diseño de los módulos educativos sobre evidencia y determinar si la utilidad didáctica que este trabajo supone por diseño y apoya en la literatura se verifica en el aula.» Y en el cierre: «...al tiempo que amplía su alcance numérico y su biblioteca de… & 15\\\addlinespace[3pt]
\end{longtable}}
\clearpage
\section{Correcciones por capítulo}
Las correcciones que no son de calibración, en el orden en que conviene hacerlas: primero las secciones donde vive el argumento (Introducción, 2.1, 3.6–3.8, Conclusiones, Resumen), luego el marco teórico, el software, los resultados numéricos y, al final, preliminares, referencias y anexos. Dentro de cada sección, por orden de página. Al comienzo de cada una se recuerdan los ID de calibración de la sección 3 que la tocan, para hacerlos en la misma sentada.
\subsection{Introducción}
{\small\color{gris}\sffamily 2 correcciones críticas o altas y 12 medias o bajas en esta sección · 6,7 h · más 4 frases de calibración (sección 3)}\par
{\small\rotulo{Primero, la calibración de la sección 3}CAL-04 (f. 2), CAL-05 (f. 4), CAL-01 (f. 5), CAL-06 (f. 6).\par}
\begin{ficha}{alta}{\casilla\textbf{TRZ-02}\quad\chip{alta}{alta}\quad Introducción — Metodología de la investigación · folio 7 · PDF 22\quad{\color{gris}\footnotesize también RED-01}\hfill 5 min}
\rotulo{Dice}\cita{«…confronta los instrumentos del software con las destrezas que un consenso publicado de expertos exige y con los principios de aprendizaje de la literatura», oración que clasifica el enfoque como documental y cierra con [9, p. 168]\allowbreak{}.}\par
\rotulo{Problema}La página 168 de Lee respalda en el resto de la tesis la evaluación cualitativa que el autor aplicó en sus cursos, no una clasificación metodológica; usada aquí queda fuera de lugar, y la oración equivalente de 2.1.1 no la lleva. Defecto menor de pertinencia, no de exactitud.\par
\rotulo{Corrección · verificar}Retirar la cita de esa oración —la clasificación del enfoque ya se sostiene en 2.1.1 con [7]\allowbreak{}, [8]\allowbreak{} y [25]\allowbreak{}— o sustituirla por la fuente metodológica que corresponda, de modo que las dos oraciones gemelas citen lo mismo.\par
\rotulo{También en}\begin{itemize}[leftmargin=1.2em,itemsep=1pt,topsep=1pt]
\item {\small §2.1.1 — Tipo y enfoque de investigación, f. 40: \cita{la oración gemela que clasifica el enfoque no lleva esa cita}}
\end{itemize}
{\small\color{gris}\rotulo{RED-01}La clasificación metodológica completa se desarrolla dos veces con redacción casi idéntica, a treinta páginas de distancia: unas veinte líneas duplicadas. §2.1.1 es la única que aporta las fuentes metodológicas ([7, pp. 90-92]\allowbreak{}, [25, p. 6]\allowbreak{}) y el modelo de simulación numérica; la versión de la Introducción es una copia sin respaldo que además adelanta lo que el capítulo 2 debe establecer.\par}
\end{ficha}
\begin{ficha}{alta}{\casilla\textbf{RED-02}\quad\chip{alta}{alta}\quad Introducción /\allowbreak{} Estructura del documento · folio 10 · PDF 25\hfill 45 min}
\rotulo{Dice}\cita{Verificado con ubicar.py: «los criterios de aceptación con que se contrasta la hipótesis» aparece tres veces (PDF 23, 25 y 52); «el modelo de simulación numérica que sostiene la verificación», dos (PDF 23 y 52); «los antecedentes y el estado del arte del software», dos (PDF 25 y 27). En PDF 25 y 52 la enumeración de los seis contenidos de §2.1 se recita completa y en el mismo orden.}\par
\rotulo{Problema}El índice del documento se recita en prosa cinco veces: «Estructura del documento», el cierre de «Metodología de la investigación», la apertura de cada capítulo y una segunda vez al abrir §2.1. Ninguna repetición añade contenido; el lector recorre el mismo sumario antes de llegar a él. Es el mayor bloque de texto parásito del documento ({\mates ≈} una página).\par
\rotulo{Corrección · reformular}Aparición canónica: «Estructura del documento» (folio 10, PDF 25), que conserva el recorrido completo. Se elimina: la enumeración de los seis apartados de §2.1 en el cierre de «Metodología de la investigación» (folio 8), sustituida por «El diseño metodológico se desarrolla en la sección 2.1»; la segunda lista de la apertura de §2.1 (folio 37), redundante con la apertura del capítulo 2 en la misma página. Se reduce a dos líneas cada apertura de capítulo, enunciando su tesis y no su índice (capítulo 2: «Este capítulo articula el modelo de investigación y su materialización en software: primero el diseño metodológico, después el software que lo concreta»), y lo mismo en la apertura de §2.2. Ahorro: una página.\par
\rotulo{También en}\begin{itemize}[leftmargin=1.2em,itemsep=1pt,topsep=1pt]
\item {\small Introducción /\allowbreak{} Metodología de la investigación (cierre), f. 8: \cita{El diseño metodológico —el modelo de simulación numérica que sostiene la verificación, las variables y su operacionalización, la matriz de consistencia, la población, los casos de estudio y su criterio de selección, el procedimiento con sus instrumentos numéri}}
\item {\small Capítulo 1, párrafo de apertura, f. 12: \cita{Este capítulo abre situando el trabajo en su contexto —los antecedentes y el estado del arte del software para la enseñanza del MEF, que diagnostican la brecha y justifican la oportunidad de EduFEM—}}
\item {\small Capítulo 2, párrafo de apertura, f. 37: \cita{Se presenta primero el diseño metodológico —el tipo de investigación y el modelo de simulación numérica que sostiene la verificación, las variables y su matriz de consistencia, la población, los casos de estudio y su criterio de selección, el procedimiento con}}
\item {\small §2.1, apertura, f. 37: \cita{esta sección lo desarrolla siguiendo la estructura de un modelo de investigación: la conceptualización y el alcance del modelo, la definición de las variables, el procedimiento y la toma de observaciones, la validación de esas observaciones y el contraste de l}}
\item {\small §2.2, apertura, f. 48: \cita{el pre-proceso interactivo, los módulos educativos, el post-proceso}}
\end{itemize}
\end{ficha}
{\small\rotulo{Párrafos que pide la defensa}\par}
\begin{ficha}{media}{\casilla\textbf{DEF-08}\quad\chip{media}{media}\quad Introducción — Alcance y limitaciones · folio 8 · PDF 23\hfill 25 min}
\rotulo{Dice}\cita{«no comprende el análisis de sistemas estructurales discretos —pórticos o reticulados—, ni la formulación de placas, cáscaras o problemas tridimensionales» (verificada literal en PDF 23 /\allowbreak{} folio 8).}\par
\rotulo{Problema}Pregunta más probable del miembro estructuralista: «yo calculo pórticos y losas; si su programa no hace nada de eso, ¿para qué me sirve en ingeniería civil?». La delimitación enumera con precisión lo que queda fuera —justo lo que un ingeniero civil calcula a diario— y en ninguna parte enumera los problemas propios de la profesión que sí caen dentro del alcance, de modo que el tribunal se lleva la impresión de que la herramienta no sirve para su disciplina.\par
\rotulo{Corrección · anadir}Añadir una sola frase en positivo tras la exclusión (folio 8, PDF 23), sin atribuir validación a lo que no se validó: «Dentro de ese dominio quedan los problemas planos habituales en ingeniería civil: en tensión plana, vigas de gran canto y muros o machones cargados en su plano; en deformación plana, secciones de presas, muros de contención y túneles, donde una rebanada representa todo el sólido. La verificación y la validación de este trabajo se realizan sobre los casos del capítulo 3; la validación con un caso civil de deformación plana con solución de referencia publicada se plantea en las recomendaciones.» Mantener los ejemplos ya presentes en el capítulo 1 («una presa o un tubo largo») y no introducir familias que el documento no menciona en ninguna parte, como las regiones D o las ménsulas cortas, salvo que se añadan también al cuerpo.\par
\rotulo{También en}\begin{itemize}[leftmargin=1.2em,itemsep=1pt,topsep=1pt]
\item {\small §3.6 Validación con la viga de Timoshenko, f. 68: \cita{para que las magnitudes del ejemplo resulten familiares al lector de ingeniería civil}}
\item {\small Limitaciones observadas, f. 85: \cita{limitan la aplicabilidad a casos de ingeniería más allá de los didácticos}}
\end{itemize}
{\footnotesize\color{gris}\rotulo{Verificación}ajustado: La cita es literal y está en PDF 23 /\allowbreak{} folio 8, y es cierto que la delimitación enumera las exclusiones con precisión. Pero se REFUTA el «en ninguna parte»: el documento sí nombra problemas civiles que caen dentro del alcance —al presentar la deformación plana dice que «corresponde, por ejemplo, a una presa o un tubo largo donde una…\par}
\end{ficha}
\begin{ficha}{media}{\casilla\textbf{DEF-09}\quad\chip{media}{media}\quad Introducción — Alcance y limitaciones · folio 8 · PDF 23\hfill 25 min}
\rotulo{Dice}\cita{«elementos cuadriláteros isoparamétricos de cuatro nodos (Q4) y de nueve nodos (Q9), que cubren respectivamente la interpolación bilineal y la cuadrática completa» (verificada literal en PDF 23 /\allowbreak{} folio 8).}\par
\rotulo{Problema}Pregunta previsible: «¿por qué solo Q4 y Q9, si el programa de 1998 que usted cita como antecedente ya tenía triángulos, Q4, Q8 y Q9?». La comparación la pone a la vista la propia tesis en 1.1 y la deja sin responder, de modo que el tribunal se lleva la idea de un retroceso de veinticinco años; el texto justifica el orden de interpolación pero nunca la ausencia de las otras familias.\par
\rotulo{Corrección · anadir}Añadir una o dos oraciones tras la justificación del orden de interpolación (folio 8, PDF 23), formuladas como decisión de diseño y no como juicio sobre otras familias: «La pareja Q4-Q9 es el par mínimo que contrasta una base polinómica incompleta con una completa dentro de una sola familia de funciones de forma y una sola regla de cuadratura por elemento, lo que mantiene las matrices del seguimiento paso a paso en un tamaño legible; los elementos triangulares y los cuadriláteros serendípitos quedan fuera de esta versión, y su incorporación se plantea en las recomendaciones (folio 92).» En el pasaje de antecedentes (folio 14, PDF 29), añadir la remisión en una línea para que la comparación con ED-Elas2D no quede abierta: que la diferencia de biblioteca es deliberada y se explica en la delimitación del alcance.\par
\rotulo{También en}\begin{itemize}[leftmargin=1.2em,itemsep=1pt,topsep=1pt]
\item {\small §1.1 Antecedentes (ED-Elas2D), f. 14: \cita{una biblioteca de elementos triangulares y cuadriláteros que incluye}}
\item {\small Recomendaciones, Más tipos de elemento, f. 93: \cita{La biblioteca de elementos puede ampliarse con elementos triangulares}}
\end{itemize}
{\footnotesize\color{gris}\rotulo{Verificación}ajustado: Las citas son literales: la delimitación en PDF 23 /\allowbreak{} folio 8 y el pasaje sobre ED-Elas2D en PDF 29 /\allowbreak{} folio 14 («una biblioteca de elementos triangulares y cuadriláteros que incluye los elementos Q4, Q8 y Q9»). Se sostiene el núcleo: el documento no justifica en ninguna parte la ausencia de triángulos y del Q8 —la búsqueda de «Q8»,…\par}
\end{ficha}
\subsubsection*{Medias y bajas}
{\small\begin{longtable}{@{}L{1.6cm}L{1.6cm}L{5.25cm}L{5.65cm}r@{}}
\toprule
\sffamily\bfseries ID & \sffamily\bfseries Dónde & \sffamily\bfseries Qué pasa & \sffamily\bfseries Qué hacer & \sffamily\bfseries min\\
\midrule\endhead
\bottomrule\endfoot
\mbox{\casilla RED-20}\newline {\color{baja}\sffamily\bfseries\scriptsize BAJA} & f. 1\newline{\scriptsize Redundancia} & El verbo «encarna» y los conectores «de forma» /\allowbreak{} «de manera» se repiten hasta convertirse en tic de redacción. No hay error de lengua, pero la monotonía léxica se nota en un texto de noventa y cinco folios y delata pasajes escritos en serie. & Pasada de revisión léxica al final del proceso, cuando el texto ya no vaya a cambiar: sustituir parte de las apariciones de «encarna»/\allowbreak{}«encarnación» por «concreta», «materializa» o «se expresa en», y variar «de forma X» /\allowbreak{} «de manera X» por el adverbio en -mente o por la reformulación directa, siempre que no se pierda precisión… & 60\\\addlinespace[3pt]
\mbox{\casilla APF-05}\newline {\color{media}\sffamily\bfseries\scriptsize MEDIA} & f. 2\newline{\scriptsize Aparato formal} & Vancouver no repite el número de referencia dentro del mismo corchete: dos localizadores de una misma fuente van agrupados. Tal como está impreso, el lector entiende que se citan dos trabajos distintos, y precisamente en la frase que apoya el planteamiento del problema sobre el consenso de expertos. & Imprimir «[3, pp. 1159, 1171]\allowbreak{}»: sustituir \textbackslash{}autocites[p.\textasciitilde{}1159]\allowbreak{}\{perezsantiago2023fem\}\allowbreak{}[p.\textasciitilde{}1171]\allowbreak{}\{perezsantiago2023fem\}\allowbreak{} por \textbackslash{}autocite[pp.\textasciitilde{}1159, 1171]\allowbreak{}\{perezsantiago2023fem\}\allowbreak{}. Revisar de paso si el patrón se repite en otros capítulos con \textbackslash{}autocites de clave única. & 15\\\addlinespace[3pt]
\mbox{\casilla RED-06}\newline {\color{media}\sffamily\bfseries\scriptsize MEDIA} & f. 3\newline{\scriptsize Redundancia} & Tres exposiciones del mismo respaldo bibliográfico. La diferencia no es solo de economía: únicamente §1.2 acota la evidencia (cualitativa, dependiente de la visualización y del apoyo instruccional, atribuida en parte a la autoría del alumno), de modo que en la Justificación y en §2.2.1 el lector recibe la afirmación de mejora sin las salvedades que la hacen defendible. & Aparición canónica: §1.2 (folio 18, PDF 33), que conserva la evidencia con sus salvedades; esas salvedades no deben perderse. La Justificación se reduce a una frase con remisión: «La literatura sobre enseñanza del método reporta que la simulación interactiva y la construcción paso a paso de las herramientas de cálculo mejoran… & 30\\\addlinespace[3pt]
\mbox{\casilla TRZ-10}\newline {\color{media}\sffamily\bfseries\scriptsize MEDIA} & f. 6\newline{\scriptsize Trazabilidad} & La cláusula (b) de la hipótesis extiende el umbral del 3 \% en la flecha también frente a SAP2000, pero 2.1.6 solo lo fija frente a la solución analítica; 3.8 aplica después ese 3 \% a los desplazamientos frente a SAP2000, con lo que el criterio que se contrasta no es el que se declaró. El enunciado cambia entre sus tres apariciones. & Añadir el criterio que falta en 2.1.6 —«y diferencia del desplazamiento vertical inferior al 3 \% frente a SAP2000 en los tres puntos de control»— o, alternativamente, acotar la cláusula en la Introducción: «[...]\allowbreak{} del 3 \% en la flecha frente a la solución analítica y del 1 \% en la tensión normal frente a la solución analítica y… & 10\\\addlinespace[3pt]
\mbox{\casilla RED-07}\newline {\color{media}\sffamily\bfseries\scriptsize MEDIA} & f. 6\newline{\scriptsize Redundancia} & La Introducción remite a §2.1.6 para los criterios y, pese a la remisión, los enuncia igualmente con sus valores numéricos, que reaparecen en forma completa cuarenta páginas después. La duplicación de cifras obliga además a mantener sincronizados dos lugares si un umbral cambia. & Aparición canónica de los valores: §2.1.6 (folio 46, PDF 61). La Introducción enuncia las tres cláusulas por su nombre y su contenido cualitativo y remite para las cifras: «…mediante tres cláusulas —(a) exposición del canal y de la respuesta, (b) corrección del motor y (c) validez de contenido y coherencia— cuyos umbrales… & 20\\\addlinespace[3pt]
\mbox{\casilla CLA-05}\newline {\color{media}\sffamily\bfseries\scriptsize MEDIA} & f. 6\newline{\scriptsize Claridad} & La oración que define el criterio de comprobación de la hipótesis —el núcleo argumental de la tesis, lo que el tribunal citará para juzgar si se probó lo que se dijo— tiene cerca de 160 palabras, tres cláusulas rotuladas, dos incisos con raya y cinco puntos y coma, de modo que ningún criterio queda visualmente aislado ni es verificable de un vistazo. & Encuadre breve más lista de tres ítems, cada cláusula cerrada con punto: «En esta etapa la hipótesis se comprueba en su parte de diseño y contenido. Los criterios se fijaron a priori en la Subsección 2.1.6 y son tres. (a) Exposición: el software expone cada etapa del canal de cálculo —en un módulo educativo o en la memoria de… & 25\\\addlinespace[3pt]
\mbox{\casilla CLA-06}\newline {\color{media}\sffamily\bfseries\scriptsize MEDIA} & f. 6\newline{\scriptsize Claridad} & «Bajo el desarrollo de EduFEM» no es una construcción del español corriente y oscurece el papel de cada variable justo en el enunciado que el tribunal leerá con más atención; que EduFEM sea la variable interviniente solo se aclara en la oración siguiente. & «Mediante el desarrollo de EduFEM —variable interviniente—, la exposición transparente e interactiva del canal de cálculo y de la respuesta estructural —variable independiente— podría mejorar el criterio para interpretar la respuesta obtenida por el MEF —variable dependiente— en los estudiantes de la Carrera de Ingeniería Civil… & 15\\\addlinespace[3pt]
\mbox{\casilla REF-02}\newline {\color{media}\sffamily\bfseries\scriptsize MEDIA} & f. 9\newline{\scriptsize Referencias} & De las 174 llamadas de cita del cuerpo, 37 van sin localizador de página; 22 de ellas, repartidas en 12 puntos del texto, sostienen una definición o una atribución concreta y por tanto lo exigen según la regla que la propia tesis declara en la Introducción (folio 11). Las otras 15 son de encuadre y están bien sin él. Las omisiones no son periféricas: incluyen la definición del bloqueo por cortante y de los modos espurios —el fenómeno del que depende todo el resultado del Capítulo 3—, el criterio con que se eligen los casos de prueba —la columna vertebral metodológica de la validación—, la afirmación de que la formulación se generaliza a tres dimensiones y las dos recomendaciones técnicas del cierre. Un tribunal que coteje el apartado de la Introducción con el aparato encuentra la norma incumplida en 12 puntos, y en ninguno puede ir a la página donde la fuente dice lo que se le atribuye. & Antes de aplicar, recontar las llamadas sin localizador con un criterio explícito (llamada = corchete, o fuente dentro del corchete) y descartar los falsos positivos del barrido —por ejemplo, el intervalo «[0, 1]\allowbreak{}»—; las cifras 174/\allowbreak{}37/\allowbreak{}22 no deben figurar en el informe sin ese recuento. Añadir después el localizador, con el… & 90\\\addlinespace[3pt]
\mbox{\casilla CLA-07}\newline {\color{media}\sffamily\bfseries\scriptsize MEDIA} & f. 10\newline{\scriptsize Claridad} & El título del Capítulo 2 y su glosa en la Introducción usan «modelo» en la misma frase para el diseño metodológico de la investigación (2.1: variables, población, instrumentos) y para el modelo numérico Q4/\allowbreak{}Q9 (2.2), lo que puede llevar al tribunal no especialista a esperar contenido de elementos finitos ya en 2.1. & Si el título del capítulo está fijado por el reglamento de la Facultad y no puede reformularse, añadir una frase de apertura en 2.1 que distinga las dos acepciones: «Este capítulo llama modelo tanto al diseño de la investigación (Sección 2.1: variables, población e instrumentos) como al modelo numérico de elementos finitos que… & 15\\\addlinespace[3pt]
\mbox{\casilla CLA-08}\newline {\color{media}\sffamily\bfseries\scriptsize MEDIA} & f. 10-11\newline{\scriptsize Claridad} & El apartado es un único párrafo de unas setecientas palabras que recorre tres capítulos, conclusiones, referencias y siete anexos, y que además aloja la declaración de las normas de citación y presentación —Vancouver con localizadores y APA 7—, donde nadie la buscará. & Dividir en un párrafo por capítulo (marco teórico, diseño e implementación, resultados, conclusiones y anexos), dejar la frase de trazabilidad de objetivos como párrafo propio y trasladar la declaración de normas a un apartado breve titulado «Normas de citación y presentación», al final de la Introducción o en el material… & 30\\\addlinespace[3pt]
\mbox{\casilla TRZ-11}\newline {\color{media}\sffamily\bfseries\scriptsize MEDIA} & f. 11\newline{\scriptsize Trazabilidad} & El primer objetivo específico se cierra de forma explícita al final de 1.1, pero el segundo no recibe oración de cierre en el Capítulo 1: su cumplimiento solo se declara después, en 3.8 y en las Conclusiones, de modo que el lector que busque el resultado del segundo objetivo en su lugar natural no lo encuentra. No se releyeron las secciones 1.3 a 1.13, conviene comprobar que ese cierre no esté redactado con otras palabras. & Cerrar el Capítulo 1 con un párrafo breve, o con una tabla expresión /\allowbreak{} módulo /\allowbreak{} capítulo de la memoria: «Con ello queda establecida la base común del motor, de los módulos educativos y de la memoria de cálculo, que es el segundo objetivo específico: cada expresión de este capítulo se retoma en el motor (sección 2.2), en el… & 20\\\addlinespace[3pt]
\mbox{\casilla TRZ-12}\newline {\color{media}\sffamily\bfseries\scriptsize MEDIA} & f. 11\newline{\scriptsize Trazabilidad} & Los Anexos C (listados de código) y F (formatos de interoperabilidad) se mencionan una sola vez en todo el cuerpo, en el párrafo que enumera los anexos al final de la Introducción, y ninguno de los dos remite a su vez a un apartado concreto del cuerpo; a diferencia de B, D, E y G, que sí se retoman en capítulos posteriores, quedan sin razón de consulta durante la lectura. & Añadir una remisión de propósito por anexo: desde la subsección 2.2.5, al presentar las funciones de forma y el ensamblaje disperso, «véase el listado completo en el Anexo C»; y desde la subsección 2.2.9, al describir los formatos .edufem, CSV y DXF, «la especificación de cada formato está en el Anexo F». Así el lector sabe… & 20\\\addlinespace[3pt]
\end{longtable}}
\subsection{Capítulo 2 · 2.1 Diseño metodológico}
{\small\color{gris}\sffamily 2 correcciones críticas o altas y 11 medias o bajas en esta sección · 5,0 h · más 1 frases de calibración (sección 3)}\par
{\small\rotulo{Primero, la calibración de la sección 3}CAL-08 (f. 45).\par}
\begin{ficha}{alta}{\casilla\textbf{TRZ-05}\quad\chip{alta}{alta}\quad §2.1.4 Población, universo de contenido y casos de estudio · folio 41 · PDF 56\quad{\color{gris}\footnotesize también DEF-05}\hfill 20 min}
\rotulo{Dice}\cita{La población del problema científico son los estudiantes de la Carrera de Ingeniería Civil de la Universidad Autónoma Tomás Frías que cursan las asignaturas en que se introduce el análisis por el MEF}\par
\rotulo{Problema}La población del estudio —y la de la prueba de campo que se diseña en las recomendaciones— se define por una perífrasis, sin nombre, sigla ni semestre de la asignatura, y oscila entre el plural («las asignaturas», 2.1.4) y el singular («la asignatura», Recomendaciones); el tribunal local conoce el plan de estudios y detectará el hueco de inmediato. La fuente LaTeX conserva en ese punto un comentario «DATO PENDIENTE» que no llega a imprimirse. Corrección de ubicación respecto del hallazgo bruto: la oración está en el folio 41 (PDF 56), no en el 42-43.\par
\rotulo{Corrección · anadir}Nombrar la asignatura con su sigla y su semestre según el plan de estudios vigente: «[...]\allowbreak{} que cursan \textless{}Asignatura, sigla, N.º semestre\textgreater{}, asignatura en que se introduce el análisis por el MEF», y usar ese mismo nombre, en singular, en las Recomendaciones (PDF 109). Si en el plan vigente el MEF se introduce en más de una asignatura, conservar el plural en los dos lugares y nombrarlas todas.\par
\rotulo{También en}\begin{itemize}[leftmargin=1.2em,itemsep=1pt,topsep=1pt]
\item {\small Recomendaciones (prueba de campo), f. 94: \cita{los estudiantes de la Carrera de Ingeniería Civil de la Universidad Autónoma Tomás Frías que cursan la asignatura en que se introduce el MEF}}
\end{itemize}
{\small\color{gris}\rotulo{DEF-05}Pregunta que un tribunal de la propia casa de estudios está en posición de hacer: «¿qué asignatura, de qué semestre y con cuántos estudiantes?». La población nunca se identifica —no se nombra la materia, ni su sigla, ni el semestre, ni la matrícula aproximada— y el documento alterna entre «las asignaturas» (folio 41) y «la asignatura» (folio 94), de modo que la población del propio plan de estudios queda indeterminada justo en el apartado que debe determinarla.\par}
{\footnotesize\color{gris}\rotulo{Verificación}ajustado: Verificadas las dos ubicaciones: la definición de la población está en 2.1.4 (PDF 56, folio 41; la oración se corta ahí por la intercalación de la Tabla 2.2, lo que es un artefacto de maquetado, no del texto) y las Recomendaciones usan el singular «la asignatura en que se introduce el MEF» (PDF 109, folio 94). La inconsistencia de número…\par}
\end{ficha}
\begin{ficha}{alta}{\casilla\textbf{TRZ-06}\quad\chip{alta}{alta}\quad §2.1.6 (criterio de la tensión recuperada) · folio 46 · PDF 61\quad{\color{gris}\footnotesize también CLA-01}\hfill 15 min}
\rotulo{Dice}\cita{una cota empírica mínima declarada como tal —tasa no inferior a 1,4 en Q4 ni a 1,9 en Q9—}\par
\rotulo{Problema}El criterio no dice entre qué mallas se mide la tasa —la Tabla 3.1 da 1,21 para el Q4 entre N = 2 y N = 4, por debajo de 1,4, y solo el pie de la Tabla D.3 aclara que se juzgan las asintóticas—, y una cota «empírica» es por naturaleza posterior a la observación, lo que choca con los «criterios fijados a priori» que el documento invoca siete veces.\par
\rotulo{Corrección · reformular}En 2.1.6: «Las tasas se juzgan en su valor asintótico, medido entre las dos mallas más finas (N = 16 y N = 32)» y «para el campo de tensiones recuperado, cuya tasa teórica no está establecida, se fija un umbral de regresión —1,4 en Q4 y 1,9 en Q9— elegido tras observar el comportamiento del motor: no es un criterio a priori, sino una salvaguarda contra retrocesos».\par
\rotulo{También en}\begin{itemize}[leftmargin=1.2em,itemsep=1pt,topsep=1pt]
\item {\small §3.2, f. 64: \cita{y por encima de la cota empírica mínima de la subsección 2.1.6}}
\item {\small Anexo D, Tabla D.3, f. 124: \cita{Tasas asintóticas de convergencia (entre las mallas}}
\end{itemize}
{\small\color{gris}\rotulo{CLA-01}Los ocho criterios numéricos con que se juzga todo el Capítulo 3 van encadenados en una sola oración de unas 180 palabras con cuatro puntos y coma y tres incisos. En ese formato ningún tribunal puede verificar uno por uno qué se prometió y qué se cumplió, que es exactamente la operación que la defensa exige; además, el único criterio de origen empírico (la cota de la tasa de tensiones) queda enterrado entre los teóricos.\par}
\end{ficha}
{\small\rotulo{Párrafos que pide la defensa}\par}
\begin{ficha}{media}{\casilla\textbf{DEF-12}\quad\chip{media}{media}\quad §2.1.1 Tipo de investigación · folio 38 · PDF 53\hfill 25 min}
\rotulo{Dice}\cita{«por su nivel es descriptiva y comparativa» (verificada literal en PDF 53 /\allowbreak{} folio 38), dentro de la clasificación que también llama «enfoque» a lo documental.}\par
\rotulo{Problema}La clasificación mezcla categorías de taxonomías distintas —«documental» es una vía de recolección y no un enfoque, y «comparativa» no figura en la escala de nivel del autor citado— y el texto no declara qué taxonomía sigue, lo que abre una objeción metodológica fácil ante el miembro que evalúa el marco.\par
\rotulo{Corrección · reformular}Declarar la taxonomía en una frase previa («se adopta la clasificación de [autor]\allowbreak{}, que distingue finalidad, naturaleza, enfoque y nivel»), reubicar «documental» como vía de recolección —«el enfoque es cuantitativo; la recolección combina la medición numérica sobre el motor con el análisis documental de los instrumentos del software»— y, si se conserva «comparativa», justificar en media línea que se usa en el sentido del contraste Q4/\allowbreak{}Q9 y no como nivel de la escala citada.\par
\end{ficha}
\begin{ficha}{media}{\casilla\textbf{DEF-13}\quad\chip{media}{media}\quad §2.1.6 Criterios de aceptación y contraste de la hipótesis · folio 46 · PDF 61\hfill 45 min}
\rotulo{Dice}\cita{«una cota empírica mínima declarada como tal —tasa no inferior a 1,4 en Q4 ni a 1,9 en Q9—» (verificada literal en PDF 61 /\allowbreak{} folio 46); la tasa medida que se contrasta con ella es 1,54 (Tabla 3.1).}\par
\rotulo{Problema}Pregunta de alto rendimiento para quien quiera atacar la V\&V: «la cota de 1,4, ¿la fijó antes o después de medir 1,54?». Los seis umbrales de aceptación de la subsección (±0,5 en las tasas; 1,4 y 1,9; 3 \%; 1 \%; 1,5 \%; 10\textasciicircum{}-8 en el equilibrio) se enuncian sin decir de dónde sale cada cifra ni con qué respaldo, y el más discutible se califica de «empírico» y queda apenas por debajo del valor obtenido, lo que deja abierta la objeción de que se fijó después de ver los resultados pese a que el documento insiste en que sus criterios son a priori.\par
\rotulo{Corrección · anadir}Dar a cada umbral su procedencia en el mismo párrafo y con la extensión mínima: el margen de +-0,5 sobre las tasas teóricas, como tolerancia habitual de la verificación de códigos, con su cita y localizador; el 3 \% de la flecha y el 1,5 \% de Cook, con la referencia de la que se toman o declarados como tolerancia de ingeniería; el residuo de 10\textasciicircum{}-8, remitiendo a la constante de cero numérico del programa que 2.1.5 ya menciona. Para 1,4 en Q4 y 1,9 en Q9, decir de qué cálculo o de qué corrida preliminar salen y remitir al guion del repositorio donde están escritas; si esa procedencia no puede documentarse, degradarlas a observación reportada sin criterio de aceptación —como ya se hace con el desempeño del solucionador— y presentar en 3.2 las tasas medidas (1,54 en Q4 y 2,00 en Q9) sin umbral.\par
\rotulo{También en}\begin{itemize}[leftmargin=1.2em,itemsep=1pt,topsep=1pt]
\item {\small §3.2 Verificación por MMS, f. 63: \cita{por encima de la cota empírica mínima de la Sección 2.1.6}}
\item {\small §2.1.6, f. 46: \cita{error de la flecha central inferior al 3 \% frente a la solución analítica con corrección de cortante}}
\item {\small §2.1.6, f. 46: \cita{desplazamiento del Q9 con N=8 dentro del 1,5 \% del valor de referencia}}
\end{itemize}
{\footnotesize\color{gris}\rotulo{Verificación}ajustado: La cita es literal y está en PDF 61 /\allowbreak{} folio 46, y el contraste de cifras se sostiene: la tasa medida del Q4 es 1,54 frente a la cota 1,4 y la del Q9 es 2,00 frente a 1,9, ambas apenas por encima. Dos correcciones de detalle en la evidencia secundaria: esa lectura está en PDF 79 /\allowbreak{} folio 64 (no PDF 78 /\allowbreak{} folio 63) y el texto dice «la cota…\par}
\end{ficha}
\begin{ficha}{media}{\casilla\textbf{DEF-14}\quad\chip{media}{media}\quad §2.1.6 Criterios de aceptación y contraste de la hipótesis · folio 46 · PDF 61\hfill 30 min}
\rotulo{Dice}\cita{«…inferior al 1 \% frente a la solución analítica y frente a SAP2000 en los tres puntos de control, y equilibrio global» (verificada literal en PDF 61 /\allowbreak{} folio 46).}\par
\rotulo{Problema}La objeción más dañina sobre la validación: «su criterio exige menos del 1 \% de error, pero su propia tabla muestra 4,56 \% en la componente transversal y 2,89 \% en el cortante; ¿no fijó el criterio sobre lo que sabía que iba a pasar?». El umbral se fija solo sobre la componente que pasa con holgura (0,04 \%) y deja fuera del criterio las componentes secundarias; el argumento de escala que lo justifica existe, pero aparece en 3.7.1, después del dato, y no en el lugar donde el criterio se enuncia.\par
\rotulo{Corrección · reformular}Añadir en 2.1.6, dentro de la oración del criterio, la razón del alcance y la remisión: «... error de la tensión normal sigma\_\allowbreak{}x —la magnitud primaria del problema de flexión— inferior al 1 \% frente a la solución analítica y frente a SAP2000 en los tres puntos de control ...», y cerrar el párrafo con «... incluidas las componentes secundarias de tensión, que no forman parte del criterio porque su magnitud en los puntos de control es muy inferior a la de sigma\_\allowbreak{}x y su error relativo está dominado por esa diferencia de escala (Subsección 3.7.1)».\par
\rotulo{También en}\begin{itemize}[leftmargin=1.2em,itemsep=1pt,topsep=1pt]
\item {\small §3.6 Validación con la viga de Timoshenko, f. 71-72: \cita{0,94 \% en sigma\_\allowbreak{}y y 2,89 \% en tau\_\allowbreak{}xy frente a la solución analítica, y 4,56 \% y 2,72 \% frente a SAP2000}}
\item {\small §3.7.1 Alcances, f. 80: \cita{cuya magnitud es entre seis y treinta y cinco veces menor que la de sigma\_\allowbreak{}x}}
\end{itemize}
{\footnotesize\color{gris}\rotulo{Verificación}ajustado: La cita es literal y está en PDF 61 /\allowbreak{} folio 46, y las cifras del contraste también (0,94 \% en sigma\_\allowbreak{}y y 2,89 \% en tau\_\allowbreak{}xy frente a la analítica, 4,56 \% y 2,72 \% frente a SAP2000, PDF 84 /\allowbreak{} folio 69). Pero el diagnóstico atribuye al texto una omisión que no comete. Primero, el criterio nombra su propio alcance donde se enuncia: «error de la…\par}
\end{ficha}
\begin{ficha}{media}{\casilla\textbf{DEF-15}\quad\chip{media}{media}\quad §2.1.5 Procedimiento, instrumentos y reproducibilidad · folio 46 · PDF 61\hfill 25 min}
\rotulo{Dice}\cita{«cada caso se regenera por completo al reejecutar su guion con las dependencias declaradas en el repositorio» (verificada literal en PDF 61 /\allowbreak{} folio 46).}\par
\rotulo{Problema}Comprobación natural de una de las banderas del trabajo: «¿podemos reproducir sus resultados? ¿qué versión exacta del programa produjo estas cifras?». La reproducibilidad se apoya en «la versión del software que acompaña a este documento» y en una dirección de repositorio, sin número de versión, etiqueta, identificador de confirmación ni depósito permanente, de modo que el tribunal no tiene forma de saber qué artefacto exacto produjo las cifras que juzga. La fórmula imprecisa se repite en tres lugares.\par
\rotulo{Corrección · anadir}Fijar la identificación una sola vez, en 2.1.1, con los datos que el proyecto ya tiene —EduFEM \textless{}número de versión del programa\textgreater{}, confirmación \textless{}hash corto\textgreater{} y fecha del repositorio citado en la Introducción, y dependencias declaradas en requirements.txt— y remitir a ella desde 2.1.5 y desde la fila «Versión evaluada» de la Tabla 2.1, indicando además el medio en que el software se entrega al tribunal. No hace falta crear una etiqueta ni un depósito permanente si no existen: bastan versión, confirmación y fecha.\par
\rotulo{También en}\begin{itemize}[leftmargin=1.2em,itemsep=1pt,topsep=1pt]
\item {\small §2.1.1, f. 38: \cita{la versión del software que acompaña a este documento}}
\item {\small Introducción — Justificación (nota al pie), f. 4: \cita{Repositorio del proyecto}}
\end{itemize}
{\footnotesize\color{gris}\rotulo{Verificación}ajustado: La cita es literal y está en PDF 61 /\allowbreak{} folio 46, y el defecto de fondo se sostiene: ni en 2.1.1, ni en 2.1.5, ni en la Tabla 2.1 aparece un número de versión, una etiqueta o una confirmación del repositorio. Dos correcciones. Recuento: la fórmula «la versión del software que acompaña a este documento» aparece dos veces —2.1.1 (PDF 53 /\allowbreak{}…\par}
\end{ficha}
\subsubsection*{Medias y bajas}
{\small\begin{longtable}{@{}L{1.6cm}L{1.6cm}L{5.25cm}L{5.65cm}r@{}}
\toprule
\sffamily\bfseries ID & \sffamily\bfseries Dónde & \sffamily\bfseries Qué pasa & \sffamily\bfseries Qué hacer & \sffamily\bfseries min\\
\midrule\endhead
\bottomrule\endfoot
\mbox{\casilla CLA-13}\newline {\color{media}\sffamily\bfseries\scriptsize MEDIA} & f. 37\newline{\scriptsize Claridad} & La oración que abre el diseño metodológico —la que el tribunal contrastará con el molde exigido— encadena cuatro criterios de clasificación, cuatro incisos entre rayas y tres citas en unas 150 palabras, de modo que ninguna de las cuatro clasificaciones se retiene. La propia Introducción (folio 7) ya presenta el mismo contenido en oraciones separadas, con lo que el capítulo empeora una redacción que el documento ya tenía resuelta. & «Por su finalidad, el trabajo es propositivo: propone un artefacto dirigido a un problema formativo. Por su naturaleza es investigación aplicada de tipo tecnológico, la que produce conocimiento operativo para transformar una realidad concreta [7, pp. 90-92]\allowbreak{}. Por el abordaje del problema combina dos vías: una cuantitativa en la… & 30\\\addlinespace[3pt]
\mbox{\casilla CLA-14}\newline {\color{media}\sffamily\bfseries\scriptsize MEDIA} & f. 37\newline{\scriptsize Claridad} & «Validez de contenido» es el concepto que sostiene la cláusula (c) de la hipótesis y aparece cinco veces en la sección sin definición conceptual ni cita: el texto da su medida operativa —una proporción de destrezas cubiertas— pero nunca dice qué pretende establecer una validez de contenido ni por qué basta un juicio documental para establecerla. Un tribunal no especialista puede leerlo como si el software hubiera sido validado en su efecto formativo, que es precisamente lo que el trabajo no midió. & Añadir una definición conceptual breve allí donde el término se opera por primera vez con peso metodológico, en §2.1.2 (folio 39), justo antes de la medida operativa: «La validez de contenido es el grado en que un instrumento representa el dominio de contenido que dice cubrir; se establece contrastando el instrumento con ese… & 25\\\addlinespace[3pt]
\mbox{\casilla TRZ-15}\newline {\color{media}\sffamily\bfseries\scriptsize MEDIA} & f. 39\newline{\scriptsize Trazabilidad} & Dos de los cinco indicadores de la variable independiente no tienen criterio en 2.1.6 ni cifra en 3.6.1: del conmutador solo se dice que lo incorporan los módulos «que exhiben formulación», sin contarlos, y la coincidencia entre lo expuesto y lo calculado se argumenta por construcción y no se mide, aunque la Tabla 2.1 le asigna como instrumento unas pruebas que solo comprueban que la memoria se genera sin error. & En 3.6.1, dar la cifra del conmutador («k de los siete módulos incorporan el conmutador entre fórmula y valor: M...») y declarar el estatuto del otro indicador: «la coincidencia se sostiene por diseño —los módulos y la memoria llaman a las mismas funciones del motor— y se comprobó en el ejemplo canónico cotejando n valores de… & 30\\\addlinespace[3pt]
\mbox{\casilla TRZ-16}\newline {\color{media}\sffamily\bfseries\scriptsize MEDIA} & f. 39\newline{\scriptsize Trazabilidad} & El indicador con que se contrasta la cláusula (a) llama «módulo interactivo» a lo que el resto del documento llama «módulo educativo» —33 apariciones frente a 6—, de modo que el indicador y la cosa medida no comparten nombre; el tercer objetivo específico y el catálogo del Anexo B.7 usan «módulo educativo». & Usar «módulo educativo» en las seis apariciones de «módulo interactivo» (folios 39, 40, 47, 73, 91 y 109). Si se quiere subrayar el carácter interactivo, decirlo como atributo y no como nombre: «las siete primeras etapas tienen módulo educativo propio, de capa interactiva sobre el lienzo». & 20\\\addlinespace[3pt]
\mbox{\casilla APF-11}\newline {\color{media}\sffamily\bfseries\scriptsize MEDIA} & f. 39\newline{\scriptsize Aparato formal} & Dos letras cargan dos significados cada una sin que ningún pasaje lo advierta: N es a la vez la función de forma de la Nomenclatura y el número de subdivisiones por lado de las mallas de refinamiento en los Capítulos 2 y 3 y en los pies del Anexo D; p es a la vez el orden polinómico del elemento y el índice del punto de Gauss. El lector que sigue la tasa de convergencia O(h\textasciicircum{}(p+1)) sobre mallas N × N se encuentra con los dos choques a la vez. & Reservar N para la interpolación y escribir «mallas estructuradas de n\_\allowbreak{}el × n\_\allowbreak{}el elementos, con n\_\allowbreak{}el {\mates ∈} \{2, 4, 8, 16, 32\}\allowbreak{}»; si se prefiere conservar N por costumbre, añadir a la Nomenclatura «N — Número de subdivisiones por lado en las mallas estructuradas de refinamiento (no confundir con N\_\allowbreak{}i)» y repetir la advertencia en la… & 45\\\addlinespace[3pt]
\mbox{\casilla TRZ-17}\newline {\color{media}\sffamily\bfseries\scriptsize MEDIA} & f. 40\newline{\scriptsize Trazabilidad} & La versión evaluada se identifica solo como «software distribuido con este documento»: sin número de versión, fecha ni identificador de revisión, los veredictos de la tabla de especificaciones y de la matriz de principios no son reproducibles contra un estado congelado del software, y la promesa de reproducibilidad de 2.1.5 tampoco nombra los guiones. Matización respecto de los hallazgos brutos: la dirección del repositorio sí figura, en nota al pie del objetivo general (folio 5); lo que falta es la versión o etiqueta y la ubicación de los guiones. & Consignar en la fila «Versión evaluada» de la Tabla 2.1, y en el Anexo A, el número de versión, la fecha y el identificador de revisión del repositorio sobre los que se construyeron los instrumentos; en 2.1.5, completar la frase: «[...]\allowbreak{} con las dependencias declaradas en el repositorio del proyecto (nota 1 del folio 5), sobre… & 15\\\addlinespace[3pt]
\mbox{\casilla CON-08}\newline {\color{media}\sffamily\bfseries\scriptsize MEDIA} & f. 41\newline{\scriptsize Consistencia} & Las tres unidades de análisis se nombran «casos de estudio» en el diseño metodológico (5 apariciones), «casos de referencia» en el Capítulo 3 y en los anexos (17), «problemas de referencia» en 1.13, 3.7.1 y los Aportes (3) y «banco de prueba» al presentar la membrana de Cook (2), lo que obliga al tribunal a comprobar por su cuenta que se trata siempre de los mismos tres. & Nombre canónico metodológico: «casos de estudio» (son las unidades de análisis del Capítulo 2 y del Capítulo 3); reservar «caso de referencia» solo para nombrar la fuente externa contra la que se contrasta -solución analítica, modelo de SAP2000-. Sustituir «problemas de referencia» y «banco de prueba» por «casos de estudio».… & 35\\\addlinespace[3pt]
\mbox{\casilla TRZ-18}\newline {\color{media}\sffamily\bfseries\scriptsize MEDIA} & f. 43\newline{\scriptsize Trazabilidad} & El texto no dice si esos tres conceptos proceden de los 49 ítems del consenso o los añade el autor; la Tabla 3.7 los codifica FEM2, FEM8 y FEM9, es decir, sí proceden del consenso, pero el lector de 2.1.4 no puede saberlo y queda sin distinguir qué parte del universo tiene criterio externo. Es una remisión faltante, no un eslabón roto, por lo que se rebaja a media la severidad del hallazgo bruto. & Explicitar la procedencia en la misma oración: «[...]\allowbreak{} más los tres ítems de conceptos teóricos del canal del mismo consenso —FEM2, FEM8 y FEM9 en la Tabla 3.7 [3, p. 1163]\allowbreak{}—: el ensamblaje con restricciones y solución, la formulación isoparamétrica de elementos bidimensionales y la integración numérica de Gauss». & 10\\\addlinespace[3pt]
\mbox{\casilla CLA-16}\newline {\color{media}\sffamily\bfseries\scriptsize MEDIA} & f. 43\newline{\scriptsize Claridad} & El pronombre «él» puede leerse como el estudiante, como el criterio o como el artefacto, y la frase resulta contradictoria en su superficie: lo que «se observa» es justamente lo que todavía no se observa. Es una frase de calibración —la que acota lo que el trabajo no midió—, de modo que conviene que sea inequívoca. & «[...]\allowbreak{} su unidad de análisis es el artefacto. El criterio del estudiante no se mide aquí: su observación queda diferida a la prueba de campo que diseñan las recomendaciones, cuya muestra será dirigida —un curso completo, seleccionado por accesibilidad y no por procedimiento probabilístico [cita]\allowbreak{}—.» & 15\\\addlinespace[3pt]
\mbox{\casilla CLA-17}\newline {\color{media}\sffamily\bfseries\scriptsize MEDIA} & f. 43\newline{\scriptsize Claridad} & La misma subsección usa «población» en dos acepciones distintas —estudiantes y problemas de cálculo— y la segunda en condicional («sería»), lo que deja en duda si esa población existe como tal o es una figura retórica, justo donde el tribunal comprueba el molde metodológico. & Nombrar la segunda con otro término y en indicativo: «El universo de referencia del experimento numérico es el conjunto de problemas de elasticidad lineal plana —en tensión o deformación plana— resolubles con elementos cuadriláteros isoparamétricos. Los tres casos que siguen no son una muestra estadística de ese universo, sino… & 15\\\addlinespace[3pt]
\mbox{\casilla CLA-18}\newline {\color{media}\sffamily\bfseries\scriptsize MEDIA} & f. 47\newline{\scriptsize Claridad} & En la misma frase el canal de cálculo tiene siete etapas y nueve, sin decir que la memoria cubre dos pasos adicionales ni cuáles son, de modo que el criterio de aceptación de la cláusula (a) parece contradecirse a sí mismo. Además, el criterio «nueve de nueve» se fija sin que el lector disponga todavía de la lista: la enumeración llega recién en 3.6.1, junto al veredicto que ese criterio debía gobernar, y la enumeración del canal que hace la Introducción usa otros rótulos. & Fijar el recuento partiendo la frase en dos y usando los rótulos de 3.6.1: «Para la exposición: módulo interactivo para cada una de las siete etapas elementales y de ensamblaje del canal —mapeo isoparamétrico, jacobiano, matriz B, matriz D, rigidez elemental por cuadratura, fuerzas nodales equivalentes y ensamblaje (M1 a M7)—.… & 25\\\addlinespace[3pt]
\end{longtable}}
\subsection{Capítulo 3 · 3.6–3.8 Validez de contenido, interpretación e hipótesis}
{\small\color{gris}\sffamily 1 correcciones críticas o altas y 10 medias o bajas en esta sección · 3,9 h · más 6 frases de calibración (sección 3)}\par
{\small\rotulo{Primero, la calibración de la sección 3}CAL-10 (f. 76), CAL-11 (f. 80), CAL-12 (f. 80), CAL-13 (f. 81), CAL-14 (f. 84), CAL-15 (f. 84).\par}
\begin{ficha}{alta}{\casilla\textbf{CON-03}\quad\chip{alta}{alta}\quad §3.7.2 Validez de contenido y coherencia · folio 80 · PDF 95\hfill 15 min}
\rotulo{Dice}\cita{«las de ranking más alto en el consenso de expertos -planificación, mallado y verificación y validación [3, p. 1168]\allowbreak{}- son justamente las que EduFEM ejercita sobre casos de referencia y con la memoria» (folio 80, PDF 95; cita verificada literal con ubicar.py).}\par
\rotulo{Problema}La frase afirma que EduFEM ejercita las tres dimensiones de ranking más alto del consenso, incluida la planificación, pero la tabla de especificaciones no contiene ningún ítem de planificación -sus cinco dimensiones son mallado, post-proceso, conceptos fundamentales, verificación y validación, y conceptos teóricos- y las Conclusiones ya solo nombran dos. Es el eslabón que sostiene la cláusula (c) de la hipótesis: el tribunal que vaya de la afirmación a la tabla no encontrará respaldo para un tercio de lo prometido.\par
\rotulo{Corrección · corregir}Sustituir la frase de 3.7.2 por una que no prometa más de lo que la Tabla 3.7 contiene: «la tabla de especificaciones muestra que el software cubre las dieciocho destrezas y conceptos que definen el criterio sobre la respuesta, y que dos de las dimensiones de mayor ranking en el consenso de expertos -mallado, y verificación y validación [3, p. 1168]\allowbreak{}- son justamente las que EduFEM ejercita sobre casos de referencia y con la memoria, es decir, las que exigen que el estudiante haga algo con el modelo y juzgue el resultado.» Si el autor desea conservar la mención a la planificación, añadir a continuación una oración que la excluya de forma explícita y verificable: «La planificación, la tercera dimensión de ranking alto, no forma parte del universo de contenido de la Tabla 3.7 y queda fuera de lo que este trabajo contrasta.» Con cualquiera de las dos variantes, la frase de Conclusiones (folio 89) ya queda alineada y no necesita cambios.\par
\rotulo{También en}\begin{itemize}[leftmargin=1.2em,itemsep=1pt,topsep=1pt]
\item {\small Conclusiones, validez de contenido, f. 89: \cita{las de ranking mas alto en ese consenso}}
\item {\small Tabla 3.7, dimensiones del universo de contenido, f. 76: \cita{mallado, post-proceso, conceptos fundamentales, verificación y validación, conceptos teóricos}}
\end{itemize}
{\footnotesize\color{gris}\rotulo{Verificación}ajustado: Confirmado en el fondo, ajustada la propuesta. La cita es literal y está en PDF 95 (folio 80). Leí la Tabla 3.7 (PDF 91, folio 76) y su universo de contenido no tiene ningún ítem de planificación: arranca en FEA13 y se agrupa en Mallado (FEA13-FEA18), Post-proceso (FEA26-FEA28), Conceptos fundamentales (FEA29-FEA31) y Verificación y…\par}
\end{ficha}
{\small\rotulo{Párrafos que pide la defensa}\par}
\begin{ficha}{media}{\casilla\textbf{DEF-20}\quad\chip{media}{media}\quad Cap. 3 — cobertura del canal de cálculo · folio 75 · PDF 90\hfill 25 min}
\rotulo{Dice}\cita{El criterio de la Subsección 2.1.6 se cumple: siete de siete etapas con módulo y nueve de nueve en la memoria}\par
\rotulo{Problema}Riesgo de que el criterio se mida contra su propio diseño: la cuenta se hace contra una partición del canal de cálculo en nueve etapas que el propio trabajo adopta y a la que luego ajusta los módulos, y el texto no dice de dónde sale esa partición. Basta que el tribunal pregunte «¿quién fijó esas nueve etapas?» para que el resultado pierda valor probatorio.\par
\rotulo{Corrección · anadir}Completar la remisión en el punto de la cuenta, con localizador: «De las nueve etapas en que la literatura clásica presenta la formulación isoparamétrica [2, pp. --; 4, pp. --]\allowbreak{} —las mismas que desarrollan las Secciones 1.3 a 1.11—, siete tienen módulo interactivo y las nueve tienen desarrollo con sustitución numérica en la memoria de cálculo». Añadir que el criterio de la Subsección 2.1.6 se fijó sobre esa partición y que su falsación sería una etapa del canal sin exposición. No afirmar una cronología que el documento no acredita (que el criterio se fijó antes de construir los módulos).\par
\rotulo{También en}\begin{itemize}[leftmargin=1.2em,itemsep=1pt,topsep=1pt]
\item {\small Cap. 3 — interpretación, f. 78: \cita{siete de siete etapas con módulo y nueve de nueve}}
\end{itemize}
{\footnotesize\color{gris}\rotulo{Verificación}ajustado: La cita es literal y está donde el hallazgo la sitúa: PDF 90 /\allowbreak{} folio 75, con segunda aparición en PDF 93 /\allowbreak{} folio 78 (tabla de principios). El riesgo de circularidad es real —la partición y el criterio los fija el mismo trabajo—, pero la premisa del diagnóstico («el texto no dice de dónde sale esa partición») no se sostiene: la oración…\par}
\end{ficha}
\subsubsection*{Medias y bajas}
{\small\begin{longtable}{@{}L{1.6cm}L{1.6cm}L{5.25cm}L{5.65cm}r@{}}
\toprule
\sffamily\bfseries ID & \sffamily\bfseries Dónde & \sffamily\bfseries Qué pasa & \sffamily\bfseries Qué hacer & \sffamily\bfseries min\\
\midrule\endhead
\bottomrule\endfoot
\mbox{\casilla TRZ-23}\newline {\color{media}\sffamily\bfseries\scriptsize MEDIA} & f. 75\newline{\scriptsize Trazabilidad} & Compartir funciones garantiza consistencia entre el documento y el motor, no corrección de los valores, y «cada valor» es una afirmación exhaustiva que la prueba invocada en la oración anterior —que la memoria se generó «sin error de compilación»— no respalda; el argumento aparece en la sección que sostiene la cláusula (a). El hallazgo bruto citaba «Anexo F»: el texto impreso dice Anexo G. & Separar la prueba del argumento por construcción, sin prometer más de lo verificado: «La memoria de cálculo se generó para el ejemplo canónico en Q4 y en Q9 sin error de compilación (tests/\allowbreak{}test\_\allowbreak{}memoria\_\allowbreak{}calculo.py) y el Anexo G reproduce su secuencia de sustitución numérica. La memoria no recalcula por su cuenta: invoca las… & 15\\\addlinespace[3pt]
\mbox{\casilla APF-14}\newline {\color{media}\sffamily\bfseries\scriptsize MEDIA} & f. 76\newline{\scriptsize Aparato formal} & La columna de evidencia mezcla dos formas de remisión -el signo de sección abreviado y la palabra completa- en doce celdas frente a ocho, mientras la Tabla 3.6 de la misma sección y toda la prosa del documento usan siempre «Sección X.X» o «Subsección X.X.X»; el lector no especialista puede creer que apuntan a objetos de naturaleza distinta cuando ambas son secciones. & Usar una sola orden de referencia en toda la columna, la que ya produce «Sección 3.x» y «Tabla 3.x» en el resto del documento, y revisar las dieciocho filas. Si el ancho de columna obliga a abreviar, usar «Subsec. 2.2.X» y aplicarlo también a la Tabla 3.6. & 20\\\addlinespace[3pt]
\mbox{\casilla APF-18}\newline {\color{baja}\sffamily\bfseries\scriptsize BAJA} & f. 78\newline{\scriptsize Aparato formal} & Tres páginas quedan a medio llenar por colocación de flotantes o por cierre de sección. No hay error de contenido —son cierres legítimos—, pero en un documento de 168 páginas destinado a impresión desaprovechan espacio y, en el caso de la Tabla 3.8, separan una tabla clave del capítulo de resultados del párrafo que la introduce. & Ajustar la colocación del flotante de la Tabla 3.8 (parámetro de posición más permisivo o adelantando el párrafo introductorio) para que no quede media página en blanco antes de ella, y redistribuir el salto de página de los folios 85 y 95 adelantando o atrasando un párrafo de la sección vecina. No tocar los folios 117, 137 y… & 30\\\addlinespace[3pt]
\mbox{\casilla RED-17}\newline {\color{media}\sffamily\bfseries\scriptsize MEDIA} & f. 79\newline{\scriptsize Redundancia} & Dos oraciones completas viajan literalmente de §3.4 a §3.7.1 y de §3.7.1 a «Aportes». No se trata del recuento de cifras en sí —el Resumen, §3.8 y las Conclusiones lo recuerdan legítimamente—, sino de oraciones interpretativas copiadas, que hacen que la sección de interpretación repita el reporte en lugar de interpretarlo. & Reportar en §3.4 (tablas y explicación del peso relativo de las componentes secundarias), interpretar en §3.7.1 sin reproducir el reporte. Se elimina de §3.7.1 la oración «cuya magnitud es entre seis y treinta y cinco veces menor que la de σx», ya presente en §3.4, y de «Aportes» la oración «…por debajo del 0,3 \% frente a la… & 20\\\addlinespace[3pt]
\mbox{\casilla RED-18}\newline {\color{media}\sffamily\bfseries\scriptsize MEDIA} & f. 80\newline{\scriptsize Redundancia} & La lectura conjunta de los dos instrumentos documentales —construidos por vías distintas y coincidentes en las mismas decisiones: conversión de elemento, refinamiento, modo crudo y suavizado, equilibrio y diagnóstico de la memoria— se expone entera en §3.7.2 y se repone entera en las Conclusiones, con frases coincidentes. & Aparición canónica: §3.7.2 (folio 80), donde el análisis corresponde. En las Conclusiones reducir el pasaje a la conclusión sin su desarrollo: «Los dos instrumentos, construidos por vías independientes, convergen en el mismo conjunto de decisiones de diseño, que la sección 3.7.2 identifica como el núcleo didáctico de la… & 15\\\addlinespace[3pt]
\mbox{\casilla CLA-23}\newline {\color{media}\sffamily\bfseries\scriptsize MEDIA} & f. 82\newline{\scriptsize Claridad} & Dentro de una subsección titulada «Limitaciones observadas», el pasaje más largo sostiene lo contrario de lo que el título anuncia y lo hace en registro defensivo, respondiendo a una objeción que el lector nunca vio formulada. El tribunal puede leer en ese tono una debilidad que el propio texto está tratando de conjurar. & Llevar la Tabla 3.9 y sus dos párrafos de medición a los resultados del motor, donde la Subsección 2.1.6 los sitúa (el desempeño del solucionador se reporta sin criterio de aceptación), y dejar en Limitaciones solo lo que limita: «El almacenamiento disperso y la factorización directa mantienen el costo dentro del rango de uso… & 40\\\addlinespace[3pt]
\mbox{\casilla TRZ-24}\newline {\color{media}\sffamily\bfseries\scriptsize MEDIA} & f. 83\newline{\scriptsize Trazabilidad} & Tres cifras de mejora se presentan como medidas pero no figuran en la tabla de tiempos ni en ningún anexo, y «el mismo guion» no identifica qué guion ni qué corrida, de modo que el dato no es reproducible por el lector. El desempeño es el único indicador sin criterio declarado, lo que limita el alcance del defecto. & Añadir a la tabla de tiempos dos columnas (tiempo de solución con COLAMD y con el ordenamiento de mínimo grado) o una nota al pie que remita al guion concreto y a la corrida, con el mismo aviso de dispersión que ya lleva la leyenda de la tabla. & 30\\\addlinespace[3pt]
\mbox{\casilla TRZ-25}\newline {\color{media}\sffamily\bfseries\scriptsize MEDIA} & f. 84\newline{\scriptsize Trazabilidad} & Las tres preguntas de investigación se responden solo por remisión, dos de ellas son de sí o no y nunca reciben un sí explícito, y las Conclusiones no vuelven a mencionarlas: el eslabón pregunta {\mates →} respuesta queda enunciado pero no escrito. & Dar la respuesta en sus propios términos: «PI-1: la organización que expone el canal es la de tres fases sobre un único lienzo, con un módulo educativo por cada una de las siete etapas elementales y de ensamblaje, las nueve etapas desarrolladas en la memoria y la respuesta en el post-proceso con sus modos crudo y suavizado.… & 20\\\addlinespace[3pt]
\mbox{\casilla CLA-24}\newline {\color{media}\sffamily\bfseries\scriptsize MEDIA} & f. 84\newline{\scriptsize Claridad} & «Superan la ida y vuelta» traslada literalmente una expresión de laboratorio de software (la prueba de ida y vuelta) y no significa nada para un lector de ingeniería civil, justo en la oración donde se da por confirmada una de las tres cláusulas comprobables de la hipótesis. & «…y los formatos de intercambio (CSV/\allowbreak{}ZIP y DXF) recuperan sin pérdida el modelo exportado, según se comprobó en la Subsección 3.3.» Ajustar el número de subsección al del apartado de robustez y gestión del modelo. & 15\\\addlinespace[3pt]
\mbox{\casilla RED-21}\newline {\color{baja}\sffamily\bfseries\scriptsize BAJA} & f. 84\newline{\scriptsize Redundancia} & La repetición de la locución de enlace borra la diferencia entre responder las preguntas de investigación y responder el problema científico, que es justamente la jerarquía que la sección debe dejar clara. & Abrir el segundo párrafo con «Respondidas las tres preguntas, queda respondida también la parte del problema científico que esta etapa alcanza.» Cambio de una línea, sin efecto sobre la extensión. & 15\\\addlinespace[3pt]
\end{longtable}}
\subsection{Conclusiones y recomendaciones}
{\small\color{gris}\sffamily 1 correcciones críticas o altas y 14 medias o bajas en esta sección · 5,2 h · más 2 frases de calibración (sección 3)}\par
{\small\rotulo{Primero, la calibración de la sección 3}CAL-03 (f. 86), CAL-16 (f. 88).\par}
\begin{ficha}{alta}{\casilla\textbf{RED-03}\quad\chip{alta}{alta}\quad Conclusiones /\allowbreak{} Limitaciones · folio 92 · PDF 107\hfill 45 min}
\rotulo{Dice}\cita{PDF 107 /\allowbreak{} folio 92 (Conclusiones, «Limitaciones»): «en deformación plana la matriz constitutiva degenera hacia la incompresibilidad cuando nu -\textgreater{} 0,5, régimen que el módulo M4 señala pero que el software no resuelve. En cuarto lugar, el solucionador emplea una factorización directa [...]\allowbreak{} Finalmente, la validación descansa en tres casos de referencia y en un único caso con material de ingeniería civil; una batería más amplia reforzaría la confianza [...]\allowbreak{} y la carga superficial linealmente variable no interviene en ningún caso de referencia.» PDF 97 /\allowbreak{} folio 82 (§3.7.4): las mismas cuatro limitaciones con la misma secuencia -«el módulo M4 lo señala visualmente al aproximarse nu a ese límite, pero el software no lo resuelve»-. PDF 108 /\allowbreak{} folio 93 (Recomendaciones): «La evidencia del capítulo 3 descansa en tres casos de referencia y en un único caso con material de ingeniería civil, la viga de hormigón. Se recomienda incorporar un caso con carga superficial linealmente variable, que ningún caso de referencia ejercita». CUARTO locus no listado en el hallazgo, verificado con ubicar.py: PDF 59 /\allowbreak{} folio 44 (§2.1.4) ya declara «la carga superficial linealmente variable no interviene en ningún caso de referencia y no tiene contraste externo; las recomendaciones proponen un caso que la incorpore», de modo que esa limitación concreta se enuncia tres veces (folios 44, 92 y 93). En cambio PDF 24 /\allowbreak{} folio 9 NO es una repetición del inventario: enuncia fronteras a priori distintas («El software no aborda análisis tridimensional, comportamiento no lineal (material o geométrico), efectos dinámicos...»).}\par
\rotulo{Problema}El inventario de limitaciones se escribe entero tres veces —Introducción, §3.7.4 y Conclusiones— y una cuarta vez, en parte, al abrir las Recomendaciones, donde el estado actual se recapitula en páginas consecutivas con oraciones que coinciden palabra por palabra antes de enunciar la propuesta. Es la redundancia más extensa del cierre ({\mates ≈} una página) y además debilita el efecto de las limitaciones al diluirlas.\par
\rotulo{Corrección · fusionar}Repartir por función. La Introducción (folio 9) conserva intactas las fronteras fijadas a priori -sin 3D, sin no lineal, sin dinámica, sin mallado automático-, que no son el mismo contenido y que el molde exige; no se tocan. §3.7.4 (folio 82) conserva las limitaciones que el desarrollo puso de manifiesto: bloqueo del Q4 sin técnicas correctivas, degeneración de D cuando nu -\textgreater{} 0,5 con la señal visual de M4, coste de la factorización directa y batería de tres casos. «Limitaciones» de las Conclusiones (folio 92) se reduce a cinco viñetas de una línea, cada una con su remisión (§3.7.4, §1.9, §2.2.5), sin volver a explicar ninguna. La apertura de «Ampliación de la batería de validación» (folio 93) arranca directamente en la propuesta -«Se recomienda incorporar un caso con carga superficial linealmente variable, que ningún caso de referencia ejercita, y al menos un caso civil en deformación plana con solución de referencia publicada...»-, suprimiendo la recapitulación del estado actual. La declaración de §2.1.4 (folio 44) sobre esa carga se conserva: allí es la salvedad de cobertura del diseño de casos, no una recapitulación. Ahorro: media página larga.\par
\rotulo{También en}\begin{itemize}[leftmargin=1.2em,itemsep=1pt,topsep=1pt]
\item {\small Recomendaciones /\allowbreak{} Ampliación de la batería de validación, f. 93: \cita{La evidencia del capítulo 3 descansa en tres casos de referencia y en un único caso con material de ingeniería civil, la viga de hormigón}}
\item {\small Introducción /\allowbreak{} Alcance y limitaciones, f. 9: \cita{El software no aborda análisis tridimensional, comportamiento no lineal}}
\item {\small §3.7.4 Limitaciones que el desarrollo puso de manifiesto, f. 82: \cita{el módulo M4 lo señala visualmente al aproximarse}}
\end{itemize}
{\footnotesize\color{gris}\rotulo{Verificación}ajustado: El núcleo del hallazgo está confirmado y es el más sólido del lote: «Limitaciones» de las Conclusiones (PDF 107 /\allowbreak{} folio 92) reescribe §3.7.4 (PDF 97 /\allowbreak{} folio 82) casi entera -bloqueo del Q4 sin técnicas correctivas, degeneración de la matriz constitutiva cuando nu tiende a 0,5 con la señal visual de M4, coste de la factorización directa,…\par}
\end{ficha}
{\small\rotulo{Párrafos que pide la defensa}\par}
\begin{ficha}{alta}{\casilla\textbf{DEF-06}\quad\chip{alta}{alta}\quad Limitaciones (Conclusiones) · folio 91 · PDF 106\hfill 30 min}
\rotulo{Dice}\cita{«validez de contenido no es efectividad: los instrumentos establecen que el software expone lo que el criterio sobre la respuesta exige, no que el estudiante lo adquiera» (verificada literal en PDF 106 /\allowbreak{} folio 91).}\par
\rotulo{Problema}La pregunta con que un docente con interés pedagógico suele cerrar la defensa: «¿qué impide que el estudiante use EduFEM como otra caja negra, más lenta que SAP2000?». El documento no aborda en ningún punto el riesgo más obvio de su propia propuesta —que el alumno recorra el pre-proceso y el post-proceso sin abrir un solo módulo—, y es la objeción que desactiva el argumento central desde dentro.\par
\rotulo{Corrección · anadir}Añadir a las limitaciones que la exposición del canal es una posibilidad que el software ofrece y no una obligación que imponga: un estudiante puede construir, resolver y leer el contorno sin abrir ningún módulo ni generar la memoria; la herramienta no fuerza el recorrido, decisión deliberada porque el orden de exploración lo fija la enseñanza y no el programa, de modo que su efecto depende de la consigna docente; los archivos de proyecto y las memorias que cada alumno produce permiten verificar si el recorrido se hizo, y la prueba de campo incorpora esa consigna como parte del tratamiento.\par
\rotulo{También en}\begin{itemize}[leftmargin=1.2em,itemsep=1pt,topsep=1pt]
\item {\small Recomendaciones, Prueba de campo, f. 94: \cita{El tratamiento reproduce el uso previsto de la herramienta}}
\end{itemize}
\end{ficha}
\begin{ficha}{media}{\casilla\textbf{DEF-21}\quad\chip{media}{media}\quad Aportes del trabajo /\allowbreak{} Limitaciones · folio 90-91 · PDF 105\hfill 30 min}
\rotulo{Dice}\cita{«El principal aporte de este trabajo es una herramienta de software libre, funcional y verificada» (verificada literal en PDF 105 /\allowbreak{} folio 90).}\par
\rotulo{Problema}Dos preguntas que se siguen naturalmente de «lo construyó y lo evaluó el autor»: «¿nadie más que usted ha usado el programa?» y «¿quién revisó su código y sus criterios?». Ni los aportes ni las limitaciones mencionan que el software no fue usado ni evaluado por nadie ajeno al autor —ni un docente, ni un estudiante, ni una prueba de usabilidad— y que el motor, los guiones de verificación y los umbrales los escribió la misma persona, sin revisión de código independiente; la afirmación de que la herramienta es «funcional» descansa por entero en quien la escribió.\par
\rotulo{Corrección · anadir}Añadir una frase a las Limitaciones, junto a las ya declaradas: «La herramienta fue ejercitada por el autor y por su batería automática de pruebas de regresión; no la operaron estudiantes ni docentes ajenos al trabajo, de modo que su usabilidad —tiempo de aprendizaje, errores de operación, claridad de los mensajes— no está establecida y forma parte de lo que la prueba de campo debe observar. El motor, los guiones de verificación y los umbrales los escribió el autor, sin revisión de código independiente; la contrastación externa se apoya en la solución analítica de Timoshenko-Goodier y en el modelo equivalente en SAP2000». No incorporar ensayos informales de última hora: añadirían al documento una afirmación no documentada.\par
\rotulo{También en}\begin{itemize}[leftmargin=1.2em,itemsep=1pt,topsep=1pt]
\item {\small Limitaciones, f. 91: \cita{validez de contenido no es efectividad: los instrumentos establecen}}
\item {\small §2.1.5, f. 46: \cita{su limitación, que se declara, es que las construye el autor}}
\end{itemize}
{\footnotesize\color{gris}\rotulo{Verificación}ajustado: La cita es literal y está en PDF 105 /\allowbreak{} folio 90. Buena parte del diagnóstico se refuta: la autoría única de los instrumentos está declarada dos veces —2.1.5, «su limitación, que se declara, es que las construye el autor» (PDF 61 /\allowbreak{} folio 46), y Limitaciones, «las tablas las construye el autor sobre un criterio externo publicado, no un…\par}
\end{ficha}
\begin{ficha}{media}{\casilla\textbf{DEF-22}\quad\chip{media}{media}\quad Recomendaciones — Prueba de campo del efecto sobre el criterio · folio 94 · PDF 109\hfill 20 min}
\rotulo{Dice}\cita{El diseño es de preprueba y posprueba con un solo grupo, o con dos cohortes si el calendario académico lo permite, en la línea de la comparación entre cohortes con que Bishay midió el efecto de sus herramientas [11, pp. 4-5, 13]\allowbreak{}.}\par
\rotulo{Problema}Todas las salvedades de la tesis remiten a esta prueba, pero con un solo grupo la variable independiente tiene un único nivel —exposición transparente— y la ganancia no puede atribuirse a la exposición del canal ni contrastarse con la caja negra: es exactamente el reparo que la propia tesis hace a Lee en 1.2, de modo que el tribunal puede devolverle su propio argumento.\par
\rotulo{Corrección · reformular}Invertir el orden de preferencia dentro de esa misma oración y añadir el aviso de atribución: «El diseño preferente es cuasiexperimental con grupo de comparación —dos paralelos o dos cohortes de la misma asignatura, uno con EduFEM y otro con el software habitual, con preprueba y posprueba en ambos—, en la línea de la comparación entre cohortes con que Bishay midió el efecto de sus herramientas [11, pp. 4-5, 13]\allowbreak{}; si el calendario solo permite un grupo, la ganancia medida no podrá atribuirse a la exposición del canal y así deberá reportarse».\par
\rotulo{También en}\begin{itemize}[leftmargin=1.2em,itemsep=1pt,topsep=1pt]
\item {\small §1.2 (el mismo reparo, aplicado a Lee), f. 18: \cita{con una evaluación cualitativa sobre sus propios cursos, sin grupo de control}}
\end{itemize}
{\footnotesize\color{gris}\rotulo{Verificación}ajustado: La cita es literal y está en PDF 109 /\allowbreak{} folio 94, pero está cortada justo antes de la salvedad que el hallazgo echa de menos: la oración completa dice «El diseño es de preprueba y posprueba con un solo grupo, o con dos cohortes si el calendario académico lo permite, en la línea de la comparación entre cohortes con que Bishay midió el…\par}
\end{ficha}
\begin{ficha}{media}{\casilla\textbf{DEF-23}\quad\chip{media}{media}\quad Recomendaciones — Prueba de campo del efecto sobre el criterio · folio 94 · PDF 109\hfill 60 min}
\rotulo{Dice}\cita{«El análisis compara la preprueba y la posprueba con estadística no paramétrica para muestras pareadas, adecuada al tamaño de un curso» (verificada literal en PDF 109 /\allowbreak{} folio 94).}\par
\rotulo{Problema}El protocolo al que la tesis difiere su afirmación central tiene población, muestra, diseño, tratamiento, instrumento y análisis, pero le faltan los seis elementos que un tribunal de metodología exige: hipótesis estadística, tamaño de muestra y potencia, validación del instrumento, criterio de éxito fijado a priori, consideraciones éticas y plazo. Si el diseño diferido está incompleto, la estrategia de diferir la comprobación pierde fuerza.\par
\rotulo{Corrección · anadir}Cerrar el apartado con un párrafo que complete el protocolo sin anticipar su resultado: hipótesis estadística (H0: la mediana de las diferencias posprueba-preprueba es nula; H1: es mayor que cero), contrastada con la prueba de Wilcoxon de rangos con signo al 5 \% e informando el tamaño del efecto; tamaño de muestra del orden de un curso, con la advertencia de que esa n solo permite detectar efectos medianos o grandes; validación previa del instrumento por juicio de docentes de la asignatura y una aplicación piloto; criterio de éxito fijado a priori, antes de aplicar la posprueba; participación voluntaria con consentimiento informado, calificación desligada de la investigación y anonimización de los datos; y la ventana académica en que se aplicaría.\par
{\footnotesize\color{gris}\rotulo{Verificación}ajustado: La cita es literal y está en PDF 109 /\allowbreak{} folio 94, y la ausencia se confirma: la búsqueda en las 168 páginas no encuentra «Wilcoxon», «consentimiento», «potencia» estadística ni «tamaño de muestra», así que los seis elementos que enumera el hallazgo no están en el documento. Pero el apartado es una recomendación y declara su propio alcance…\par}
\end{ficha}
\begin{ficha}{media}{\casilla\textbf{DEF-24}\quad\chip{media}{media}\quad Conclusiones y recomendaciones · Recomendaciones y trabajo futuro (punto de inserción: entre «Sobre el software» y «Sobre la evidencia»); la evidencia citada está en «Sobre la evidencia» · ítem «Prueba de campo del efecto sobre el criterio». Cita secundaria corregida: tabla de la sección 3.6 (no 3.7.2), folio 79 /\allowbreak{} PDF 94 · folio 94 · PDF 109\hfill 35 min}
\rotulo{Dice}\cita{«El tratamiento reproduce el uso previsto de la herramienta: construir el ejemplo canónico, validar su salud, resolverlo, recorrer los módulos M1 a M7 sobre un elemento propio, generar la memoria y reproducir a mano una etapa, y convertir el modelo de Q4 a Q9 y refinar la malla sobre la membrana de Cook» (folio 94 /\allowbreak{} PDF 109). La frase que encuadra la sección y que la propuesta debe respetar es: «A partir de las limitaciones anteriores se plantean tres líneas de continuación, una por capítulo, cada una con sus ítems concretos» (folio 93 /\allowbreak{} PDF 108).}\par
\rotulo{Problema}Pregunta natural sobre un software educativo: «¿cómo se usaría esto en la cátedra? ¿tiene una guía o una secuencia de prácticas para el docente?». El único lugar donde el documento describe cómo se usaría EduFEM en clase es el tratamiento de una prueba de campo futura, de modo que el trabajo entrega una herramienta a la cátedra sin secuencia didáctica ni guía docente; el manual del Anexo B enseña a operar el programa, no a enseñar con él.\par
\rotulo{Corrección · anadir}Abrir una cuarta subsección breve, «Sobre el uso en la cátedra», entre «Sobre el software» y «Sobre la evidencia», y ajustar la frase marco de «tres líneas de continuación, una por capítulo» a «cuatro líneas de continuación: una por capítulo y una sobre el uso del software en la cátedra». En ella, un único ítem —«Guía de prácticas para la cátedra»— que se limite a reordenar como plan de aula el material que el trabajo ya produjo, sin atribuir efectos: «El presente trabajo entrega la herramienta y su manual de operación (Anexo B), pero no un material didáctico que oriente su uso en la asignatura. Se recomienda redactar una guía de prácticas de laboratorio que organice en tres sesiones la secuencia de uso previsto que este trabajo ya especifica como tratamiento de la prueba de campo: (i) construcción del modelo, validación de salud y lectura de la deformada y de las reacciones; (ii) recorrido de los módulos M1 a M7 sobre un elemento del modelo propio y reproducción a mano, con la memoria de cálculo a la vista, de una de sus etapas; (iii) conversión de Q4 a Q9 y refinamiento sobre la membrana de Cook, con contraste de tensiones crudas frente a suavizadas. Cada práctica debería fijar su consigna, el archivo .edufem y la memoria en PDF como producto entregable, y apoyarse en las capturas del Anexo B. Esa guía es además la condición práctica para aplicar la prueba de campo que se describe a continuación, ya que fija el tratamiento en términos operativos para el docente; su eficacia didáctica, como la de la herramienta, queda sujeta a esa prueba».\par
\rotulo{También en}\begin{itemize}[leftmargin=1.2em,itemsep=1pt,topsep=1pt]
\item {\small §3.7.2, f. 77: \cita{se ejercitan sobre casos de referencia con secuencias de mallas}}
\end{itemize}
{\footnotesize\color{gris}\rotulo{Verificación}ajustado: El hecho de base se verifica: la evidencia es literal y está donde dice (PDF 109 /\allowbreak{} folio 94, subsección «Sobre la evidencia», ítem «Prueba de campo del efecto sobre el criterio»), y en todo el documento no existe ninguna guía docente ni secuencia de prácticas: las recomendaciones se organizan en tres líneas («Sobre la formulación»,…\par}
\end{ficha}
\subsubsection*{Medias y bajas}
{\small\begin{longtable}{@{}L{1.6cm}L{1.6cm}L{5.25cm}L{5.65cm}r@{}}
\toprule
\sffamily\bfseries ID & \sffamily\bfseries Dónde & \sffamily\bfseries Qué pasa & \sffamily\bfseries Qué hacer & \sffamily\bfseries min\\
\midrule\endhead
\bottomrule\endfoot
\mbox{\casilla CON-14}\newline {\color{media}\sffamily\bfseries\scriptsize MEDIA} & f. 86\newline{\scriptsize Consistencia} & La Introducción declara expresamente que su enunciado «es la única formulación del problema que el documento emplea», y las Conclusiones lo reescriben: suprimen la población y la gestión, cambian «software de análisis estructural» por «software» y, sobre todo, pasan del condicional «podrá afectar» al indicativo «afecta», que da por existente un efecto que este trabajo no midió. El Resumen añade además a la hipótesis una cláusula -«sobre el modelo del propio estudiante»- ausente del enunciado canónico. Con tres versiones de las dos oraciones más citadas de la tesis, el tribunal no puede comprobar la promesa de formulación única. & Un solo cambio, en las Conclusiones: citar el problema textualmente o marcarlo como paráfrasis conservando el condicional, «y a la interrogante que planteó como problema -de qué manera el uso del software de análisis estructural como caja negra podrá afectar en el bajo criterio para interpretar la respuesta estructural obtenida… & 20\\\addlinespace[3pt]
\mbox{\casilla TRZ-26}\newline {\color{media}\sffamily\bfseries\scriptsize MEDIA} & f. 87\newline{\scriptsize Trazabilidad} & El producto del primer objetivo -los atributos que la herramienta debe tener- no tiene una lista canónica: el cierre de 1.1 enumera licencia libre, interfaz en español, memoria trazable, fenómenos numéricos observables y verificación publicada, y la conclusión enumera otros cuatro, con solo dos en común, de modo que no se puede comprobar que el software cumpla lo que el diagnóstico pidió. & Fijar al final de 1.1 una lista numerada única de atributos, reutilizarla literalmente como requisitos en 2.2.1 y citarla por número en la conclusión del primer objetivo. Queda por comprobar la lista de requisitos de 2.2.1, que el agente que reportó el hallazgo no llegó a leer. & 25\\\addlinespace[3pt]
\mbox{\casilla CLA-25}\newline {\color{media}\sffamily\bfseries\scriptsize MEDIA} & f. 88\newline{\scriptsize Claridad} & «Empíricamente» sugiere evidencia experimental a diez líneas de la frase que declara que no hay mediciones físicas; la evidencia es numérica. Una sola palabra, pero en Conclusiones y en contradicción aparente con la salvedad de alcance. & Sustituir «ilustrar empíricamente» por «ilustrar numéricamente». & 5\\\addlinespace[3pt]
\mbox{\casilla ACC-11}\newline {\color{media}\sffamily\bfseries\scriptsize MEDIA} & f. 88-89\newline{\scriptsize Accesibilidad} & «Los cuatro principios de aprendizaje» se invocan cuatro veces entre el Resumen y las Conclusiones sin enumerarse nunca en esos capítulos, de modo que el lector que solo lea el Resumen y las Conclusiones —lectura habitual en un tribunal— no puede juzgar la afirmación que sostiene la cláusula de coherencia del diseño. & Enumerarlos en la primera mención de cada capítulo, entre rayas o entre paréntesis, con el nombre exacto que emplea el capítulo 1, y remitir a la sección donde se fundamentan; en las menciones siguientes basta la remisión. & 20\\\addlinespace[3pt]
\mbox{\casilla CON-15}\newline {\color{media}\sffamily\bfseries\scriptsize MEDIA} & f. 88\newline{\scriptsize Consistencia} & Las Conclusiones extienden a todas las reacciones el «a precisión de máquina» que el cuerpo reserva a la suma horizontal, y suprimen el residuo relativo de 1,7 x 10\textasciicircum{}-13, que es el dato medido de la componente vertical y el único que el lector puede contrastar con el cuerpo. & Sustituir por: «y la suma de las reacciones verticales equilibra la carga total con un residuo relativo de 1,7 x 10\textasciicircum{}-13, mientras la suma horizontal es nula a precisión de máquina». & 10\\\addlinespace[3pt]
\mbox{\casilla CON-16}\newline {\color{media}\sffamily\bfseries\scriptsize MEDIA} & f. 90\newline{\scriptsize Consistencia} & Dos defectos en la misma oración. Primero, la cota se extiende a «los problemas de referencia estudiados» sin acotarla a la viga de Timoshenko, y la membrana de Cook -uno de esos problemas- da para la magnitud primaria, el desplazamiento u\_\allowbreak{}y, un error del -0,594 \% con Q4 y 2178 GDL, y de más del 50 \% en la malla gruesa: leída literalmente, la frase la contradice la propia Tabla 3.4. El folio 80 arrastra la misma generalización, aunque al menos acota las magnitudes. Segundo, llama «validación» a lo que la página anterior llama «contrastación» tras declarar que la validación en sentido estricto queda fuera del alcance, con lo que el lector que solo lea Aportes se lleva una palabra más fuerte que la que el trabajo se concede. & Escribir en Aportes del trabajo: «La validación -en la acepción declarada en 1.13: contraste con referencias externas de exactitud conocida- acota el error de la tensión normal y de la flecha por debajo del 0,3 \% frente a la solución analítica y del 0,6 \% frente al modelo equivalente en SAP2000, en la viga de Timoshenko, que es… & 20\\\addlinespace[3pt]
\mbox{\casilla ACC-12}\newline {\color{media}\sffamily\bfseries\scriptsize MEDIA} & f. 91\newline{\scriptsize Accesibilidad} & Es la única mención individual de un módulo en todo el capítulo y aparece sin decir qué es M4, cuando los párrafos anteriores solo citaron el bloque «M1 a M7» y «M0». & «...régimen que el módulo de la matriz constitutiva (M4) señala con una advertencia explícita, pero que el software no resuelve.» & 10\\\addlinespace[3pt]
\mbox{\casilla CLA-26}\newline {\color{media}\sffamily\bfseries\scriptsize MEDIA} & f. 92\newline{\scriptsize Claridad} & «De forma desfavorable» y «mallas masivas» no dan ni el orden de crecimiento ni el umbral a partir del cual el esquema deja de servir, que es lo único accionable de la limitación y lo que el tribunal preguntará. & «…pero cuyo coste en memoria y tiempo crece más que linealmente con el número de grados de libertad, de modo que por encima del orden de magnitud verificado en el Capítulo 3 (Tabla 3.9) el esquema directo deja de ser practicable.» Si la Tabla 3.9 permite fijar la cifra, escribirla. & 15\\\addlinespace[3pt]
\mbox{\casilla TRZ-27}\newline {\color{media}\sffamily\bfseries\scriptsize MEDIA} & f. 93\newline{\scriptsize Trazabilidad} & Limitaciones y recomendaciones no se corresponden en tres puntos: se recomienda un caso externo en deformación plana sin haber declarado como limitación que ese estado solo se verifica con el MMS (Timoshenko y Cook son de tensión plana); la limitación de que las tablas las construye el autor sin panel independiente no tiene recomendación que la atienda; y el origen del criterio externo —expertos mexicanos de ingeniería mecánica— no se declara como límite de transferibilidad. & En Limitaciones, añadir: «la deformación plana solo se verifica con el método de soluciones manufacturadas: los dos casos de contraste externo son de tensión plana, y la viga de Timoshenko se resolvió únicamente con Q9», y una frase sobre el perfil del panel de expertos del que procede el criterio externo. En Recomendaciones,… & 15\\\addlinespace[3pt]
\mbox{\casilla CLA-27}\newline {\color{media}\sffamily\bfseries\scriptsize MEDIA} & f. 94\newline{\scriptsize Claridad} & El protocolo completo —población, muestra, diseño, tratamiento, instrumento, análisis— se entrega como un solo bloque de prosa. Como es la prueba a la que la tesis remite el efecto sobre el estudiante, el tribunal necesita auditarlo elemento por elemento, y en este formato no puede hacerlo ni advertir cuáles elementos faltan. & Convertir el párrafo en una lista rotulada con una entrada por elemento —Población, Muestra, Diseño, Tratamiento, Instrumento, Análisis, Criterio de éxito, Consideraciones éticas, Cronograma—, sin cambiar el contenido. El formato deja a la vista los elementos que hoy no están y refuerza que el efecto queda deferido a esa prueba. & 30\\\addlinespace[3pt]
\mbox{\casilla REF-04}\newline {\color{media}\sffamily\bfseries\scriptsize MEDIA} & f. 94\newline{\scriptsize Referencias} & Cinco claves del .bib llevan un año distinto del que imprime la entrada. La clave no llega al lector, de modo que el PDF no muestra la discrepancia, pero sí señala el riesgo que sí es verificable: que los localizadores de página se hayan tomado de una edición distinta de la que la lista declara. Vancouver liga el localizador a la edición consignada, y estas cinco entradas concentran citas paginadas de precisión fina ([1, pp. 69-71]\allowbreak{}, [15, art. 21, pp. 39-43]\allowbreak{}, [24, pp. 106-107]\allowbreak{}, [16, pp. 232-234]\allowbreak{}, [4, pp. 476-477]\allowbreak{}) cuyo cotejo depende de acertar el ejemplar. & Comprobar, para las cinco entradas, que las páginas citadas correspondan al ejemplar declarado en la lista (la carpeta tesis/\allowbreak{}respaldo\_\allowbreak{}citas ya guarda ese cotejo para otras fuentes) y, hecho eso, renombrar las claves al año de la entrada —zienkiewicz2005fem, bathe1996fem, hughes1987fem, timoshenko1951elasticity,… & 40\\\addlinespace[3pt]
\mbox{\casilla REF-05}\newline {\color{media}\sffamily\bfseries\scriptsize MEDIA} & f. 94\newline{\scriptsize Referencias} & La fuente que sostiene el marco metodológico completo —objeto y campo, planteamiento del problema, hipótesis y sistema de variables, es decir, justo lo que el tribunal examina con más detalle— es un documento inédito sin pie de imprenta ni fecha, citado cuatro veces con páginas exactas. La declaración entre corchetes es correcta en Vancouver/\allowbreak{}NLM y la tesis es transparente al respecto, pero el lector no dispone de ninguna vía para recuperar el ejemplar y comprobar las páginas 7, 8, 13 y 21. & Mantener la entrada tal como está y añadir a la nota una vía de recuperación concreta (repositorio o cátedra de la Carrera de Ingeniería Civil de la UATF que lo distribuye, o «ejemplar en poder del autor, disponible a solicitud»). En paralelo, respaldar las cuatro afirmaciones que hoy dependen solo de esta fuente con las dos… & 30\\\addlinespace[3pt]
\mbox{\casilla CON-22}\newline {\color{baja}\sffamily\bfseries\scriptsize BAJA} & f. 94\newline{\scriptsize Consistencia} & Higiene dispar de campos entre entradas del mismo tipo: los dos informes técnicos de Sandia se tratan de forma distinta; el manual de CSI, que se distribuye en línea, no lleva URL ni fecha de consulta que permitan localizar la revisión 18; y el ISBN de Reddy está mal formado («007-124473-5» por «0-07-124473-5»). Ninguno de los tres altera lo que el lector ve en la lista compuesta —Vancouver no imprime el ISBN y la lista ya identifica cada obra—, pero sí la uniformidad del archivo de referencias. & Añadir a salari2000mms el DOI y la URL de OSTI con su fecha de consulta, con el mismo formato y el mismo comentario de verificación que lleva stimpson2007verdict; fijar un criterio único para los recursos en línea (o URL con fecha de consulta en todos, o en ninguno) y aplicarlo también a csi2017sap2000; corregir el ISBN de… & 15\\\addlinespace[3pt]
\mbox{\casilla CLA-29}\newline {\color{baja}\sffamily\bfseries\scriptsize BAJA} & f. 95\newline{\scriptsize Claridad} & Cierre retórico que no añade información y termina el documento con una metáfora en lugar de con el compromiso concreto de la prueba de campo, que es lo que el tribunal debe retener sobre el alcance de lo comprobado. & «De las tres líneas, la prueba de campo es la que decide el alcance del trabajo: las dos primeras amplían el artefacto; solo la tercera puede establecer si la exposición del canal mejora el criterio con que el estudiante interpreta la respuesta estructural.» & 10\\\addlinespace[3pt]
\end{longtable}}
\subsection{Resumen}
{\small\color{gris}\sffamily 2 correcciones críticas o altas y 4 medias o bajas en esta sección · 2,4 h · más 0 frases de calibración (sección 3)}\par
\begin{ficha}{critica}{\casilla\textbf{TRZ-01}\quad\chip{critica}{crítica}\quad Resumen · folio i · PDF 2\quad{\color{gris}\footnotesize también CON-01, CLA-02}\hfill 10 min}
\rotulo{Dice}\cita{la viga de Timoshenko arrojó errores del 0,04 \% en la tensión normal y del 0,26 \% en la flecha frente a la solución analítica, y diferencias de hasta el 0,56 \% frente a SAP2000}\par
\rotulo{Problema}El Resumen enuncia «hasta el 0,56 \%» como cota máxima de la diferencia frente a SAP2000 sin acotarla a la magnitud medida, mientras que la Tabla 3.2 registra diferencias frente a SAP2000 del 4,56 \% en σy y del 2,72 \% en τxy, y 3.8 sí acota la cifra («en desplazamientos»); quien lea solo el Resumen se lleva una cota de exactitud ocho veces menor que la mayor diferencia medida, y es la única afirmación del Resumen que la propia tabla que la respalda desmiente.\par
\rotulo{Corrección · reformular}«[...]\allowbreak{} y diferencias de hasta el 0,21 \% en la tensión normal y del 0,56 \% en el desplazamiento frente a SAP2000; en las componentes secundarias de tensión, de magnitud entre seis y treinta y cinco veces menor, las diferencias frente a SAP2000 llegan al 4,56 \%.»\par
\rotulo{También en}\begin{itemize}[leftmargin=1.2em,itemsep=1pt,topsep=1pt]
\item {\small §3.4 y Tabla 3.2, f. 68: \cita{en el punto B, 0,94 \% en σy y 2,89 \% en τxy frente a la solución analítica, y 4,56 \% y 2,72 \% frente a SAP2000}}
\item {\small §3.8 (cláusula b), f. 84: \cita{y de hasta el 0,21 \% en σx y el 0,56 \% en desplazamientos frente a SAP2000}}
\end{itemize}
{\small\color{gris}\rotulo{CON-01}El Resumen da el 0,56 \% como diferencia máxima frente a SAP2000 sin el calificador «en desplazamientos» que sí llevan las Conclusiones (folio 88) y sin la salvedad de las componentes secundarias que sí lleva 3.7.1 (folio 80), mientras la Tabla 3.2 reporta frente a la misma referencia 4,56 \% en sigma\_\allowbreak{}y y 2,72 \% en tau\_\allowbreak{}xy: leída al pie de la letra, la cifra de cierre del Resumen -la primera que lee el tribunal y la única que muchos leerán- es falsa. Es el único lugar del documento donde la cifra aparece sin acotar; los demás ya están bien calibrados, de modo que basta importar su redacción.\par}
{\small\color{gris}\rotulo{CLA-02}Los tres casos de verificación y validación —método de soluciones manufacturadas, viga de Timoshenko y membrana de Cook— se apilan en una sola oración de unas 140 palabras encadenada con dos puntos y punto y coma, de modo que se pierde que son tres experimentos distintos con conclusiones distintas; ocurre en la página que el tribunal lee primero y a veces única.\par}
{\footnotesize\color{gris}\rotulo{Verificación}ajustado: Las tres citas son literales y están en sus páginas (Resumen, PDF 2; Tabla 3.2 y su comentario, PDF 83, folio 68; 3.8, PDF 99, folio 84). Pero el diagnóstico sobrecarga el defecto: la oración del Resumen enumera dos magnitudes —«errores del 0,04 \% en la tensión normal y del 0,26 \% en la flecha frente a la solución analítica»— y la cota…\par}
\end{ficha}
\begin{ficha}{alta}{\casilla\textbf{ACC-01}\quad\chip{alta}{alta}\quad Resumen · folio i · PDF 2\hfill 30 min}
\rotulo{Dice}\cita{«se confirmaron las tasas de convergencia asintóticas teóricas del desplazamiento —orden dos en norma L2 y orden uno en seminorma H1 para Q4; tres y dos para Q9— y la convergencia del campo de tensiones recuperado» (PDF 2, literal); cierre: «las quince destrezas y los tres conceptos del universo tienen instrumento» y «cada uno de los cuatro principios de aprendizaje que fundamentan el diseño tiene una encarnación comprobable» (PDF 3).}\par
\rotulo{Problema}El Resumen es la única página que el tribunal completo lee entera y enuncia el resultado central de verificación en un formalismo (norma L2, seminorma H1) del que no se dice qué mide ni por qué importa; el cierre añade «universo» y «encarnación comprobable», vocabulario de la teoría de la medición educativa que no es terminología académica reconocible para un lector de ingeniería civil. La página más leída del documento resulta así la menos legible.\par
\rotulo{Corrección · reformular}Decir la magnitud antes que el nombre técnico: «se confirmaron las tasas de convergencia que la teoría predice: al reducir a la mitad el tamaño de los elementos, el error del desplazamiento se divide por cuatro y el de su gradiente por dos con el Q4, y por ocho y por cuatro con el Q9 -órdenes dos y uno, y tres y dos, medidos en la norma L2 del desplazamiento y la seminorma H1 de su gradiente-; también converge el campo de tensiones recuperado». En el cierre, conservar «universo» (es el término de la tabla de especificaciones y reaparece en la subsección 3.6.2) pero glosarlo en su primera aparición: «los tres conceptos del universo de contenido -el conjunto de destrezas y conceptos que el consenso de expertos fija como referencia- tienen instrumento en EduFEM»; sustituir «encarnación comprobable» por «realización identificable en el software», y glosar el instrumento en aposición: «una tabla de especificaciones -el instrumento que confronta, ítem por ítem, lo que el software ofrece con lo que un criterio externo exige-». Si se adopta el cambio de «encarnación comprobable», replicarlo en PDF 99 /\allowbreak{} folio 84 para no dejar dos formulaciones distintas.\par
\rotulo{También en}\begin{itemize}[leftmargin=1.2em,itemsep=1pt,topsep=1pt]
\item {\small Resumen (cierre), f. ii: \cita{las quince destrezas y los tres conceptos del universo tienen instrumento}}
\end{itemize}
{\footnotesize\color{gris}\rotulo{Verificación}ajustado: Evidencia verificada y literal en su página: el pasaje de las tasas está en el Resumen (PDF 2 /\allowbreak{} folio i) y el cierre «las quince destrezas y los tres conceptos del universo tienen instrumento en EduFEM, y cada uno de los cuatro principios de aprendizaje que fundamentan el diseño tiene una encarnación comprobable» en PDF 3 /\allowbreak{} folio ii. No…\par}
\end{ficha}
{\small\rotulo{Párrafos que pide la defensa}\par}
\begin{ficha}{alta}{\casilla\textbf{DEF-01}\quad\chip{alta}{alta}\quad Resumen (y 2.1.5, Introducción, Limitaciones) · folio ii · PDF 3\hfill 45 min}
\rotulo{Dice}\cita{«las destrezas de interpretación y verificación de resultados que un consenso publicado de 67 expertos exige al egresado» (verificada literal en PDF 3 /\allowbreak{} folio ii); la salvedad, «los ítems que se toman son contenido universal del método», en PDF 60 /\allowbreak{} folio 45.}\par
\rotulo{Problema}Pregunta previsible: «sus 67 expertos son mexicanos y de ingeniería mecánica, y el propio estudio dice que no son generalizables; ¿por qué valen para ingeniería civil en Bolivia?». El criterio externo, único sustento de la validez de contenido de todo el trabajo, se presenta en el Resumen, la Introducción y las Conclusiones como lo que «los expertos exigen al egresado»; la salvedad aparece una sola vez, en 2.1.5, donde se resuelve con una aserción del autor que no se apoya en nada, y no llega a las Limitaciones. El Resumen tampoco dice que la tabla la construye el propio autor.\par
\rotulo{Corrección · reformular}(1) En el Resumen y en la Introducción, nombrar origen y autoría: «una tabla de especificaciones, construida por el autor, que contrasta los instrumentos del software con las destrezas de interpretación y verificación de resultados que un consenso publicado de 67 expertos mexicanos de la industria y la academia considera necesarias para el egresado»; y en la Introducción (folio 2, PDF 17) cambiar «confirma que» por «reporta que» —en 1.2 el verbo ya es «encuentra que» y no se toca—. (2) En 2.1.5 (folio 45, PDF 60), sustituir la aserción de universalidad por un argumento sostenido: que los ítems adoptados nombran operaciones del propio método (mallado, post-proceso, conceptos fundamentales, verificación y validación), respaldándolo con localizador en un texto clásico del MEF ya presente en la bibliografía; y reconocer que lo intransferible es la ponderación entre ítems, que esta tabla no realiza porque solo registra cobertura. Si además se comprueba —ítem por ítem, sobre [3, p. 1163]\allowbreak{}— que ningún ítem dependiente del contexto entró en el universo, decirlo; si no se comprueba, no afirmarlo. (3) Añadir a las Limitaciones (folio 91, PDF 106), junto a la frase que ya reconoce la autoría de las tablas, una oración sobre la procedencia del criterio externo y su alcance declarado.\par
\rotulo{También en}\begin{itemize}[leftmargin=1.2em,itemsep=1pt,topsep=1pt]
\item {\small §2.1.5 Procedimiento e instrumentos, f. 45: \cita{los ítems que se toman son contenido universal del método}}
\item {\small Introducción — Planteamiento del problema, f. 2: \cita{Un estudio con expertos de la industria y de la academia confirma que ese criterio}}
\item {\small Limitaciones, f. 91: \cita{validez de contenido no es efectividad}}
\end{itemize}
{\footnotesize\color{gris}\rotulo{Verificación}ajustado: Verificado y se sostiene en alta. Las dos citas son literales y están donde se dice (Resumen, PDF 3 /\allowbreak{} folio ii; salvedad de 2.1.5, PDF 60 /\allowbreak{} folio 45), igual que las otras dos ubicaciones (Introducción, PDF 17 /\allowbreak{} folio 2; Limitaciones, PDF 106 /\allowbreak{} folio 91). El texto de 2.1.5 confirma el núcleo del hallazgo: reconoce la limitación —«expertos…\par}
\end{ficha}
\subsubsection*{Medias y bajas}
{\small\begin{longtable}{@{}L{1.6cm}L{1.6cm}L{5.25cm}L{5.65cm}r@{}}
\toprule
\sffamily\bfseries ID & \sffamily\bfseries Dónde & \sffamily\bfseries Qué pasa & \sffamily\bfseries Qué hacer & \sffamily\bfseries min\\
\midrule\endhead
\bottomrule\endfoot
\mbox{\casilla CLA-04}\newline {\color{media}\sffamily\bfseries\scriptsize MEDIA} & f. i\newline{\scriptsize Claridad} & Del mapeo al ensamblaje hay siete módulos (M1 a M7); el octavo (M0) trata la calidad de malla y no expone ninguna etapa del canal. Además «cada etapa» sobrepasa lo verificado, porque dos de las nueve etapas del canal no tienen módulo y solo se desarrollan en la memoria de cálculo, como reconocen 3.6.1 y las limitaciones. & «ocho módulos educativos —uno de calidad de malla y siete que exponen paso a paso las etapas del cálculo, del mapeo isoparamétrico al ensamblaje global—». & 5\\\addlinespace[3pt]
\mbox{\casilla ACC-05}\newline {\color{media}\sffamily\bfseries\scriptsize MEDIA} & f. i\newline{\scriptsize Accesibilidad} & Tres términos metodológicos que el tribunal necesita para juzgar el trabajo —validez de contenido, tabla de especificaciones y variable interviniente— se usan en el Resumen y en la Introducción y solo se definen en la sección 2.1, entre treinta y tres y treinta y ocho folios más adelante. No inducen a error, pero obligan a leer el Resumen dos veces. & Glosar cada uno en su primera aparición, con una aposición breve, y no repetir la definición después. En el Resumen: «con una tabla de especificaciones —el instrumento que confronta, ítem por ítem, lo que el software ofrece con lo que un criterio externo exige—» y «la validez de contenido, es decir, el grado en que los… & 30\\\addlinespace[3pt]
\mbox{\casilla APF-03}\newline {\color{media}\sffamily\bfseries\scriptsize MEDIA} & f. i-ii\newline{\scriptsize Aparato formal} & El resumen incumple la norma de presentación que el propio documento declara, en extensión (casi el doble del máximo) y en número de palabras clave. No es un defecto de contenido, pero es de los primeros que un tribunal contrasta contra la norma citada. & Recortar a unas 250 palabras conservando brecha, artefacto, hipótesis y su alcance, y comprimiendo la batería de verificación y validación a una frase por caso. Reducir a cinco palabras clave, fundiendo «elementos isoparamétricos Q4/\allowbreak{}Q9» en «método de elementos finitos» o suprimiendo «memoria de cálculo». Si el reglamento de la… & 30\\\addlinespace[3pt]
\mbox{\casilla RED-05}\newline {\color{media}\sffamily\bfseries\scriptsize MEDIA} & f. ii\newline{\scriptsize Redundancia} & La salvedad de que el efecto sobre el estudiante no se mide y queda para una prueba de campo se enuncia con oración completa una docena de veces. Es la garantía de calibración del documento y su ausencia sería un defecto mayor que su exceso, pero repetida en cada subsección se lee como disculpa y pierde fuerza donde de verdad importa. Se conservan íntegras las apariciones funcionales (Resumen, hipótesis, alcance, criterios de aceptación, contraste y cierre de Conclusiones). & Aparición canónica del enunciado completo: cinco nodos —Resumen, Introducción/\allowbreak{}Hipótesis, §2.1.6, §3.8 y cierre de Conclusiones—, que es donde el lector toma la afirmación al pie de la letra. En §1.2, §2.1.1, §2.1.4, §2.1.5, §3.7.2, Introducción/\allowbreak{}Metodología y Conclusiones/\allowbreak{}Limitaciones reducirla a media línea con remisión: «el… & 40\\\addlinespace[3pt]
\end{longtable}}
\subsection{Capítulo 1 · 1.1–1.2 Estado del arte y fundamentos pedagógicos}
{\small\color{gris}\sffamily 2 correcciones críticas o altas y 10 medias o bajas en esta sección · 5,6 h · más 1 frases de calibración (sección 3)}\par
{\small\rotulo{Primero, la calibración de la sección 3}CAL-07 (f. 17).\par}
\begin{ficha}{alta}{\casilla\textbf{TRZ-03}\quad\chip{alta}{alta}\quad §1.2 Fundamentos pedagógicos de la enseñanza del MEF · folio 18 · PDF 33\hfill 15 min}
\rotulo{Dice}\cita{Ese supuesto descansa en cuatro principios que la literatura sobre enseñanza del MEF documenta, apoyados cada uno en la investigación general sobre el aprendizaje con simulaciones y ejemplos}\par
\rotulo{Problema}La agrupación en cuatro principios es una síntesis propia construida a partir de fuentes heterogéneas, pero se presenta con la autoridad de un resultado publicado; como la matriz de principios es el instrumento con que se contrasta la cláusula (c) de la hipótesis, esa cláusula queda apoyada en un criterio externo que, tal como se enuncia, no existe.\par
\rotulo{Corrección · reformular}Sustituir la atribución conservando el resto de la oración: «Ese supuesto descansa en cuatro principios que este trabajo formula a partir de la literatura sobre enseñanza del MEF y de la investigación general sobre el aprendizaje con simulaciones y con ejemplos resueltos; la agrupación en cuatro principios es propia —ninguna de las fuentes citadas la enuncia como tal— y cada principio se apoya en las fuentes que se indican en su párrafo». Repetir el carácter propio del instrumento al presentar la Tabla 3.8 («matriz construida por el autor a partir de los principios de la Sección 1.2») y añadir una frase en Limitaciones que diga que los dos instrumentos de validez de contenido los construye el autor: la tabla de especificaciones sobre ítems del consenso publicado, y la matriz de principios sobre un marco propio.\par
\rotulo{También en}\begin{itemize}[leftmargin=1.2em,itemsep=1pt,topsep=1pt]
\item {\small §3.6.2, matriz de principios (Tabla 3.8), f. 77: \cita{los cuatro principios de aprendizaje tienen encarnación comprobable}}
\end{itemize}
{\footnotesize\color{gris}\rotulo{Verificación}ajustado: La cita es literal y está en su página (PDF 33, folio 18). El defecto existe: «cuatro principios que la literatura sobre enseñanza del MEF documenta» atribuye a fuentes externas una agrupación que es del autor, y la matriz de principios (Tabla 3.8) es uno de los dos instrumentos con que se declara confirmada la cláusula (c) en 3.8 (PDF…\par}
\end{ficha}
\begin{ficha}{alta}{\casilla\textbf{ACC-02}\quad\chip{alta}{alta}\quad §1.3 · folio 20 · PDF 35\hfill 45 min}
\rotulo{Dice}\cita{«La forma fuerte exige continuidad de segundo orden del campo de desplazamientos u, requisito difícil de satisfacer de manera global.» (PDF 35 /\allowbreak{} folio 20, literal); la sección 1.3 abre con «Sea un dominio bidimensional Omega subconjunto de R2 con frontera Gamma = Gamma\_\allowbreak{}u unión Gamma\_\allowbreak{}t, donde Gamma\_\allowbreak{}u es la porción con desplazamientos prescritos (restricciones de Dirichlet) y Gamma\_\allowbreak{}t la porción con tracciones impuestas (Neumann).»}\par
\rotulo{Problema}El capítulo entra en el MEF por la formulación variacional sin tender el puente desde lo único que el lector destinatario ya domina —el análisis matricial de pórticos— hacia el medio continuo, de modo que «forma fuerte», «forma débil» y «desplazamientos virtuales admisibles» aparecen sin anclaje. Es el punto donde el tribunal de ingeniería civil pierde el hilo del argumento y ya no lo recupera en todo el capítulo 1.\par
\rotulo{Corrección · anadir}Anteponer a la sección 1.3 un párrafo puente, sin la frase final del borrador original (que duplicaría la explicación de la forma débil ya presente en el texto): «En el análisis matricial de pórticos, cada barra aporta una matriz de rigidez, esas matrices se ensamblan en una matriz global y el sistema K u = F devuelve los desplazamientos de los nudos. El MEF hace lo mismo con un sólido continuo, con una diferencia: una barra tiene solución exacta conocida, mientras que un trozo de chapa no la tiene, y por eso su rigidez se obtiene integrando sobre el área del elemento. Las secciones siguientes recorren ese cálculo pieza por pieza.»\par
{\footnotesize\color{gris}\rotulo{Verificación}ajustado: Confirmado en lo sustancial: la sección 1.3 abre directamente con «Sea un dominio bidimensional Omega subconjunto de R2 con frontera Gamma = Gamma\_\allowbreak{}u unión Gamma\_\allowbreak{}t...», y el tramo «requisito difícil de satisfacer de manera global» es literal en PDF 35 /\allowbreak{} folio 20. No hay ningún puente desde el análisis matricial de pórticos -lo único que…\par}
\end{ficha}
\subsubsection*{Medias y bajas}
{\small\begin{longtable}{@{}L{1.6cm}L{1.6cm}L{5.25cm}L{5.65cm}r@{}}
\toprule
\sffamily\bfseries ID & \sffamily\bfseries Dónde & \sffamily\bfseries Qué pasa & \sffamily\bfseries Qué hacer & \sffamily\bfseries min\\
\midrule\endhead
\bottomrule\endfoot
\mbox{\casilla APF-06}\newline {\color{media}\sffamily\bfseries\scriptsize MEDIA} & f. 12\newline{\scriptsize Aparato formal} & Tres títulos combinan un enunciado general con una enumeración tras dos puntos y superan las diez-catorce palabras, mientras el resto del documento tiene títulos de tres a seis. En un índice de una línea por entrada esos tres se comprimen o cuesta escanearlos, justo en el primer contacto del tribunal con la estructura del trabajo. & Acortar a la idea central y mover el detalle al primer párrafo de cada sección: «Marco teórico del Método de los Elementos Finitos» (Capítulo 1), «Resultados del software» (3.6, dejando cobertura, validez y coherencia para las dos subsecciones que ya las nombran) y «Bloqueo y modos espurios» (1.9). & 15\\\addlinespace[3pt]
\mbox{\casilla TRZ-13}\newline {\color{media}\sffamily\bfseries\scriptsize MEDIA} & f. 14\newline{\scriptsize Trazabilidad} & El texto invoca «esta revisión» como respaldo del vacío que justifica todo el trabajo y como cierre del primer objetivo específico, pero en ninguna parte declara su alcance —qué fuentes se consultaron, con qué términos, en qué período, cuántos recursos se examinaron y con qué criterios de inclusión—, de modo que la afirmación de vacío no es auditable ni reproducible por el tribunal. & Insertar antes de la presentación de los antecedentes un párrafo breve con el contenido real de la búsqueda —repositorios consultados, términos empleados, período, número de publicaciones y de programas examinados y criterios de inclusión y exclusión—, aclarando que las dos categorías de propósito general (bibliotecas de… & 60\\\addlinespace[3pt]
\mbox{\casilla CLA-09}\newline {\color{media}\sffamily\bfseries\scriptsize MEDIA} & f. 14\newline{\scriptsize Claridad} & Oración de 78 palabras con cuatro incisos anidados que además dice dos veces lo mismo: el inciso inicial ya había establecido la interfaz en castellano que «sino también, en buena medida, el idioma» repite. Es el antecedente más directo del trabajo, de modo que la comparación debería leerse de una vez. & «ED-Elas2D comparte con EduFEM el dominio —elasticidad plana en tensión y deformación plana—, la organización en pre-proceso, proceso y post-proceso y una biblioteca de elementos que incluye los cuadriláteros Q4, Q8 y Q9. También, al menos en parte, el idioma: se desarrolló en una institución española y las capturas de pantalla… & 15\\\addlinespace[3pt]
\mbox{\casilla RED-08}\newline {\color{media}\sffamily\bfseries\scriptsize MEDIA} & f. 15\newline{\scriptsize Redundancia} & La misma lista de cuatro prestaciones se despliega en el marco teórico y otra vez en los aportes, a setenta y cinco folios de distancia, sin que la segunda añada matiz alguno. & Aparición canónica: §1.1, folio 15, donde la enumeración responde al hueco detectado en el antecedente. En «Aportes» suprimir la lista y dejar la formulación breve: «Frente al antecedente más directo, cuyo post-proceso devuelve la respuesta sin explicarla, EduFEM la expone con su procedimiento y sus medios de verificación… & 15\\\addlinespace[3pt]
\mbox{\casilla APF-07}\newline {\color{media}\sffamily\bfseries\scriptsize MEDIA} & f. 15\newline{\scriptsize Aparato formal} & El documento remite a sí mismo sin número de sección, tabla ni página justo en los puntos donde el tribunal querría ir a comprobar: el objetivo específico que una sección afirma cumplir, la delimitación del alcance, el dato de desempeño que respalda una limitación y las dos líneas de extensión futura que el Capítulo 1 promete encontrar en las recomendaciones. En un documento de 168 páginas, remitir al capítulo completo obliga a recorrer veinticinco páginas para verificar una cifra. & Sustituir cada remisión por una referencia numerada con \textbackslash{}autoref a la sección, tabla o figura concreta, y reservar el \textbackslash{}autoref de capítulo para cuando la referencia sea realmente al conjunto. En el caso del objetivo específico, reproducir además su enunciado: «Este diagnóstico cumple el primer objetivo específico —caracterizar… & 60\\\addlinespace[3pt]
\mbox{\casilla CON-05}\newline {\color{media}\sffamily\bfseries\scriptsize MEDIA} & f. 16\newline{\scriptsize Consistencia} & La celda afirma un atributo que el propio texto declara no documentado por la fuente -«a juzgar por las capturas de pantalla de su publicación, que no declara el idioma»-, con lo que incumple la regla que la nota de la tabla enuncia para los datos no documentados y presenta como dato una inferencia del autor. & Cambiar la celda por «No consta (las capturas de la publicación muestran una interfaz en castellano)», que es lo que la fuente permite sostener; o mantener «Español (parcial)» y añadir en la nota de la tabla la salvedad de que el dato se infiere de las capturas de pantalla. & 10\\\addlinespace[3pt]
\mbox{\casilla APF-08}\newline {\color{media}\sffamily\bfseries\scriptsize MEDIA} & f. 17\newline{\scriptsize Aparato formal} & La sigla V\&V se estrena en una tabla del folio 17 y reaparece en tres tablas del Capítulo 2, pero solo se desarrolla en el folio 61. Un tribunal no especialista atraviesa cuarenta y cuatro folios con una sigla sin desplegar en la tabla que sostiene el diagnóstico de la brecha. & Desarrollarla en su primera aparición y suprimir el desarrollo tardío: en la nota al pie de la Tabla 1.1 escribir «V\&V: verificación y validación (Sección 1.13)» y dejar en el folio 61 «una batería sistemática de V\&V, siguiendo la distinción metodológica clásica…». Opción más limpia para el tribunal: escribir «verificación y… & 15\\\addlinespace[3pt]
\mbox{\casilla RED-09}\newline {\color{media}\sffamily\bfseries\scriptsize MEDIA} & f. 19\newline{\scriptsize Redundancia} & El mismo hallazgo de Pérez-Santiago y Campos se enuncia entero cuatro veces, dos de ellas (§1.1 y §1.2) casi palabra por palabra y con la misma referencia a la misma página. Solo la de §1.2 cumple función argumental: es la que declara que las destrezas del consenso son la definición operativa del criterio, de modo que las otras tres gastan el hallazgo antes de que sirva para algo. & Aparición canónica: §1.2, cuarto principio (folio 19, PDF 34), con [3, pp. 1159, 1168]\allowbreak{} y la declaración de definición operativa del criterio. En §1.1 (folio 13) conservar una sola oración con su cita -«un estudio con expertos de la industria y de la academia señala que estos cursos privilegian las destrezas prácticas sobre los… & 30\\\addlinespace[3pt]
\mbox{\casilla ACC-06}\newline {\color{media}\sffamily\bfseries\scriptsize MEDIA} & f. 19\newline{\scriptsize Accesibilidad} & ICAP es una sigla que el documento nunca desarrolla ni recoge en la tabla de siglas de la Nomenclatura, y que reaparece cincuenta y ocho folios más adelante, en el contraste de la cláusula de coherencia, cuando el lector ya no tiene a mano la explicación. Solo aparece tres veces en las 168 páginas (PDF 34, 92 y, como título de la referencia de Chi y Wylie, PDF 113). & Desarrollarla en su primera aparición: «el marco ICAP —por las iniciales en inglés de interactivo, constructivo, activo y pasivo— de la psicología del aprendizaje explica...», ordenando los cuatro modos como los ordena la sigla. Añadir a la tabla de siglas la entrada «ICAP — Marco de modos de compromiso cognitivo: interactivo,… & 15\\\addlinespace[3pt]
\mbox{\casilla CON-06}\newline {\color{media}\sffamily\bfseries\scriptsize MEDIA} & f. 19\newline{\scriptsize Consistencia} & El destinatario se llama «estudiante» en el problema, la hipótesis y las variables (62 apariciones), «alumno» en los fundamentos pedagógicos y en parte del Capítulo 2 (22 apariciones, seis de ellas en el folio 19) y «usuario» en el Capítulo 2 y los anexos (26 apariciones), con lo que el sujeto sobre el que se define la variable dependiente pierde nitidez justo donde el tribunal comprueba la población. & Nombre canónico: «estudiante», porque es el término con que están escritos el problema científico, la hipótesis y la población. Dejar «alumno» solo dentro de una cita o paráfrasis de una fuente que lo emplee, y «usuario» solo en los anexos de manual y cuando se hable de la interacción con la interfaz con independencia de quien… & 40\\\addlinespace[3pt]
\end{longtable}}
\subsection{Capítulo 1 · 1.3–1.13 Teoría del MEF y V\&V}
{\small\color{gris}\sffamily 2 correcciones críticas o altas y 9 medias o bajas en esta sección · 7,8 h · más 0 frases de calibración (sección 3)}\par
\begin{ficha}{alta}{\casilla\textbf{ACC-03}\quad\chip{alta}{alta}\quad §1.13 · folio 35 · PDF 50\hfill 60 min}
\rotulo{Dice}\cita{«Al refinar la malla, estas normas deben decrecer con las tasas asintóticas teóricas [...]\allowbreak{} lo que confirma la correcta implementación» (el tramo «lo que confirma la correcta implementación» es literal en PDF 50 /\allowbreak{} folio 35); en 3.1-3.2 la misma terminología reaparece sin glosa: «La exactitud frente a soluciones de referencia se cuantifica con la norma...» (literal en PDF 77 /\allowbreak{} folio 62; la unidad la había declarado en PDF 76 /\allowbreak{} folio 61).}\par
\rotulo{Problema}La norma L2, la seminorma H1, la tasa de convergencia asintótica, el orden polinómico p y el método de soluciones manufacturadas son los instrumentos con que se mide la evidencia principal del trabajo —la cláusula de corrección del motor— y nunca se explican en castellano llano, ni donde se definen (1.13) ni donde se presentan los resultados (3.1-3.2). El tribunal no especialista no puede juzgar si la verificación es exigente, qué significa que una tasa sea de orden dos o tres, ni por qué el Q9 alcanza con 578 GDL una exactitud que el Q4 no logra con 2178.\par
\rotulo{Corrección · anadir}En 1.13, intercalar antes de las ecuaciones de error (sin glosar el MMS, que el propio párrafo ya explica): «Para comparar la solución del programa con la exacta no basta mirar un punto: hace falta una medida única del error en toda la pieza. La norma L2 del desplazamiento es esa medida -una media cuadrática de la diferencia, ponderada por el área- y la seminorma H1 hace lo mismo con los gradientes del desplazamiento, es decir, con las deformaciones, que son lo que determina las tensiones. Al reducir a la mitad el tamaño h de los elementos, ambas medidas deben caer en una proporción fija que la teoría predice: esa proporción es la tasa de convergencia, y comprobarla es la prueba más exigente de que el programa resuelve lo que dice resolver.» En 3.1-3.2, tras la primera mención: «En términos llanos, la norma L2 mide el desajuste promedio entre el desplazamiento calculado y el exacto en todo el dominio, y la seminorma H1 el de sus derivadas, es decir, el de las deformaciones y tensiones que el ingeniero acaba leyendo; por eso converge más despacio y es la exigente.» Y al presentar las tasas de la Tabla 3.1: «El orden fija cuánto cuesta ganar exactitud: con orden 2, partir cada elemento en dos divide el error por cuatro; con orden 3, por ocho.» No repetir la explicación de la cuadratura con un punto de Gauss adicional: ya está en el folio 35.\par
\rotulo{También en}\begin{itemize}[leftmargin=1.2em,itemsep=1pt,topsep=1pt]
\item {\small §3.1-3.2, f. 62: \cita{La exactitud frente a soluciones de referencia se cuantifica con la norma L2 del desplazamiento y la seminorma H1 de su gradiente}}
\item {\small §3.2, f. 63: \cita{El elemento Q4 alcanza asintóticamente tasas de O(h\textasciicircum{}2,00) en L2 y O(h\textasciicircum{}1,00) en H1}}
\end{itemize}
{\footnotesize\color{gris}\rotulo{Verificación}ajustado: Las tres citas son literales y están en sus páginas: «lo que confirma la correcta implementación» en PDF 50 /\allowbreak{} folio 35 (1.13), «La exactitud frente a soluciones de referencia se cuantifica con la norma L2 del desplazamiento y la seminorma H1 de su gradiente» en PDF 77 /\allowbreak{} folio 62 (3.1) y «El elemento Q4 alcanza asintóticamente tasas de…\par}
\end{ficha}
\begin{ficha}{alta}{\casilla\textbf{TRZ-04}\quad\chip{alta}{alta}\quad §1.13 Verificación y validación (membrana de Cook) · folio 36 · PDF 51\quad{\color{gris}\footnotesize también RED-12}\hfill 30 min}
\rotulo{Dice}\cita{El valor de referencia de uso convencional para ese desplazamiento, 23,96, no procede de las obras citadas —ninguna de ellas trata el caso— sino de la práctica de la comunidad}\par
\rotulo{Problema}La cifra que el capítulo 3 usa como patrón de comparación del criterio de Cook queda atribuida a «la práctica de la comunidad», sin fuente verificable, y la salvedad «ninguna de ellas trata el caso» contradice la oración inmediatamente anterior, que cita esas mismas obras ([2, pp. 98-99]\allowbreak{} y [16, p. 243]\allowbreak{}) precisamente para describir el comportamiento del Q4 y del Q9 en ese caso.\par
\rotulo{Corrección · corregir}Separar las dos cosas: «Las obras citadas describen el caso, pero ninguna publica un valor de referencia para ese desplazamiento; la cifra de uso convencional, 23,96, procede de la literatura de verificación de elementos finitos [cita con localizador]\allowbreak{}. Este trabajo no la toma como dato externo: la corrobora con su propia extrapolación de Richardson (sección 3.5)». Si no se localiza una fuente publicada que la sostenga, eliminar la atribución a «la práctica de la comunidad» y presentar 23,96 únicamente como resultado propio coincidente con el valor de uso extendido.\par
\rotulo{También en}\begin{itemize}[leftmargin=1.2em,itemsep=1pt,topsep=1pt]
\item {\small §1.13 (oración inmediatamente anterior), f. 36: \cita{es el banco de prueba clásico de la literatura para evaluar la sensibilidad de los elementos a la distorsión geométrica y al bloqueo por cortante [...]\allowbreak{} [2, pp. 98-99]\allowbreak{} [...]\allowbreak{} [16, p. 243]\allowbreak{}}}
\item {\small §3.5 (criterio de Cook), f. 71: \cita{El elemento Q4 subestima la respuesta en más de un}}
\end{itemize}
{\small\color{gris}\rotulo{RED-12}La membrana de Cook se presenta dos veces desde cero —autoría, condición de banco de prueba clásico y mecanismo del bloqueo— y el tratamiento del valor de referencia 23,96 con su corroboración por extrapolación de Richardson se anticipa en el marco teórico y se repite entero en §3.5. El marco teórico termina adelantando un resultado propio del capítulo de resultados.\par}
\end{ficha}
{\small\rotulo{Párrafos que pide la defensa}\par}
\begin{ficha}{media}{\casilla\textbf{DEF-10}\quad\chip{media}{media}\quad §1.10 Imposición de restricciones y solución del sistema · folio 29 · PDF 44\hfill 15 min}
\rotulo{Dice}\cita{Un modelo bien restringido —al menos tres GDL que impidan los movimientos de cuerpo rígido— hace que Kff sea definida positiva y, por tanto, no singular}\par
\rotulo{Problema}La oración es correcta porque la cláusula relativa hace el trabajo, pero se lee con facilidad como si bastaran tres restricciones cualesquiera; el caso de tres apoyos paralelos o concurrentes, que un tribunal de ingeniería civil reconoce al instante, queda sin mencionar justo donde el texto enuncia la condición.\par
\rotulo{Corrección · corregir}Añadir el inciso de disposición: «…al menos tres GDL, dispuestos de modo que ninguna combinación de ellos permita las dos traslaciones ni la rotación del sólido como cuerpo rígido; tres apoyos paralelos o concurrentes no bastan…».\par
\end{ficha}
\begin{ficha}{media}{\casilla\textbf{DEF-11}\quad\chip{media}{media}\quad §1.11 Recuperación y suavizado de tensiones · folio 31 · PDF 46\hfill 25 min}
\rotulo{Dice}\cita{El orden de convergencia que resulta de esta secuencia de muestreo, extrapolación y promediado no se toma de la literatura: se mide en el capítulo de resultados con el método de las soluciones manufacturadas}\par
\rotulo{Problema}El texto declara con honestidad que la elección de muestreo no es la óptima y remite la consecuencia al Capítulo 3, pero no dice qué se pierde, de modo que la pregunta obvia —por qué no se muestrea en los puntos superconvergentes— queda sin respuesta en el lugar donde el lector la formula.\par
\rotulo{Corrección · anadir}Cerrar el razonamiento antes de la remisión: el coste esperable es que las tensiones nodales no hereden la precisión superior de los puntos óptimos y que la extrapolación amplifique las diferencias entre puntos de Gauss —los pesos de la matriz del Q4 llegan a 1,87 y -0,5—; a cambio, una única ruta de cálculo sirve a los dos tipos de elemento y conserva la variación del campo dentro de cada uno.\par
\end{ficha}
\subsubsection*{Medias y bajas}
{\small\begin{longtable}{@{}L{1.6cm}L{1.6cm}L{5.25cm}L{5.65cm}r@{}}
\toprule
\sffamily\bfseries ID & \sffamily\bfseries Dónde & \sffamily\bfseries Qué pasa & \sffamily\bfseries Qué hacer & \sffamily\bfseries min\\
\midrule\endhead
\bottomrule\endfoot
\mbox{\casilla TRZ-14}\newline {\color{media}\sffamily\bfseries\scriptsize MEDIA} & f. 22\newline{\scriptsize Trazabilidad} & Dieciocho de las veintinueve ecuaciones numeradas del Capítulo 1 —D en tensión y deformación plana, funciones de forma de Q4 y Q9, Jacobiano y su determinante e inversa, cinemática, matriz B, integral de rigidez, ensamblaje, partición, extrapolación inversa, promediado, Jacobiano escalado, estiramiento, descomposición bilineal y normas de error— no vuelven a citarse en ningún punto posterior, de modo que su número no cumple función de referencia y el hilo entre la derivación teórica y su cómputo real no queda cerrado en ninguna dirección. & Adoptar una política única: o se añaden remisiones explícitas a esas ecuaciones desde la subsección 2.2.5, desde el Anexo G —que instancia numéricamente D, J, B y la matriz de rigidez elemental sobre el ejemplo canónico— y desde los módulos M1 a M5 que exponen cada matriz; o se retira el número a las que son puramente… & 45\\\addlinespace[3pt]
\mbox{\casilla ACC-07}\newline {\color{media}\sffamily\bfseries\scriptsize MEDIA} & f. 23\newline{\scriptsize Accesibilidad} & Patrón transversal: el término técnico se usa en su primera aparición sin una aposición de media línea que diga qué es y por qué importa, y su definición —cuando existe— llega decenas de folios después o solo en la nomenclatura. Los lectores de unidad inventariaron una treintena de casos repartidos en siete secciones. El efecto acumulado es que el tribunal no especialista no puede juzgar si la propiedad invocada es relevante ni si el criterio aplicado es exigente; en 1.1, además, impide ver en qué consiste el vacío que la tesis dice cubrir. & Añadir a cada término, en su primera aparición, media línea de glosa entre rayas o paréntesis y la remisión numerada al apartado donde se desarrolla, sin repetir la definición después ni recurrir a notas al pie. Adoptar tal cual: «propiedad de Kronecker (cada función vale uno en su propio nodo y cero en los demás, de modo que… & 120\\\addlinespace[3pt]
\mbox{\casilla APF-09}\newline {\color{media}\sffamily\bfseries\scriptsize MEDIA} & f. 23\newline{\scriptsize Aparato formal} & Diez leyendas de figura y tabla son meros rótulos que no se pueden leer sin el cuerpo: no dicen qué representan las cajas ni en qué sentido corren las flechas, de qué módulo es la captura ni sobre qué modelo, qué magnitud va en cada eje, en qué unidades está —ni que los casos del MMS y de Cook son adimensionales mientras Timoshenko va en kgf/\allowbreak{}cm2 y cm—, ni de qué modelo son los nodos que tabula el Anexo G. La leyenda es lo único que acompaña a la figura cuando el tribunal la hojea suelta, y aquí falla en el Capítulo 2 (cuatro de cinco), en las tablas y figuras de resultados del Capítulo 3 (cuatro) y en el anexo que se presenta como memoria autocontenida. & Reescribir cada leyenda como enunciado autosuficiente, con la fórmula «qué se muestra + sobre qué modelo + cómo leerlo». Modelos: «Arquitectura por capas de EduFEM. Cada flecha indica dependencia: una capa solo invoca a las que están por debajo, y todas comparten por referencia el mismo modelo de proyecto; el motor numérico no… & 75\\\addlinespace[3pt]
\mbox{\casilla REF-03}\newline {\color{media}\sffamily\bfseries\scriptsize MEDIA} & f. 23\newline{\scriptsize Referencias} & Dos atribuciones concretas se dan sin ninguna cita, no ya sin localizador: la comparación técnica entre el Q9 y el Q4 en §1.5, que es una afirmación de la literatura y no un resultado propio de este trabajo, y la paternidad del caso de prueba en §3.5, que atribuye la membrana a R. D. Cook sin remitir a la publicación en que la propuso. En un documento que declara dar localizador en toda cita que sostenga una atribución concreta, la ausencia total de cita es más visible que la falta de página. & En §1.5, cerrar la oración con la página en que Hughes o Bathe comparan el orden de interpolación y la tolerancia a la distorsión ([16, pp. 192, 208]\allowbreak{} y [4, pp. 476-477]\allowbreak{} ya se citan en el Capítulo 1 para ese mismo contenido). En §3.5, o bien citar la fuente original del caso con su página, o bien reformular sin atribuir autoría:… & 20\\\addlinespace[3pt]
\mbox{\casilla RED-10}\newline {\color{media}\sffamily\bfseries\scriptsize MEDIA} & f. 26\newline{\scriptsize Redundancia} & El mecanismo de los modos espurios se explica dos veces casi con las mismas palabras y con la misma cita de página, y el papel de la membrana de Cook frente al bloqueo por cortante se repite entre §1.9 y §1.13 con remisiones cruzadas circulares, de modo que el lector no sabe cuál de las dos secciones es la fuente. & Aparición canónica del mecanismo: §1.9 (folio 27), donde el fenómeno se trata como tal. §1.8 se reduce a su función original, justificar el orden de integración: «Una subintegración puede dejar la matriz de rigidez con rango deficiente y admitir modos de deformación de energía nula, los modos espurios (hourglass) que se… & 25\\\addlinespace[3pt]
\mbox{\casilla CLA-12}\newline {\color{media}\sffamily\bfseries\scriptsize MEDIA} & f. 33\newline{\scriptsize Claridad} & «Puede combinar» oculta si el programa lo hace o no —el párrafo siguiente revela que sí—, y «ortogonales» se usa en sentido figurado a pocas líneas de «mapeo ortogonal» en sentido geométrico, para nombrar la misma propiedad que el cierre del apartado llama «complementarias». El lector no especialista no puede saber si son dos propiedades o una. & «El programa asigna además a cada elemento un estado global de calidad —buena, aceptable o mala— combinando cuatro métricas.» Sustituir «dos métricas ortogonales» por «dos métricas complementarias, que miden propiedades independientes: una los ángulos y otra las proporciones», reservar «ortogonal» para su sentido geométrico y… & 15\\\addlinespace[3pt]
\mbox{\casilla CON-07}\newline {\color{media}\sffamily\bfseries\scriptsize MEDIA} & f. 34\newline{\scriptsize Consistencia} & La Sección 1.13 declara una convención para «en el resto del documento» y el documento la incumple después en el término más delicado del trabajo: la misma raíz designa además la validez de contenido (folio 7), el validador de salud del modelo (folio 56), el cuadrado de validación del MMS (folio 108, que es un caso de verificación, no de validación) y, en las Conclusiones, la validación experimental que se declara fuera de alcance. El lector no especialista no puede saber, al leer «validado», cuál de los cuatro sentidos está en juego. & Ampliar en una oración la declaración de 1.13, sin tocar los nombres que el software ya usa: «En lo que sigue, «validación» designa ese contraste con referencias externas de exactitud conocida. El documento emplea además «validez de contenido» (o «validación de contenido») para el contraste documental del Capítulo 3, que no es… & 40\\\addlinespace[3pt]
\mbox{\casilla RED-11}\newline {\color{media}\sffamily\bfseries\scriptsize MEDIA} & f. 35\newline{\scriptsize Redundancia} & La precaución de integrar las normas de error con un punto de Gauss adicional por dirección se enuncia cinco veces con su paréntesis y su justificación completa. Es un detalle de implementación que basta establecer una vez; repetido, hace pensar al lector que se trata de decisiones distintas. & Aparición canónica: §1.13 (folio 35, PDF 50), donde se establece el método de verificación, con la justificación completa. En §2.1.5, §3.1, §3.2 y el anexo D dejar únicamente «(integradas con la cuadratura de orden superior de la sección 1.13)», suprimiendo paréntesis y justificación. La aparición de §3.1 se resuelve junto con… & 20\\\addlinespace[3pt]
\mbox{\casilla APF-10}\newline {\color{media}\sffamily\bfseries\scriptsize MEDIA} & f. 35\newline{\scriptsize Aparato formal} & En la ecuación de las normas de error no se definen h (tamaño característico de malla), el subíndice h de la solución numérica ni el asterisco del campo de tensiones recuperado, y lo que el texto llama «norma H1» se escribe con las barras simples de la seminorma. La incoherencia entre el nombre y la notación es un flanco fácil para un tribunal técnico y se repara con dos líneas. & Añadir tras la ecuación: «donde u es la solución exacta, u\_\allowbreak{}h la calculada sobre una malla de tamaño característico h, y el asterisco distingue el campo de tensiones recuperado —extrapolado y promediado— del evaluado directamente en los puntos de Gauss». Y llamar a la segunda medida por su nombre: «la seminorma H1, que mide… & 20\\\addlinespace[3pt]
\end{longtable}}
\subsection{Capítulo 2 · 2.2 El software EduFEM}
{\small\color{gris}\sffamily 0 correcciones críticas o altas y 10 medias o bajas en esta sección · 5,7 h · más 1 frases de calibración (sección 3)}\par
{\small\rotulo{Primero, la calibración de la sección 3}CAL-09 (f. 56).\par}
{\small\rotulo{Párrafos que pide la defensa}\par}
\begin{ficha}{media}{\casilla\textbf{DEF-16}\quad\chip{media}{media}\quad §2.2.4 Modelo de datos e indexación · folio 52 · PDF 67\hfill 45 min}
\rotulo{Dice}\cita{«Una decisión de diseño crítica es la separación explícita entre la identidad de un nodo y su índice en el sistema algebraico» (verificada literal en PDF 67 /\allowbreak{} folio 52).}\par
\rotulo{Problema}Cuatro subsecciones del capítulo de implementación (2.2.2, 2.2.3, 2.2.4 y 2.2.9) se agotan en el cómo técnico y terminan sin la frase que ate la decisión al argumento de la tesis, de modo que el tribunal no especialista no puede distinguir lo que es decisión al servicio del objetivo de lo que es artesanía de programación. Las subsecciones 2.2.1, 2.2.6 y 2.2.7 sí tienen ese cierre, lo que hace más visible la ausencia.\par
\rotulo{Corrección · anadir}Cerrar cada una de esas cuatro subsecciones con una oración de propósito; para 2.2.4, que esa separación es lo que hace que los números de nodo que el estudiante ve en la tabla, en el lienzo, en el módulo M7 y en la memoria sigan designando al mismo nodo después de borrar o renumerar, y que las filas de la matriz que M7 resalta sean trazables hasta ese nodo. Si a alguna no se le encuentra ese cierre, trasladarla al anexo del manual.\par
\rotulo{También en}\begin{itemize}[leftmargin=1.2em,itemsep=1pt,topsep=1pt]
\item {\small §2.2.2, f. 49: \cita{La implementación se realizó íntegramente en Python, lenguaje que combina}}
\item {\small §2.2.5, f. 54: \cita{la factorización LU dispersa de SuperLU se invoca con el ordenamiento}}
\item {\small §2.2.9, f. 60: \cita{un archivo JSON producido por la serialización del modelo y escrito de forma atómica}}
\end{itemize}
\end{ficha}
\subsubsection*{Medias y bajas}
{\small\begin{longtable}{@{}L{1.6cm}L{1.6cm}L{5.25cm}L{5.65cm}r@{}}
\toprule
\sffamily\bfseries ID & \sffamily\bfseries Dónde & \sffamily\bfseries Qué pasa & \sffamily\bfseries Qué hacer & \sffamily\bfseries min\\
\midrule\endhead
\bottomrule\endfoot
\mbox{\casilla TRZ-19}\newline {\color{media}\sffamily\bfseries\scriptsize MEDIA} & f. 48\newline{\scriptsize Trazabilidad} & La subsección abre con un «objetivo central» propio, que no es el objetivo general del trabajo, y deriva de él los cuatro requisitos rectores sin citar ni el objetivo general ni los cuatro principios de la sección 1.2; en sustancia el contenido coincide —el primer requisito es que cada operación pueda inspeccionarse y no quede oculta tras una caja negra—, pero el eslabón objetivo general {\mates →} requisitos {\mates →} decisiones de diseño no está escrito justo donde el tribunal lo va a recorrer. & Basta con enlazar la primera oración con el objetivo general, sin tocar el resto del párrafo: «El objetivo general enunciado en la introducción -desarrollar un software educativo de elementos finitos para el análisis estructural- se traduce, en el plano del artefacto, en cuatro requisitos rectores que condicionaron todas las… & 30\\\addlinespace[3pt]
\mbox{\casilla ACC-08}\newline {\color{media}\sffamily\bfseries\scriptsize MEDIA} & f. 48\newline{\scriptsize Accesibilidad} & La sección describe nueve componentes del programa sin recordar antes el alcance del artefacto —análisis estático lineal, elasticidad plana, elementos Q4 y Q9, pequeños desplazamientos—, de modo que el tribunal empieza a leer la arquitectura del software sin tener a la vista qué problemas puede y no puede resolver. La delimitación existe en la Introducción, así que basta con recuperarla aquí en una frase y remitir. & Cerrar el párrafo introductorio de 2.2 con una frase de alcance y su remisión: «El programa resuelve el problema estático lineal de la elasticidad plana —tensión plana y deformación plana— con elementos isoparamétricos cuadriláteros de cuatro y nueve nodos y la hipótesis de pequeños desplazamientos; quedan fuera el análisis… & 20\\\addlinespace[3pt]
\mbox{\casilla CON-09}\newline {\color{media}\sffamily\bfseries\scriptsize MEDIA} & f. 48\newline{\scriptsize Consistencia} & EduFEM se nombra de seis maneras -herramienta (57), software (106), programa (51), aplicación (19), artefacto (15) y recurso (20)-, y las seis conviven dentro del mismo capítulo, lo que diluye el sujeto del que se predica cada afirmación; en particular «artefacto» es término técnico de la investigación tecnológica en 2.1 y se usa después como mero sinónimo. & Fijar tres usos y no más: «EduFEM» cuando se nombre el producto; «la herramienta» como sustituto en el cuerpo; «el artefacto» solo en la Sección 2.1, donde es el término de la investigación tecnológica, declarándolo allí una vez («el artefacto, es decir, EduFEM»). Sustituir «la aplicación» y «el programa» por «EduFEM» o «la… & 60\\\addlinespace[3pt]
\mbox{\casilla CON-20}\newline {\color{baja}\sffamily\bfseries\scriptsize BAJA} & f. 52\newline{\scriptsize Consistencia} & Un mismo objeto recibe dos nombres en el mismo párrafo -«mapa» y «diccionario»- y el mecanismo de deshacer se llama alternativamente «instantánea» y «estado capturado»: inconsistencias menores que obligan al lector no especialista a preguntarse si se habla de dos cosas distintas. & Nombres canónicos: «mapa de índices» en toda la subsección 2.2.4, suprimiendo «diccionario», e «instantánea» en 2.2.6, con la glosa «copia completa del estado del proyecto en un instante dado» en su primera aparición. & 10\\\addlinespace[3pt]
\mbox{\casilla RED-13}\newline {\color{media}\sffamily\bfseries\scriptsize MEDIA} & f. 53\newline{\scriptsize Redundancia} & La separación entre identidad del nodo e índice de grado de libertad, que es una decisión de diseño y no un resultado, se explica cuatro veces: donde se decide, donde se prueba y dos veces al concluir. & Aparición canónica: §2.2.4 (folio 53), donde la decisión de diseño se toma y se fundamenta. En §3.3 sustituir la explicación por «La segunda prueba ejercita la separación entre identidad e índice de grado de libertad descrita en la sección 2.2.4» y pasar directamente al resultado (equivalencia inferior a 1e-9 con… & 20\\\addlinespace[3pt]
\mbox{\casilla RED-14}\newline {\color{media}\sffamily\bfseries\scriptsize MEDIA} & f. 54\newline{\scriptsize Redundancia} & El estado actual del solucionador —ensamblaje por tripletes COO, conversión a CSR, factorización LU dispersa y ordenamiento de mínimo grado— se describe cuatro veces fuera de la sección que lo decide, incluida la recomendación, que recapitula lo implementado antes de proponer el cambio. & Aparición canónica: §2.2.5 (folio 54), que describe la decisión de implementación. §3.7.4 se limita a cuantificar su efecto (1,7, 2,0 y 2,9 veces) sin volver a describir el esquema; Conclusiones/\allowbreak{}Limitaciones lo menciona en una línea con remisión a §2.2.5; y la recomendación «Solucionador para problemas de gran porte» arranca… & 25\\\addlinespace[3pt]
\mbox{\casilla ACC-09}\newline {\color{media}\sffamily\bfseries\scriptsize MEDIA} & f. 54\newline{\scriptsize Accesibilidad} & Patrón: una veintena de términos de programación y estructuras de datos —factorización LU dispersa, ordenamiento de mínimo grado, COLAMD, Cuthill-McKee inverso, fill-in, COO, CSR, índice espacial, firma topológica, idempotente, función pura, serialización, escritura atómica, interpolación de Gouraud, triángulos baricéntricos, rasterizado, TSV, JSON— se usan sin una sola frase que diga qué se gana con ellos, ante un tribunal de ingeniería civil que no cursó estructuras de datos. En 3.7.4 el resultado de rendimiento (una reducción de 1,7 veces) queda sin significado legible, y el Anexo C, que prometía condensar la formulación, acaba enseñando álgebra dispersa. & Limitar la intervención a decir, una sola vez y en llano, qué es factorizar y qué es una matriz dispersa, y llevar los nombres de algoritmos, formatos e identificadores del código a nota al pie. En 2.2.5, anteponer a la frase actual: «El sistema se resuelve factorizando la matriz de rigidez: se la descompone en dos matrices… & 90\\\addlinespace[3pt]
\mbox{\casilla TRZ-20}\newline {\color{media}\sffamily\bfseries\scriptsize MEDIA} & f. 55\newline{\scriptsize Trazabilidad} & La cifra se enuncia sin decir sobre qué modelo se midió ni remitir a ningún punto del capítulo 3, de modo que el lector no puede saber si es una medición, una estimación de orden de magnitud o un supuesto de diseño. & Declarar la medición («entre 20 y 60 kB para los modelos de verificación del capítulo 3, de hasta N elementos») o marcarla como estimación («del orden de decenas de kilobytes, estimado a partir del tamaño del archivo de proyecto sin la solución»). Si no hay medición que la respalde, suprimir el dato y dejar la afirmación… & 15\\\addlinespace[3pt]
\mbox{\casilla ACC-10}\newline {\color{media}\sffamily\bfseries\scriptsize MEDIA} & f. 56\newline{\scriptsize Accesibilidad} & Una sola oración de unas ciento cincuenta palabras concentra tres juegos de umbrales, introduce cuatro símbolos —q\_\allowbreak{}SJ, R\_\allowbreak{}J, AR, T\_\allowbreak{}R— que no se glosan en ningún punto de la sección y remata con una lista posicional de cuatro cifras sueltas que obliga al lector a reconstruir de memoria a qué índice corresponde cada una, con riesgo de emparejar mal umbral e índice. & Partir la oración en dos frases —una que declare que los umbrales son convención del programa apoyada en los mínimos de la biblioteca Verdict, y otra que justifique la elevación del corte del jacobiano escalado— y trasladar los valores a una tabla de cuatro filas con las columnas: índice, símbolo, qué mide, umbral «bueno»,… & 45\\\addlinespace[3pt]
\mbox{\casilla TRZ-21}\newline {\color{media}\sffamily\bfseries\scriptsize MEDIA} & f. 57\newline{\scriptsize Trazabilidad} & La pieza que sostiene la cláusula (a) se presenta como un hecho consumado —ocho módulos, ni más ni menos— sin remitir a la tabla de especificaciones con que después se valida el contenido ni explicar de qué se deriva la selección; el contraste existe en 3.6.2, pero no hay señal que lo anuncie donde el lector lo necesita. & Anteponer una frase de derivación: «Los ocho módulos no se eligieron por conveniencia de implementación: cada uno corresponde a una operación del canal de cálculo que el universo de contenido de la subsección 2.1.4 identifica como exigible a partir del consenso de expertos recogido por Pérez-Santiago y Campos, y la subsección… & 25\\\addlinespace[3pt]
\end{longtable}}
\subsection{Capítulo 3 · 3.1–3.5 Verificación y validación numérica}
{\small\color{gris}\sffamily 1 correcciones críticas o altas y 12 medias o bajas en esta sección · 4,5 h · más 1 frases de calibración (sección 3)}\par
{\small\rotulo{Primero, la calibración de la sección 3}CAL-02 (f. 72).\par}
\begin{ficha}{alta}{\casilla\textbf{TRZ-07}\quad\chip{alta}{alta}\quad §3.4 (evidencia citada); el defecto se reparte entre 3.3, 3.4, 3.5 y 3.8 · folio 68 · PDF 83\hfill 60 min}
\rotulo{Dice}\cita{Los errores de EduFEM frente a la solución analítica permanecen por debajo del 0,05 \% en los tres puntos, con un máximo de 0,0414 \% en el punto B; la concordancia con SAP2000 es del mismo orden, con un máximo de 0,2066 \%}\par
\rotulo{Problema}De los diecisiete criterios de aceptación fijados en 2.1.6, cinco tienen cifra medida pero ningún veredicto —tasa mínima de la tensión recuperada, residuo de equilibrio menor que 1e-8, Cook con Q9 y N = 8 dentro del 1,5 \%, idempotencia de la reimportación DXF y generación de la memoria en Q4 y Q9— y un sexto (flecha del Q4 menor que la del Q9 con la misma malla) se juzga con otra comparación, a igualdad de GDL; además, las cuatro secciones de resultados reportan cifras sin enfrentarlas a su umbral y 3.2 remite a la cota empírica mínima sin dar su valor, de modo que la declaración global «la cláusula (b) queda confirmada» descansa en criterios que nunca se juzgan uno a uno. Todas las cifras cumplen su umbral: lo que falta es el veredicto escrito.\par
\rotulo{Corrección · anadir}Sustituir la prosa de las cláusulas (a) y (b) de 3.8 por una tabla «criterio /\allowbreak{} umbral /\allowbreak{} valor medido /\allowbreak{} dónde /\allowbreak{} veredicto» con todos los criterios de 2.1.6, y cerrar las secciones 3.3, 3.4 y 3.5 con una oración del molde «Criterio fijado en la Subsección 2.1.6: [umbral]\allowbreak{}. Observado: [cifra]\allowbreak{}. Se cumple». Oraciones que faltan: «el residuo de equilibrio, 1,7 × 10{\mates ⁻}¹³, es inferior al 10{\mates ⁻}{\mates ⁸} fijado en la Subsección 2.1.6»; «con N = 8 el Q9 queda a {\mates −}0,144 \% de la referencia, dentro del 1,5 \% fijado en la Subsección 2.1.6, y la flecha del Q4 es menor que la del Q9 en las cinco mallas, que es la firma del bloqueo por cortante»; «la reimportación DXF no creó entidades nuevas y la memoria se generó para Q4 y Q9 sin error». En 3.2, donde el veredicto ya está escrito, basta con dar la cifra: «[...]\allowbreak{} y por encima de la cota empírica mínima de la Subsección 2.1.6 (1,4 en Q4 y 1,9 en Q9)».\par
\rotulo{También en}\begin{itemize}[leftmargin=1.2em,itemsep=1pt,topsep=1pt]
\item {\small §3.2 (tasa de la tensión recuperada), f. 64: \cita{y por encima de la cota empírica mínima de la subsección 2.1.6}}
\item {\small §3.8 (cláusula b), f. 84: \cita{la membrana de Cook reproduce el bloqueo por cortante del elemento Q4 frente a la convergencia del Q9, cuantificado a igualdad de grados de libertad}}
\item {\small §3.4 (equilibrio de reacciones), f. 69: \cita{La suma de las reacciones verticales en los dos apoyos equilibra la carga total}}
\item {\small §3.5 (membrana de Cook), f. 71: \cita{El elemento Q4 subestima la respuesta en más de un}}
\end{itemize}
{\footnotesize\color{gris}\rotulo{Verificación}ajustado: La evidencia es literal y está en su página (PDF 83, folio 68), y el defecto de fondo existe: hay criterios de aceptación de 2.1.6 con cifra medida y sin veredicto escrito. Pero dos de los elementos del diagnóstico se caen al cotejarlos: (i) la tasa mínima de la tensión recuperada SÍ tiene veredicto en 3.2 —«por encima de la cota…\par}
\end{ficha}
{\small\rotulo{Párrafos que pide la defensa}\par}
\begin{ficha}{media}{\casilla\textbf{DEF-17}\quad\chip{media}{media}\quad Cap. 3 — entradilla · folio 61 · PDF 76\hfill 35 min}
\rotulo{Dice}\cita{«la distinción metodológica clásica entre verificación y validación» (verificada literal en PDF 76 /\allowbreak{} folio 61), seguida de «presenta dos validaciones independientes».}\par
\rotulo{Problema}La entradilla adopta la distinción de Oberkampf y a renglón seguido llama «validación» a dos contrastes sin dato experimental alguno —una solución analítica del mismo modelo matemático y un valor de referencia numérico—, que en ese marco son verificación de la solución. La acotación existe, pero en las Conclusiones (folio 89) y no en el punto de uso, donde el especialista en métodos numéricos la echará de menos.\par
\rotulo{Corrección · reformular}Añadir en la entradilla la acepción que el trabajo adopta —verificación como contraste con soluciones exactas del propio modelo matemático y validación, en sentido acotado, como contraste con referencias externas de aceptación general— y consignar allí mismo que la validación en sentido estricto, contra datos experimentales, queda fuera del alcance y se recoge entre las limitaciones; ajustar «dos validaciones independientes» a «dos contrastes con referencias externas».\par
\rotulo{También en}\begin{itemize}[leftmargin=1.2em,itemsep=1pt,topsep=1pt]
\item {\small §3.4 (título), f. 67: \cita{Validación con la viga de Timoshenko}}
\item {\small §3.5 (título), f. 71: \cita{Validación con la membrana de Cook y el bloqueo por cortante}}
\end{itemize}
\end{ficha}
\begin{ficha}{media}{\casilla\textbf{DEF-18}\quad\chip{media}{media}\quad §3.2 Verificación por soluciones manufacturadas · folio 64 · PDF 79\hfill 25 min}
\rotulo{Dice}\cita{«Que no alcance O(h\textasciicircum{}2) se atribuye a la capa de contorno, donde los nodos reciben el promedio de un solo lado y pierden la cancelación de errores que los nodos interiores obtienen de sus cuatro elementos vecinos» (verificada literal en PDF 80 /\allowbreak{} folio 65; el hallazgo bruto la situaba en PDF 79 /\allowbreak{} folio 64).}\par
\rotulo{Problema}Afirmación sobrecalibrada en el cuerpo: la explicación del orden anómalo del Q4 se presenta como causa establecida cuando es una conjetura que el propio trabajo podía contrastar recomputando la norma sin la franja de contorno, y el texto no dice que no se hizo. En U6 no hay tabla, figura ni guion que respalde el mecanismo propuesto.\par
\rotulo{Corrección · reformular}Marcar el estatus de la explicación sin ampliarla: «Una explicación plausible, que este trabajo no contrastó, es la capa de contorno, donde los nodos reciben el promedio de un solo lado y pierden la cancelación de errores que los nodos interiores obtienen de sus cuatro elementos vecinos». Y dejar la comprobación pendiente en una línea de trabajo futuro: recomputar la norma L2 del campo de tensiones excluyendo la franja de contorno y verificar si el orden se aproxima entonces a 2.\par
{\footnotesize\color{gris}\rotulo{Verificación}ajustado: La cita es literal, pero está en PDF 79 /\allowbreak{} folio 64, no en PDF 80 /\allowbreak{} folio 65: la «corrección» que traía el hallazgo invirtió la ubicación buena. El defecto existe —la capa de contorno se ofrece como causa del orden anómalo sin evidencia ni contraste, y el texto no dice que no se comprobó—, pero está atenuado por su propio entorno: la…\par}
\end{ficha}
\begin{ficha}{media}{\casilla\textbf{DEF-19}\quad\chip{media}{media}\quad Tabla de síntesis del capítulo 3, fila «Timoshenko sigma\_\allowbreak{}x (error máx.)» · folio 72 · PDF 87\hfill 45 min}
\rotulo{Dice}\cita{Timoshenko sigma\_\allowbreak{}x (error máx.) — 0,0414 \%}\par
\rotulo{Problema}El «error máximo» se calcula sobre tres puntos elegidos sin criterio declarado, nunca sobre el dominio completo, y esa cifra local asciende a la tabla de síntesis rotulada «error máx.», donde un lector la leerá como el error máximo del software en el caso. Es una etiqueta que induce a error sobre una de las cifras que el tribunal citará.\par
\rotulo{Corrección · corregir}Cambiar el rótulo de esa fila por «Timoshenko sigma\_\allowbreak{}x — error máximo entre los tres puntos de contraste» y añadir a la nota al pie de la tabla de síntesis que los errores se evalúan en los puntos A, B y C de la Tabla 3.2 y no sobre todos los nodos de la malla. No hace falta tocar la Subsección 3.4: ya declara cómo se ubican los puntos y ya acota la cifra a esos tres.\par
\rotulo{También en}\begin{itemize}[leftmargin=1.2em,itemsep=1pt,topsep=1pt]
\item {\small Tabla de síntesis del capítulo, f. 72: \cita{Timoshenko sigma\_\allowbreak{}x (error máx.) --- 0,0414 \%}}
\end{itemize}
{\footnotesize\color{gris}\rotulo{Verificación}ajustado: La cita es literal y está en PDF 83 /\allowbreak{} folio 68 (es el pie de la Tabla 3.2), pero dos de las tres piezas del diagnóstico se refutan contra el documento. (i) El criterio de elección de los puntos sí se declara: «La Tabla 3.2 compara la tensión normal sigma\_\allowbreak{}x en tres puntos de referencia de la viga —uno en la fibra extrema del centro del…\par}
\end{ficha}
\subsubsection*{Medias y bajas}
{\small\begin{longtable}{@{}L{1.6cm}L{1.6cm}L{5.25cm}L{5.65cm}r@{}}
\toprule
\sffamily\bfseries ID & \sffamily\bfseries Dónde & \sffamily\bfseries Qué pasa & \sffamily\bfseries Qué hacer & \sffamily\bfseries min\\
\midrule\endhead
\bottomrule\endfoot
\mbox{\casilla RED-15}\newline {\color{media}\sffamily\bfseries\scriptsize MEDIA} & f. 61\newline{\scriptsize Redundancia} & El lector recorre el índice del capítulo por partida doble dentro de un mismo párrafo, sin que la segunda enumeración añada nada a la primera salvo el detalle. & Conservar solo la enumeración detallada («La exposición recorre primero…») y dejar la primera oración en su función de encuadre: «Este capítulo presenta y analiza los resultados de la investigación según el diseño metodológico del capítulo anterior.» Ahorro: cuatro líneas. & 15\\\addlinespace[3pt]
\mbox{\casilla RED-16}\newline {\color{media}\sffamily\bfseries\scriptsize MEDIA} & f. 62\newline{\scriptsize Redundancia} & §3.1 reproduce casi íntegro el contenido de §2.1.5: los mismos instrumentos deterministas y reejecutables, la misma precaución de cuadratura con el mismo paréntesis, la misma prueba de reproducción exacta del extrapolador y la misma declaración de autoría de las tablas. La sección declara ella misma que el dato ya fue declarado en §2.1.5, con lo que el propio texto admite la duplicación. & Aparición canónica: §2.1.5 (folios 45-46, PDF 60-61), donde el procedimiento se define. En §3.1 (folio 62) conservar el primer párrafo -qué magnitudes se observan en este capítulo y con qué indicadores- y sustituir el segundo por: «La confiabilidad de estas mediciones se apoya en las tres vías descritas en la subsección 2.1.5:… & 30\\\addlinespace[3pt]
\mbox{\casilla CON-10}\newline {\color{media}\sffamily\bfseries\scriptsize MEDIA} & f. 65\newline{\scriptsize Consistencia} & El párrafo anuncia dos pruebas y la sección desarrolla tres; la tercera es precisamente la que sostiene los criterios de interoperabilidad del tercer objetivo específico, de modo que el lector que se fíe del anuncio no la espera donde el trabajo la necesita. & Sustituir por: «Tres pruebas de regresión resultan especialmente relevantes: el ciclo de transformación entre tipos de elemento, la separación entre identidad de nodo e índice de grado de libertad, y la interoperabilidad de los formatos de intercambio.» & 5\\\addlinespace[3pt]
\mbox{\casilla CON-21}\newline {\color{baja}\sffamily\bfseries\scriptsize BAJA} & f. 65\newline{\scriptsize Consistencia} & La leyenda redondea a 1,5 la tasa que la tabla y el texto dan como 1,54 en tres apariciones, y la llama «observada»: un lector atento no sabrá cuál de las dos es la medida. & Escribir «pendiente de referencia O(h\textasciicircum{}1,5), próxima a la tasa medida de 1,54», o bien usar 1,54 en la leyenda. & 5\\\addlinespace[3pt]
\mbox{\casilla TRZ-28}\newline {\color{baja}\sffamily\bfseries\scriptsize BAJA} & f. 66\newline{\scriptsize Trazabilidad} & La remisión al «tercer objetivo específico» no lleva número ni referencia cruzada, de modo que el lector debe volver a buscarlo en la Introducción. & Sustituir por una remisión con el enunciado abreviado: «que exige el tercer objetivo específico, relativo al intercambio de datos con otras herramientas (Introducción, Objetivos)». & 10\\\addlinespace[3pt]
\mbox{\casilla CLA-21}\newline {\color{media}\sffamily\bfseries\scriptsize MEDIA} & f. 69\newline{\scriptsize Claridad} & El dato aparece sin decir para qué se menciona ni qué se concluye de él, cuando es justamente la clave que explica por qué las dos referencias analíticas difieren entre sí: el lector no especialista lo lee como un dato descriptivo más. & Enlazarlo con el párrafo: «Con una relación luz-peralte de 11,7, la viga es lo bastante esbelta para que la deformación por cortante aporte poco, pero no tanto como para despreciarla: de ahí la diferencia entre el 0,2633 \% de error frente a la solución con corrección de cortante y el 1,85 \% frente a la fórmula de… & 15\\\addlinespace[3pt]
\mbox{\casilla CON-11}\newline {\color{media}\sffamily\bfseries\scriptsize MEDIA} & f. 71\newline{\scriptsize Consistencia} & El mismo folio cifra la incertidumbre de la referencia en 0,05 \% (leyenda) y fija el umbral de significación en 0,1 \% (nota al pie) sin explicar la duplicación, y el valor calculable a partir de la propia Tabla 3.4 es (23,961 - 23,949)/\allowbreak{}23,96 = 0,05 \%. La ambigüedad alimenta además el error de CON-02 en la página siguiente, de modo que conviene corregir los dos pasajes a la vez. & Unificar en el valor calculado: «Se adopta 23,96 como referencia con la salvedad de que su incertidumbre, del orden del 0,05 \% -la diferencia entre las dos mallas Q9 más finas-, es comparable al error de esas mallas: por debajo de ese 0,05 \% las diferencias no son distinguibles de la incertidumbre de la propia referencia.» & 10\\\addlinespace[3pt]
\mbox{\casilla CON-12}\newline {\color{media}\sffamily\bfseries\scriptsize MEDIA} & f. 71\newline{\scriptsize Consistencia} & La sección de Cook separa las coordenadas con coma mientras el resto del capítulo usa punto y coma, de modo que en un texto con coma decimal «(48, 52)» se lee como un único número -cuarenta y ocho coma cincuenta y dos- en vez de un par de coordenadas, justo en el punto donde se mide la cantidad de interés. & Unificar en punto y coma, como ya hacen 3.2 y 3.4: (0; 0), (48; 44), (48; 60), (0; 44) y el punto de interés (48; 52). Corregir también la leyenda y el rótulo de la Tabla 3.4 y la remisión de la tabla de síntesis. & 15\\\addlinespace[3pt]
\mbox{\casilla CON-13}\newline {\color{media}\sffamily\bfseries\scriptsize MEDIA} & f. 71\newline{\scriptsize Consistencia} & El documento declara como regla propia que dice «restricción» y no «condición de contorno», y veintitrés folios después la incumple en la discusión de la membrana de Cook, el único lugar del documento donde reaparece la expresión proscrita (búsqueda exhaustiva: dos apariciones en 168 páginas, la del folio 48 solo para descartarla). & Sustituir por «donde la restricción cambia bruscamente de desplazamiento impuesto a borde libre». Queda así una sola aparición de «condición de contorno» en todo el documento: la del folio 48, que la menciona para descartarla. & 5\\\addlinespace[3pt]
\mbox{\casilla APF-12}\newline {\color{media}\sffamily\bfseries\scriptsize MEDIA} & f. 71\newline{\scriptsize Aparato formal} & La justificación del valor de referencia del único caso sin solución analítica cerrada viaja entera en letra pequeña, en una sola cadena de seis oraciones. La leyenda de la Tabla 3.4 sí declara la incertidumbre del orden de 0,05 por ciento, de modo que el lector no queda a ciegas, pero la razón de esa cifra —y la verificación propia que la respalda— es material de cuerpo: es lo que separa «comparamos con un número de la literatura» de «comprobamos ese número por nuestra cuenta». & Subir la nota al cuerpo como párrafo propio inmediatamente anterior a la tabla, dividida en tres oraciones: qué es el valor de referencia y por qué no es exacto; la corroboración propia por extrapolación de Richardson, con el valor obtenido y su cita con página; y la incertidumbre resultante con su consecuencia operativa, que… & 30\\\addlinespace[3pt]
\mbox{\casilla CLA-22}\newline {\color{media}\sffamily\bfseries\scriptsize MEDIA} & f. 72\newline{\scriptsize Claridad} & Las leyendas de la Figura 3.4 y de la Figura D.2 no se limitan a identificar lo que la imagen muestra: cierran con una conclusión interpretativa que en APA 7 corresponde al cuerpo del texto, no al título del elemento flotante. Además duplican una lectura que el párrafo anterior ya hace. Es una desviación formal, no un error de contenido. & Recortar las leyendas a lo descriptivo: «Convergencia de uy de la membrana de Cook para Q4 y Q9 frente al número de grados de libertad» y «Perfil de σx(y) en x = 0 (centro del vano) de la viga de Timoshenko». La lectura interpretativa ya está en el cuerpo en ambos casos, de modo que basta suprimirla de la leyenda; si se quiere… & 15\\\addlinespace[3pt]
\mbox{\casilla APF-13}\newline {\color{media}\sffamily\bfseries\scriptsize MEDIA} & f. 72\newline{\scriptsize Aparato formal} & Dos piezas de resultados están alojadas fuera de las secciones de resultados. La tabla comparativa Q4/\allowbreak{}Q9 —el elemento más citable del Capítulo 3— queda dentro de la sección de un caso particular, anunciada por una línea y sin conclusión que la lea. Y las mediciones de desempeño, que la Tabla 2.1 declara como indicador observado, aparecen por primera vez dentro de las limitaciones, con razones de 1,7, 2,0 y 2,9 que son resultados nuevos y no la discusión de un resultado ya presentado. & Llevar la tabla comparativa y su anuncio a la sección de interpretación de resultados y acompañarla de dos oraciones de lectura: cuándo conviene el Q4 (mallas gruesas de exploración, menor coste por elemento) y cuándo el Q9 (flexión dominante, geometría distorsionada, exactitud por grado de libertad). Llevar la tabla de tiempos… & 55\\\addlinespace[3pt]
\end{longtable}}
\subsection{Preliminares e índices}
{\small\color{gris}\sffamily 0 correcciones críticas o altas y 2 medias o bajas en esta sección · 2,2 h · más 0 frases de calibración (sección 3)}\par
\subsubsection*{Medias y bajas}
{\small\begin{longtable}{@{}L{1.6cm}L{1.6cm}L{5.25cm}L{5.65cm}r@{}}
\toprule
\sffamily\bfseries ID & \sffamily\bfseries Dónde & \sffamily\bfseries Qué pasa & \sffamily\bfseries Qué hacer & \sffamily\bfseries min\\
\midrule\endhead
\bottomrule\endfoot
\mbox{\casilla APF-04}\newline {\color{media}\sffamily\bfseries\scriptsize MEDIA} & f. xii-xiv\newline{\scriptsize Aparato formal} & La lista de símbolos y siglas promete reunir la notación del documento y no cumple en ninguno de los dos sentidos. Por defecto omite más de veinte símbolos que aparecen en ecuaciones numeradas del Capítulo 1 y en las tablas de resultados (X\_\allowbreak{}e, u\_\allowbreak{}e, n, δu, M y E, f\_\allowbreak{}e\textasciicircum{}vol, F\_\allowbreak{}n, F\_\allowbreak{}Γ, la partición K\_\allowbreak{}ff/\allowbreak{}K\_\allowbreak{}fr/\allowbreak{}K\_\allowbreak{}rf/\allowbreak{}K\_\allowbreak{}rr y u\_\allowbreak{}f/\allowbreak{}u\_\allowbreak{}r, L, q\_\allowbreak{}s y q\_\allowbreak{}e, n\_\allowbreak{}e, R\_\allowbreak{}J, AR, T\_\allowbreak{}R, X\_\allowbreak{}1 a X\_\allowbreak{}3, r\_\allowbreak{}i, u\_\allowbreak{}h, σ*\_\allowbreak{}h, u\_\allowbreak{}M) y, sobre todo, el operador de ensamblaje de la ecuación (1.17), que tampoco nombra el texto: la misma ecuación escribe el primer miembro con una A caligráfica muda y el segundo con el sumatorio corriente, contraste que un tribunal no especialista no puede resolver. Por exceso incluye GUI, que el documento no emplea nunca —siempre escribe «interfaz gráfica»—, coloca TP y DP entre las siglas cuando solo existen como subíndices de la matriz constitutiva, y gasta dos líneas en explicar por qué no usa los códigos FEA/\allowbreak{}FEM, aclaración que pertenece a la tabla de especificaciones. A la vez faltan CAD, ICMJE y NLM, que sí se usan: CAD dentro de la propia lista de abreviaturas, e ICMJE y NLM en el folio 11 sin desarrollar. & Nombrar el operador en la frase que precede a la ecuación (1.17): «La matriz de rigidez global se obtiene con el operador de ensamblaje, que denotamos A y que no es otra cosa que la suma de cada aporte elemental en las filas y columnas que le corresponden según el mapeo de GDL:», y dar de alta en la tabla de símbolos la entrada… & 90\\\addlinespace[3pt]
\mbox{\casilla CON-04}\newline {\color{media}\sffamily\bfseries\scriptsize MEDIA} & f. xiii\newline{\scriptsize Consistencia} & Cinco colisiones o descuidos de notación frente a la Nomenclatura: Q designa la compacidad mientras Q4 y Q9 designan la familia de elementos; E designa el módulo de elasticidad en 1.4 y la matriz de extrapolación en 1.11, distinguidas solo por la negrita; n es a la vez número de nodos del elemento e índice del sumatorio de ensamblaje; el subíndice de w\_\allowbreak{}p reutiliza la p del orden polinómico; y las normas del error llevan «según magnitud» en la columna de unidades, que no es una unidad. Además la lista no declara su criterio de ordenación, que no es alfabético, de modo que el lector que busque un símbolo suelto no lo encuentra. & Nombres canónicos recomendados: compacidad q\_\allowbreak{}C (coherente con q\_\allowbreak{}SJ) o sin símbolo; peso de la cuadratura w\_\allowbreak{}g (coherente con n\_\allowbreak{}g); matriz de extrapolación P, de proyección a los nodos, o E con advertencia expresa en su primera aparición y en la Nomenclatura; índice j en el sumatorio de ensamblaje. Sustituir «según magnitud» por… & 40\\\addlinespace[3pt]
\end{longtable}}
\subsection{Anexos}
{\small\color{gris}\sffamily 0 correcciones críticas o altas y 9 medias o bajas en esta sección · 3,6 h · más 0 frases de calibración (sección 3)}\par
{\small\rotulo{Párrafos que pide la defensa}\par}
\begin{ficha}{media}{\casilla\textbf{DEF-25}\quad\chip{media}{media}\quad Anexo C — apertura · folio 118 · PDF 133\hfill 10 min}
\rotulo{Dice}\cita{«Se reproducen aquí fragmentos representativos del motor numérico, escogidos por condensar la formulación isoparamétrica descrita» (verificada literal en PDF 133 /\allowbreak{} folio 118).}\par
\rotulo{Problema}El anexo no aclara si los fragmentos están copiados literalmente del repositorio o adaptados y abreviados para la lectura, de modo que el tribunal puede preguntar si el código mostrado es exactamente el que ejecuta el software; en un trabajo cuya bandera es la reproducibilidad, la ambigüedad es gratuita.\par
\rotulo{Corrección · definir}Agregar la cláusula explícita —«se transcriben tal como figuran en el repositorio, sin abreviar ni renombrar variables»— o, si se editaron por espacio, decirlo y señalar qué se omitió.\par
\end{ficha}
\subsubsection*{Medias y bajas}
{\small\begin{longtable}{@{}L{1.6cm}L{1.6cm}L{5.25cm}L{5.65cm}r@{}}
\toprule
\sffamily\bfseries ID & \sffamily\bfseries Dónde & \sffamily\bfseries Qué pasa & \sffamily\bfseries Qué hacer & \sffamily\bfseries min\\
\midrule\endhead
\bottomrule\endfoot
\mbox{\casilla RED-19}\newline {\color{media}\sffamily\bfseries\scriptsize MEDIA} & f. 109\newline{\scriptsize Redundancia} & Los ocho módulos educativos se describen tres veces —§2.2.7, §B.3 y §B.7— y su tabla se imprime dos veces con el mismo contenido. §B.3 anticipa cuatro páginas antes lo que §B.7 desarrolla, incluida la descripción de la capa superpuesta y del conmutador Fórmula–Valores. & Aparición canónica de la descripción detallada: §B.7. Fundir las dos tablas en una sola —la del anexo, que ya incluye fase, concepto y atajo— y sustituir la tabla del capítulo 2 por la remisión a ella, conservando en §2.2.7 solo el párrafo que explica el papel de los módulos en la arquitectura. En §B.3 dejar el atajo y la… & 30\\\addlinespace[3pt]
\mbox{\casilla CON-17}\newline {\color{media}\sffamily\bfseries\scriptsize MEDIA} & f. 109\newline{\scriptsize Consistencia} & El concepto vertebrador del trabajo, «canal de cálculo» -27 apariciones y parte del objetivo general-, alterna con «canal del MEF» (folio 17), «flujo del Método de Elementos Finitos» (folio 109) y «proceso de cálculo» (folios 2 y 86), justo donde el lector busca confirmar que se le habla de lo mismo. La misma oración del folio 109 es además la única de las ocho que escribe «Método de Elementos Finitos» sin artículo, frente a las siete «Método de los Elementos Finitos» del resto. & Nombre canónico: «canal de cálculo», sin excepción. En el Anexo B.7, «distribuidos según las fases del canal de cálculo del MEF», que resuelve de paso el artículo del nombre desarrollado; en la nota de la Tabla 1.1, «el recurso recorre el canal de cálculo de extremo a extremo»; sustituir las dos apariciones de «proceso de… & 30\\\addlinespace[3pt]
\mbox{\casilla CON-18}\newline {\color{media}\sffamily\bfseries\scriptsize MEDIA} & f. 110\newline{\scriptsize Consistencia} & El mismo componente se llama «validador de salud» en el cuerpo (5 apariciones, folios 33, 44, 56, 73 y 78, entre ellas la celda de evidencia de la matriz de principios) y «comprobador de salud» en el Anexo B que lo documenta (8 apariciones), de modo que el lector que siga la evidencia del Capítulo 3 hasta el anexo cree encontrarse con dos instrumentos distintos. Además «validador» refuerza la ambigüedad de «validación» señalada en CON-03. & Nombre canónico: «comprobador de salud», el que usan la interfaz y el anexo de referencia; sustituirlo en las cinco apariciones del cuerpo (folios 33, 44, 56, 73 y 78). Conservar «salud del modelo» solo como nombre del diálogo. Añadir en la celda de evidencia de la matriz de principios el localizador «Anexo B.8», que hoy falta. & 25\\\addlinespace[3pt]
\mbox{\casilla APF-15}\newline {\color{media}\sffamily\bfseries\scriptsize MEDIA} & f. 110-114\newline{\scriptsize Aparato formal} & Las capturas de los módulos M1 a M6 se presentan en bloque y quedan luego solas con su leyenda: ninguna se referencia individualmente ni desde la Tabla 2.5 ni desde la Subsección 2.2.7, que es donde el lector del cuerpo esperaría el puente al anexo. El lector no sabe qué debe mirar en cada captura, y la leyenda carga sola con toda la explicación de un módulo interactivo reducido a una imagen fija. & Añadir en la Subsección 2.2.7 o en la Tabla 2.5 una remisión al bloque: «la interfaz de cada módulo se muestra en el Anexo B, Figuras B.7 a B.12». Y en el anexo, después de cada figura (o cada dos), una oración breve que la cite por su número y señale el rasgo puntual que el lector debe mirar: «la Figura B.9 muestra, en la… & 40\\\addlinespace[3pt]
\mbox{\casilla RED-22}\newline {\color{baja}\sffamily\bfseries\scriptsize BAJA} & f. 125\newline{\scriptsize Redundancia} & El anexo D reproduce la tabla de tensiones en los puntos A, B y C que la tabla 3.x del capítulo 3 ya presenta completa con sus errores frente a las dos referencias, en lugar de limitarse a ampliar la evidencia que el cuerpo no reporta. La frase metodológica sobre la cuadratura que también se repite allí se trata en RED-12. & Aparición canónica de la tabla: capítulo 3, sección 3.4 (folio 68). En el anexo D suprimir la tabla repetida y sustituirla por la remisión «los valores en los puntos de referencia figuran en la tabla 3.x»; reservar el anexo para las componentes y configuraciones que el cuerpo no reporta. Ahorro: una tabla y un cuarto de página.… & 25\\\addlinespace[3pt]
\mbox{\casilla CON-23}\newline {\color{baja}\sffamily\bfseries\scriptsize BAJA} & f. 125\newline{\scriptsize Consistencia} & La prosa del Anexo D redondea a un decimal (0,9 \%) el error de sigma\_\allowbreak{}y en el punto B que la Tabla 3.2 y la propia Tabla D.4 reportan con dos decimales (0,94 \%), una precisión distinta para el mismo dato a pocas líneas de la tabla que lo contiene; el recálculo independiente da 0,939 \%. & Escribir «0,94 \%» en la oración del Anexo D, como en la Tabla 3.2 y en la Tabla D.4, en vez de redondear a un decimal solo en ese punto del texto. & 5\\\addlinespace[3pt]
\mbox{\casilla ACC-13}\newline {\color{media}\sffamily\bfseries\scriptsize MEDIA} & f. 128\newline{\scriptsize Accesibilidad} & El párrafo de apertura, obra del autor, corrige solo tres erratas puntuales y una frase técnica, pero no advierte que el resto de las diez páginas incrustadas conserva mayúsculas sin tilde, minúscula inconsistente de «sap2000» y errores de digitación: el lector recibe sin aviso general un registro muy distinto al del resto de la tesis y puede atribuírselo al autor. & Sustituir la lista de tres erratas por una salvedad general: «se conserva además el estilo tipográfico y ortográfico original del documento —mayúsculas sin tilde, abreviaturas propias y erratas de digitación—, ajeno a las convenciones del resto de la tesis». & 15\\\addlinespace[3pt]
\mbox{\casilla CON-19}\newline {\color{media}\sffamily\bfseries\scriptsize MEDIA} & f. 130\newline{\scriptsize Consistencia} & El Anexo E -PDF 143-154, folios 128-139; no confundir con el Anexo D, PDF 136-142- rompe a la vez la terminología, la notación y la tipografía del documento: dice «esfuerzo» donde el resto dice «tensión» (7 de las 8 apariciones de «esfuerzo» del documento están aquí), compone los símbolos con una fuente matemática distinta y en negrita, usa letras que la Nomenclatura no declara (c, l, I, V, H, f'c) y titula en mayúsculas sin tildes. Al ser un documento externo reproducido, el contraste se explica, pero el trabajo no lo advierte en ninguna parte y el lector lo toma por texto propio descuidado. & Opción mínima y recomendada: abrir el anexo con una nota que declare su naturaleza y su glosario -«Este anexo reproduce el documento de cálculo elaborado con SAP2000 que origina las referencias del Anexo D. Se conserva su redacción original; en él, «esfuerzo» equivale a «tensión» y las letras c, l, I, V y H designan el… & 25\\\addlinespace[3pt]
\mbox{\casilla APF-16}\newline {\color{media}\sffamily\bfseries\scriptsize MEDIA} & f. 148\newline{\scriptsize Aparato formal} & El anexo se presenta como memoria paso a paso trazable a la formulación teórica, pero el puente entre la sustitución numérica y la definición simbólica solo existe en dos de siete eslabones. El lector que quiera comprobar de dónde sale una matriz tiene que buscarla en el Capítulo 1 por su aspecto, que es justamente el trabajo que el anexo dice ahorrarle. & Agregar una referencia de ecuación junto a cada fórmula reutilizada —D, J y J\textasciicircum{}-1, B, k\_\allowbreak{}e y la reducción K\_\allowbreak{}ff—, con el mismo formato que ya usan la extrapolación y von Mises, de modo que cada paso numérico del anexo apunte a su ecuación numerada del Capítulo 1. & 20\\\addlinespace[3pt]
\end{longtable}}
\subsection{Transversal}
{\small\color{gris}\sffamily 0 correcciones críticas o altas y 7 medias o bajas en esta sección · 3,5 h · más 0 frases de calibración (sección 3)}\par
\subsubsection*{Medias y bajas}
{\small\begin{longtable}{@{}L{1.6cm}L{1.6cm}L{5.25cm}L{5.65cm}r@{}}
\toprule
\sffamily\bfseries ID & \sffamily\bfseries Dónde & \sffamily\bfseries Qué pasa & \sffamily\bfseries Qué hacer & \sffamily\bfseries min\\
\midrule\endhead
\bottomrule\endfoot
\mbox{\casilla TRZ-08}\newline {\color{media}\sffamily\bfseries\scriptsize MEDIA} & f. 5 (evidencia); 89-90 (párrafo de cierre, que sí lleva la salvedad)\newline{\scriptsize Trazabilidad} & La finalidad con que cierra el objetivo general es un efecto sobre el estudiante que el trabajo no midió; el final de 3.8 sí lleva la salvedad («que ese criterio mejore de hecho [...]\allowbreak{} no se contrasta en esta etapa»), pero las Conclusiones declaran el objetivo cumplido dos veces sin repetirla y el título de la tesis tampoco contiene esa finalidad, de modo que el eslabón objetivo general {\mates →} resultado {\mates →} conclusión se cierra con más de lo que se contrastó. & Añadir la remisión in situ al objetivo general, sin tocar el molde: «[...]\allowbreak{} obtenida por el MEF. El alcance de esta etapa llega al diseño y al contenido del software; la mejora del criterio se deja a la prueba de campo que se diseña en las recomendaciones (véase Alcance y limitaciones)». No modificar las Conclusiones: tanto el… & 30\\\addlinespace[3pt]
\mbox{\casilla TRZ-09}\newline {\color{media}\sffamily\bfseries\scriptsize MEDIA} & f. 43 y 77\newline{\scriptsize Trazabilidad} & La cuenta de los 49 ítems del consenso no cierra: el universo toma 18 (quince destrezas FEA y tres conceptos FEM) y la nota de la Tabla 3.7 excluye 22, de modo que nueve quedan sin contabilizar —FEA1, FEA2, FEA4, FEA5, FEA11, FEA12, FEA20, FEA24 y FEM3 si la numeración de los códigos es correlativa, supuesto que no se verificó contra la fuente—, y tampoco se aclara si TYPE1 a TYPE7 forman parte de los 49. El hueco deja abierta la objeción de selección a conveniencia justo sobre el universo en que descansa la cláusula (c). & Cerrar la cuenta por escrito en 2.1.4 con una oración del tipo: «De los 49 ítems del consenso, 18 forman el universo —quince destrezas y tres conceptos— y los 31 restantes quedan fuera del alcance declarado; la nota de la Tabla 3.7 los lista uno a uno». Antes de escribirla, cotejar la numeración real de los códigos en las… & 45\\\addlinespace[3pt]
\mbox{\casilla RED-04}\newline {\color{media}\sffamily\bfseries\scriptsize MEDIA} & f. 13 (recorte) /\allowbreak{} 18 (canonica)\newline{\scriptsize Redundancia} & El diagnóstico de la caja negra —textos abstractos por un lado, paquetes que ocultan el procedimiento por otro, software didáctico como puente— se desarrolla desde cero cinco veces. Las dos del capítulo 1 (§1.1 y §1.2, a cinco páginas) son casi literales entre sí e incluso repiten la misma tríada de referencias. Ubicación corregida respecto del reporte bruto: la cita no está en PDF 17, sino en PDF 28 y 33. & Invertir la asignación respecto de la propuesta original. Aparición canónica: §1.2, primer principio (folio 18, PDF 33), que se conserva íntegra porque sostiene el principio de visibilidad del procedimiento y es la única con localizadores de página ([9, p. 157; 10, p. 872; 11, p. 2]\allowbreak{}). El recorte va en §1.1 (folio 13, PDF 28):… & 30\\\addlinespace[3pt]
\mbox{\casilla ACC-04}\newline {\color{media}\sffamily\bfseries\scriptsize MEDIA} & f. 16-17\newline{\scriptsize Accesibilidad} & La fila que más diferencia a EduFEM de los recursos comparados —y con ella el vacío que justifica la tesis— se lee con códigos internos de módulo, símbolos sin glosa y nombres de casos de verificación que el tribunal no especialista no puede interpretar sin haber leído antes los capítulos 2 y 3. Una tabla se consulta fuera del hilo del texto: si no se sostiene sola, el argumento que carga se pierde. & Añadir al pie de la Tabla 1.1 (o ampliar la nota existente) las glosas mínimas de los códigos y símbolos que aparecen en sus celdas: «M0 a M7 son los ocho módulos educativos de EduFEM, descritos en la sección 2.2. det J es el determinante del Jacobiano, indicador de la distorsión del elemento (sección 1.6). kappa\_\allowbreak{}2 es el número… & 25\\\addlinespace[3pt]
\mbox{\casilla APF-01}\newline {\color{media}\sffamily\bfseries\scriptsize MEDIA} & f. 15-16\newline{\scriptsize Aparato formal} & La tabla que sostiene el diagnóstico de la brecha —y que las Conclusiones vuelven a invocar como evidencia— mezcla sin declararlo dos bases de evidencia de rango muy distinto: la columna de EduFEM la evalúa el autor sobre su propio producto, con acceso directo al código, y las demás se infieren de publicaciones de terceros. Además gradúa la fila «Cálculo paso a paso» con una escala verbal de cuatro niveles (Alta, Media-alta, Media, Baja) que no se define en ninguna parte del documento. Leída como está, la tabla se presenta al tribunal como una comparación simétrica de once filas cuando no lo es, y el criterio que separa «Media-alta» de «Media» queda a discreción del autor sobre su propia herramienta. & Añadir dos frases a la nota al pie de la Tabla 1.1, a continuación de la que ya explica «no consta»: «La columna de EduFEM procede de la verificación directa del programa desarrollado en este trabajo; las demás, de las publicaciones citadas en cada encabezado, de modo que la comparación no es simétrica en su base de evidencia.… & 25\\\addlinespace[3pt]
\mbox{\casilla APF-02}\newline {\color{media}\sffamily\bfseries\scriptsize MEDIA} & f. 86-89\newline{\scriptsize Aparato formal} & Los cinco párrafos de conclusiones específicas sí remiten al capítulo y a la sección que los desarrolla, pero ninguno nombra ni numera el objetivo específico al que responde: abren con paráfrasis («el diagnóstico de la brecha», «la fundamentación teórica»). La correspondencia uno a uno que la frase de apertura promete queda por tanto sin verificar, y es exactamente el eslabón que un tribunal recorre en la defensa —objetivo específico n, conclusión n—. Con cinco objetivos y cinco párrafos, basta con rotularlos para cerrar la trazabilidad. & Rotular cada párrafo con el objetivo que cierra, conservando la remisión al capítulo que ya tiene y copiando el verbo del enunciado original del objetivo, para que el cotejo sea palabra por palabra: «Objetivo específico 1 —diagnosticar la brecha— (Capítulo 1, Sección 1.1): la revisión de los antecedentes…», y así hasta el… & 30\\\addlinespace[3pt]
\mbox{\casilla REF-01}\newline {\color{media}\sffamily\bfseries\scriptsize MEDIA} & f. ii (ancla); 10 y 89-90 (apariciones con defecto); 7 y 94 (apariciones ya trazadas)\newline{\scriptsize Referencias} & La cláusula que acota el alcance de la hipótesis —la afirmación más consultada del documento, y la que separa lo comprobado de lo pendiente— remite a «la literatura» en cinco puntos de máximo peso (Resumen, Introducción, 2.1, Conclusiones, Recomendaciones) y solo en uno de ellos, el de la Introducción, indica dónde buscarla (remisión interna a la Sección 1.2); en ninguno lleva número Vancouver. En la misma página de Conclusiones, además, el «consenso publicado de 67 expertos» que sostiene el resultado central de validez de contenido aparece sin el número [3]\allowbreak{}, que sí acompaña a esa fuente en el cuerpo. El tribunal que quiera verificar sobre qué se apoya la parte no medida de la hipótesis no tiene por dónde empezar en el punto exacto donde se la enuncia. & Dar traza puntual sólo donde hoy no la hay, reutilizando las referencias y los localizadores que la Sección 1.2 y la subsección 3.6.2 ya emplean (verificar cada número y cada página contra el texto ya citado; no introducir páginas nuevas sin cotejar la fuente). (1) Conclusiones, folio 89: «…los tres conceptos del canal que un… & 25\\\addlinespace[3pt]
\end{longtable}}
\subsection{Portada}
{\small\color{gris}\sffamily 0 correcciones críticas o altas y 1 medias o bajas en esta sección · 0,2 h · más 0 frases de calibración (sección 3)}\par
\subsubsection*{Medias y bajas}
{\small\begin{longtable}{@{}L{1.6cm}L{1.6cm}L{5.25cm}L{5.65cm}r@{}}
\toprule
\sffamily\bfseries ID & \sffamily\bfseries Dónde & \sffamily\bfseries Qué pasa & \sffamily\bfseries Qué hacer & \sffamily\bfseries min\\
\midrule\endhead
\bottomrule\endfoot
\mbox{\casilla APF-17}\newline {\color{baja}\sffamily\bfseries\scriptsize BAJA} & f. 1\newline{\scriptsize Aparato formal} & Falta un dato que el formato de presentación de la Carrera puede exigir. No es un defecto confirmado: depende del reglamento vigente, que esta auditoría no consultó. & Comprobar en el reglamento de grado de la Carrera de Ingeniería Civil de la UATF si la portada debe consignar tutor o asesor y, en su caso, añadirlo con el formato que el reglamento indique. Si no lo exige, no hay nada que corregir. & 15\\\addlinespace[3pt]
\end{longtable}}
\clearpage
\section{Patrones transversales}
Defectos que aparecen más de tres veces en el documento: se corrigen con una decisión y una búsqueda, no frase por frase. Aquí van los de claridad y edición, que no pertenecen a ningún capítulo; los patrones de trazabilidad, redundancia, consistencia, accesibilidad, aparato formal y referencias están en su capítulo (sección 4), bajo el folio de su primer ejemplo, y los de calibración, en la sección 3. La lista completa, con recuento y ejemplos, está en la sección 6 de AUDITORIA-REDACCION.md.
\begin{ficha}{media}{\casilla\textbf{CLA-03}\quad\chip{media}{media}\quad Resumen (primera aparición) y 2.2 · folio i · PDF 2\hfill 25 min}
\rotulo{Dice}\cita{«La hipótesis sostiene que la exposición transparente e interactiva del canal de cálculo y de la respuesta estructural, sobre el modelo del propio estudiante, permitirá mejorar ese criterio» (PDF 2, verificado literal; el término reaparece en PDF 21, folio 6, al enunciar la hipótesis). Su única glosa llega dieciséis páginas después (PDF 18).}\par
\rotulo{Problema}«Canal de cálculo» es un término acuñado por el autor y eje de todo el documento —problema, variable independiente, hipótesis, objetivos y PI-1—, pero se usa por primera vez en el Resumen sin glosa alguna y su única definición aparece recién en el objeto de estudio (folio 3). Además, el Capítulo 2 lo diluye con dos usos metafóricos: «un canal de operaciones encadenadas» (el motor) y «un único canal de generación» (la composición del PDF), de modo que el término que nombra la variable independiente compite con dos sentidos genéricos.\par
\rotulo{Corrección · definir}Glosar en la primera aparición del Resumen: «... la exposición transparente e interactiva del canal de cálculo —la cadena de operaciones que va del mapeo isoparamétrico del elemento a la recuperación de tensiones en los nodos— y de la respuesta estructural ...», y repetir la glosa breve en la primera aparición de la Introducción, antes de usar el término para nombrar la variable independiente. Para desambiguar, sustituir solo «un único canal de generación» (folio 60) por «un único procedimiento de generación». NO tocar «un canal de operaciones encadenadas» (folio 53): ahí 'canal' nombra exactamente el canal de cálculo del método aplicado al motor, y cambiarlo rompe la continuidad del término entre el marco y la descripción del software.\par
\rotulo{También en}\begin{itemize}[leftmargin=1.2em,itemsep=1pt,topsep=1pt]
\item {\small Introducción /\allowbreak{} Objeto de estudio (única glosa), f. 3: \cita{el canal de cálculo que va del mapeo isoparamétrico a la recuperación de tensiones}}
\item {\small §2.2.5 Motor de cálculo del MEF, f. 53: \cita{El cálculo se organiza como un canal de operaciones encadenadas}}
\item {\small §2.2.9 Memoria de cálculo e interoperabilidad, f. 60: \cita{ofrece dos estilos que comparten un único canal de generación}}
\end{itemize}
{\footnotesize\color{gris}\rotulo{Verificación}ajustado: Las citas son literales y están en su página: PDF 2 (folio i) y PDF 21 (folio 6) para el uso del término, PDF 18 (folio 3) para la glosa, PDF 68 (folio 53) y PDF 75 (folio 60) para los otros usos de 'canal'. Pero la lectura del contexto rebaja el hallazgo en dos puntos. (1) No es un término clave sin definir: la glosa del objeto de…\par}
\end{ficha}
\begin{ficha}{media}{\casilla\textbf{CLA-11}\quad\chip{media}{media}\quad §1.11 y 1.12 (marco teórico) · folio 30 · PDF 45\hfill 60 min}
\rotulo{Dice}\cita{«por lo que EduFEM la incorpora en todas sus rutas de cálculo (puntos de Gauss, nodos extrapolados, promedio nodal, sonda y memoria de cálculo)» (verificado literal en PDF 45); «el módulo M0 emplea la forma por vértices de Verdict» y «el validador de salud y la memoria de cálculo la evalúan en los puntos de Gauss del elemento» (verificados literales en PDF 48). El software se presenta recién en la Sección 2.2 (PDF 62 y ss.).}\par
\rotulo{Problema}El marco teórico nombra cinco componentes del software —«sonda», «memoria de cálculo», «validador de salud», «módulo M0», «prueba de regresión»— que solo se presentan en el Capítulo 2, de modo que en su primera aparición el lector no especialista no puede saber a qué se refieren. El mismo desborde afecta a la estructura: la Sección 1.11 ocupa cuatro páginas y alterna teoría general (factorización, tensiones principales, superconvergencia) con decisiones y pruebas del programa (orden de numeración de los puntos de Gauss, prueba de regresión, criterio de recálculo de von Mises), de modo que se pierde la línea del marco teórico y se difumina la frontera entre lo que la teoría exige y lo que este trabajo decidió.\par
\rotulo{Corrección · mover}Acción: definir/\allowbreak{}glosar en la primera aparición, no mover. En el Capítulo 1, sustituir por su contenido genérico con remisión los tres nombres que llegan sin presentar: «en todas las rutas de cálculo del programa —tensiones en los puntos de Gauss, valores extrapolados a los nodos, promedio nodal, consulta puntual y memoria de cálculo (§2.2)—» (folio 30); «el módulo M0» por «el módulo introductorio de calidad de malla (§2.2)» (folio 33); «el validador de salud» por «la rutina de comprobación del modelo (§2.2)» (folio 33). No hace falta tocar «memoria de cálculo» (ya presentada en el Resumen y la Introducción) ni «prueba de regresión» (se explica en su propia oración). Opcionalmente, subdividir §1.11 en 1.11.1 «Resolución del sistema y reacciones» y 1.11.2 «Recuperación de tensiones» para aligerar sus cuatro páginas. NO trasladar a §2.2 el párrafo del orden de numeración de los puntos de Gauss ni el criterio de recálculo de las magnitudes no lineales: el primero sostiene un argumento teórico sobre la matriz de extrapolación, y la apertura del capítulo ya declara que la exposición sigue el orden del canal de cálculo del programa.\par
\rotulo{También en}\begin{itemize}[leftmargin=1.2em,itemsep=1pt,topsep=1pt]
\item {\small §1.12 Calidad geométrica de la malla, f. 33: \cita{el módulo M0 emplea la forma por vértices de Verdict}}
\item {\small §1.12 Calidad geométrica de la malla, f. 33: \cita{El validador de salud y la memoria de cálculo la evalúan en los puntos de Gauss del elemento}}
\item {\small §1.11 Resolución y recuperación de tensiones, f. 32: \cita{una prueba de regresión comprueba que}}
\end{itemize}
{\footnotesize\color{gris}\rotulo{Verificación}ajustado: Las tres citas son literales y están en sus páginas (PDF 45 /\allowbreak{} folio 30 y PDF 48 /\allowbreak{} folio 33), y el Capítulo 2 presenta el software recién en PDF 62 y ss., así que el adelanto de nombres existe. Pero el hallazgo está sobrecargado y hay que partirlo. (1) De los cinco componentes que enumera, dos no son términos sin presentar: «memoria de…\par}
\end{ficha}
\begin{ficha}{media}{\casilla\textbf{CLA-15}\quad\chip{media}{media}\quad PATRÓN: oraciones de 70 palabras o más (Introducción, Cap. 1 y Cap. 2; el recuento abarca todo el cuerpo) · folio 41 · PDF 56\hfill 240 min}
\rotulo{Dice}\cita{«El problema científico, en la formulación única de la Introducción —¿de qué manera el uso del software de análisis estructural como caja negra podrá afectar en el bajo criterio [...]\allowbreak{} en la gestión 2026?—, se atiende mediante el objetivo general de desarrollar EduFEM y se contrasta con la hipótesis de que la exposición transparente e interactiva del canal y de la respuesta permitirá mejorar ese criterio, cuya parte comprobable en esta etapa son las tres cláusulas de exposición, corrección y validez de contenido.» (verificado literal en PDF 56)}\par
\rotulo{Problema}El cuerpo acumula 101 oraciones de 70 palabras o más según la medición automática (Cap. 1: 28; Cap. 2: 25; Cap. 3: 28; Introducción: 9; Conclusiones: 11), de las cuales 62 caen en la zona de este informe. El molde es siempre el mismo: dos puntos o punto y coma que introducen una cláusula o un ítem en vez de cerrar la oración, y uno o dos incisos entre rayas anidados que separan el sujeto de su verbo tres o cuatro líneas. El patrón golpea precisamente las frases estructurales —resumen, hipótesis, tipo de investigación, criterios de aceptación, matriz de consistencia—, que son las que el tribunal no especialista debe poder citar sin releer.\par
\rotulo{Corrección · reformular}Mantener las tres reglas de poda —(a) cerrar con punto donde hoy hay dos puntos o punto y coma que introducen una cláusula o un ítem; (b) un solo inciso entre rayas por oración; (c) toda enumeración de cuatro elementos o más sale a lista o tabla— y aplicarlas con prioridad a las frases estructurales ya levantadas en CLAa-01, CLAa-03, CLAa-05, CLAa-07 y CLAa-10. Presentar el alcance sin cifra: en lugar de «101 oraciones de 70 palabras o más», decir «oraciones de más de 70 palabras en los pasajes estructurales de la Introducción y de los capítulos 1 y 2, con cuatro casos verificados de 100 a 180 palabras (folios i, 37, 41 y 46)», y, si se quiere una cifra, medirla sobre los .tex excluyendo tablas, leyendas y listas. Para el folio 41, la reescritura propuesta con un cierre que remite a la prueba de campo: «El problema científico es el enunciado único de la Introducción: ¿de qué manera el uso del software de análisis estructural como caja negra podrá afectar en el bajo criterio para interpretar la respuesta estructural obtenida por el MEF en los estudiantes de la Carrera de Ingeniería Civil de la Universidad Autónoma Tomás Frías en la gestión 2026? Este problema se atiende mediante el objetivo general de desarrollar EduFEM y se contrasta con la hipótesis de que la exposición transparente e interactiva del canal y de la respuesta permitiría mejorar ese criterio. De esa hipótesis, esta etapa comprueba solo las tres cláusulas de exposición, corrección y validez de contenido; el efecto sobre el estudiante queda para la prueba de campo diseñada en las recomendaciones.»\par
\rotulo{También en}\begin{itemize}[leftmargin=1.2em,itemsep=1pt,topsep=1pt]
\item {\small Resumen, f. i: \cita{y la membrana de Cook evidenció el bloqueo por cortante (shear-locking) del Q4 frente a la}}
\item {\small §2.1.1 Tipo de investigación, f. 37: \cita{Por su finalidad, el trabajo es de carácter propositivo}}
\item {\small §2.1.6 Criterios de aceptación, f. 46: \cita{tasas de convergencia del desplazamiento observadas dentro de}}
\end{itemize}
{\footnotesize\color{gris}\rotulo{Verificación}ajustado: El defecto existe y la cita es literal y está donde se dice (PDF 56, folio 41, 2.1.3 Matriz de consistencia; la oración mide unas 105 palabras). Verifiqué además tres instancias más, todas literales y en su página: Resumen (PDF 2, folio i), 2.1.1 Tipo de investigación (PDF 52, folio 37) y 2.1.6 Criterios de aceptación (PDF 61, folio 46,…\par}
\end{ficha}
\begin{ficha}{media}{\casilla\textbf{CLA-19}\quad\chip{media}{media}\quad §3.2-3.5 (verificación y validación) · folio 63 · PDF 78\hfill 45 min}
\rotulo{Dice}\cita{Folio 63: «Se adoptó como configuración base el cuadrado unitario [0, 1]\allowbreak{}2 en tensión plana, con material E = 1,0, ν = 0,3 y espesor t = 1,0 […]\allowbreak{} Se resolvieron mallas estructuradas de N × N elementos para N {\mates ∈} \{2, 4, 8, 16, 32\}\allowbreak{} en cuatro configuraciones».}\par
\rotulo{Problema}El impersonal con «se» recorre las secciones 3.2 a 3.5 (al menos doce construcciones en posición de decisión metodológica: se adoptó, se resolvieron, se midieron, se eligió, se empleó, se verificó, se atribuye, se realiza, se comprueba, se evalúan, se integran, se registra) y borra la distinción entre lo que decidió el autor, lo que calcula el software y lo que sostiene la literatura, justo en el capítulo donde el tribunal necesita saber quién responde por cada dato. No induce a error de contenido, pero deja sin sujeto la cadena de responsabilidad de la verificación.\par
\rotulo{Corrección · reformular}Fijar tres sujetos y usarlos con disciplina en 3.2-3.5: «el autor adoptó /\allowbreak{} eligió /\allowbreak{} no contrastó» para las decisiones de diseño experimental, «EduFEM resuelve /\allowbreak{} reporta /\allowbreak{} recupera» para lo que hace el programa y «la literatura predice /\allowbreak{} advierte» para lo tomado de fuentes. Basta convertir la decena de impersonales localizados; conviene empezar por los tres puntos donde la ambigüedad tiene consecuencias: «Se adoptó como configuración base» (el autor), «Se empleó una única malla Q9» (el autor) y «se atribuye a la capa de contorno» (atribución que conviene declarar como interpretación del autor o como predicción de la literatura, con su localizador).\par
\rotulo{También en}\begin{itemize}[leftmargin=1.2em,itemsep=1pt,topsep=1pt]
\item {\small §3.2 Verificación de la solución, f. 64: \cita{se atribuye a la capa de contorno}}
\item {\small §3.3 Robustez y gestión del modelo, f. 65: \cita{se verificó la robustez de las operaciones de gestión del modelo}}
\item {\small §3.4 Validación con la viga de Timoshenko, f. 67: \cita{Se empleó una única malla Q9 estructurada de 56 x 8 elementos}}
\end{itemize}
\end{ficha}
\begin{ficha}{media}{\casilla\textbf{CLA-20}\quad\chip{media}{media}\quad Patrón de zona: 3.4, 3.7-3.8, Conclusiones y recomendaciones, Anexos D-E (ancla en 3.4) · folio 67 · PDF 82\hfill 180 min}
\rotulo{Dice}\cita{Folio 67 (PDF 82), Sección 3.4: «Se eligió una viga simplemente apoyada de longitud L = 14 m, peralte h = 1,20 m y base b = 0,40 m, sometida a una carga uniforme q = 5000 kgf/\allowbreak{}m, modelada en el sistema técnico (kgf, cm) con material E = 217 370 kgf/\allowbreak{}cm2 ({\mates ≈} 21,32 GPa) y ν = 0,20; ese módulo corresponde a la correlación módulo-resistencia E = 15 000 {\mates √}f'c (en kgf/\allowbreak{}cm2), la forma en unidades técnicas de la expresión que la norma ACI 318 propone para hormigón de peso normal, evaluada en f'c = 210 kgf/\allowbreak{}cm2; el análisis se realiza en régimen elástico lineal, sin considerar la fisuración del material»: un solo período de unas 95 palabras con dos punto y coma y dos incisos entre paréntesis (no entre rayas). Otros cuatro períodos largos verificados en su página: folio 80 /\allowbreak{} PDF 95 («la convergencia lenta y monótona por debajo de la referencia…»), folio 86 /\allowbreak{} PDF 101 («El trabajo regresa así a la doble brecha…»), folio 87 /\allowbreak{} PDF 102 («se construyó un núcleo numérico completo en NumPy y SciPy…», \textasciitilde{}95 palabras, con la enumeración de las etapas del canal entre rayas) y folio 128 /\allowbreak{} PDF 143 («el segundo término de la expresión de Timoshenko-Goodier…»). Recuento propio de rayas sobre la fuente LaTeX: 381 en total, {\mates ≈}328 en el cuerpo, repartidas de manera uniforme entre las diez unidades.}\par
\rotulo{Problema}En toda la zona (Cap. 3, Conclusiones y anexos redactados) el texto encadena períodos de 70 a 100 palabras que incrustan entre rayas dos y tres cláusulas subordinadas, de modo que el lector debe sostener a la vez la decisión de diseño, su justificación normativa, la remisión cruzada y la salvedad; el conteo automático de rayas de inciso en el cuerpo da 324 (3,4 por página, con picos de 8 a 10 en al menos doce páginas). Es precisamente el registro que un tribunal no especialista en elementos finitos no puede seguir en una lectura. El recuento bruto que hablaba de «una sola oración de 215 palabras» en el folio 87 corresponde al párrafo completo, no a un período: el período de apertura del balance del desarrollo son unas 95 palabras, cerradas en «identificadores de nodo no contiguos».\par
\rotulo{Corrección · reformular}Acotar la intervención a los cinco períodos localizados y aplicar un criterio de corte menos rígido: partir en punto y seguido todo período que pase de unas 60 palabras o que anide un inciso dentro de otra subordinada, sin tocar ninguna idea ni ninguna cifra. Los cinco puntos son folio 67 (3.4, datos de la viga), folio 73-75 (3.6.1, cobertura del canal, período que cruza dos folios con el pie de la Figura 3.5 en medio), folio 80 (3.7.1, separar el enunciado del bloqueo de su explicación), folio 87 (Conclusiones, partir el balance del desarrollo en núcleo numérico, aplicación, módulos educativos e intercambio de datos) y folio 128 (Anexo E, partir la nota de erratas en las tres ideas que encadena). Reescritura modelo para el folio 67, fiel al original: «Se eligió una viga simplemente apoyada de 14 m de luz, 1,20 m de peralte y 0,40 m de base, sometida a una carga uniforme de 5000 kgf/\allowbreak{}m. El modelo se planteó en unidades técnicas (kgf, cm), con E = 217 370 kgf/\allowbreak{}cm2 (unos 21,32 GPa) y coeficiente de Poisson 0,20. Ese módulo es el que la expresión que la norma ACI 318 propone para hormigón de peso normal, escrita en unidades técnicas como E = 15 000 {\mates √}f'c, arroja para una resistencia de 210 kgf/\allowbreak{}cm2. El análisis es elástico lineal y no considera la fisuración del hormigón.» No fijar un techo de 45 palabras ni proscribir el inciso de raya como norma de capítulo: la raya está repartida por igual en todo el documento ({\mates ≈}328 en el cuerpo, con más densidad en la unidad de la Sección 2.1 que en la del Capítulo 3), de modo que esa regla obligaría a reescribir la tesis entera para corregir una incomodidad de lectura que en ningún caso examinado induce a error. Si el autor quiere una pasada uniforme, que sea una lectura en voz alta de Capítulo 3, Conclusiones y anexos redactados, cortando solo donde el período obligue a releer.\par
\rotulo{También en}\begin{itemize}[leftmargin=1.2em,itemsep=1pt,topsep=1pt]
\item {\small §3.7.1 Interpretación de los resultados, f. 80: \cita{la convergencia lenta y monótona por debajo de la referencia es el comportamiento característico de un elemento que resulta excesivamente rígido cuando debe representar flexión mediante una interpolación bilineal incapaz de reproducir el modo de curvatura pura}}
\item {\small Conclusiones (apertura), f. 86: \cita{El trabajo regresa así a la doble brecha que lo originó}}
\item {\small Conclusiones (balance del desarrollo), f. 87: \cita{se construyó un núcleo numérico completo en NumPy y SciPy}}
\item {\small Anexo E (nota de erratas), f. 128: \cita{el segundo término de la expresión de Timoshenko-Goodier deja en los apoyos una distribución de σx autoequilibrada}}
\end{itemize}
{\footnotesize\color{gris}\rotulo{Verificación}ajustado: La cita es literal y está donde se dice (PDF 82 /\allowbreak{} folio 67), y el defecto existe: el período de apertura de la Sección 3.4 son unas 95-100 palabras con dos punto y coma. Los otros cuatro ejemplos también están en su página (folio 80 /\allowbreak{} PDF 95, folio 86 /\allowbreak{} PDF 101, folio 87 /\allowbreak{} PDF 102, folio 128 /\allowbreak{} PDF 143), de modo que el patrón se sostiene…\par}
\end{ficha}
\begin{ficha}{media}{\casilla\textbf{TIP-01}\quad\chip{media}{media}\quad Anexo E — documento de verificación en SAP2000 incrustado (PDF 143-154) · folio 128-139 · PDF 145\hfill 45 min}
\rotulo{Dice}\cita{«tensión en dirección x [kg/\allowbreak{}cm2]\allowbreak{} x=0 /\allowbreak{} 0 20 40 60 80 100 120 140» y «tensión en dirección x [kg/\allowbreak{}cm2]\allowbreak{} x= l/\allowbreak{}2=7m» (PDF 145, folio 130); «cortante en dirección x [kg/\allowbreak{}cm2]\allowbreak{} x=+7m /\allowbreak{} 15.0000 10.0000 5.0000 0.0000» (PDF 150, folio 135); «dezplazamiento de la línea neutra y=0» (PDF 153, folio 138). Las tres ubicaciones verificadas con ubicar.py.}\par
\rotulo{Problema}Las páginas PDF 143-154 (folios 128-139) son un documento externo de verificación en SAP2000 incrustado tal cual, y el lector no tiene forma de saberlo: entra sin un párrafo que diga qué es, quién lo elaboró y en qué condiciones se reproduce. De ahí que arrastre convenciones ajenas al resto de la tesis —punto decimal en los paneles de perfiles de tensión y de cortante (0.0000, 15.0000, 0.4) frente a la coma decimal del cuerpo y de las tablas del Anexo D, unidades en kg/\allowbreak{}cm2, rótulos en mayúsculas sin tilde, imágenes a baja resolución— y erratas de rotulado como «Dezplazamiento» y «linea» sin tilde. Sin la nota de procedencia, esas marcas se leen como descuidos de la tesis y no como el formato propio de una fuente externa. Nota de ubicación: los informes brutos atribuían estas páginas al Anexo D; el Anexo D ocupa PDF 136-142.\par
\rotulo{Corrección · reformular}Encabezar el Anexo E con un párrafo propio que declare qué es (verificación independiente del caso de la viga realizada con SAP2000), quién la elaboró y cuándo, que se reproduce tal como fue entregada y que por eso conserva sus propias convenciones —punto decimal, kg/\allowbreak{}cm2, rótulos en mayúsculas sin tilde—, distintas de las del cuerpo; añadir una línea de equivalencia con las unidades declaradas en la Nomenclatura. Si en cambio se decide integrarlo al estilo de la tesis, regenerar los paneles con coma decimal y a resolución de impresión, y corregir las erratas de los rótulos («Desplazamiento de la línea neutra y=0»). Verificar en la misma pasada que las llamadas del cuerpo a este material digan «Anexo E».\par
\rotulo{También en}\begin{itemize}[leftmargin=1.2em,itemsep=1pt,topsep=1pt]
\item {\small Anexo E, f. 135: \cita{cortante en dirección x [kg/\allowbreak{}cm2]\allowbreak{} x=+7m /\allowbreak{} 15.0000 10.0000 5.0000 0.0000}}
\item {\small Anexo E, f. 138: \cita{dezplazamiento de la linea neutra y=0}}
\end{itemize}
\end{ficha}
\begin{ficha}{baja}{\casilla\textbf{TIP-02}\quad\chip{baja}{baja}\quad Términos ingleses en Resumen, Capítulo 1 y Capítulo 3 · folio i · PDF 2\hfill 30 min}
\rotulo{Dice}\cita{«el bloqueo por cortante (shear-locking) del Q4 frente a la convergencia rápida del Q9» (folio i) en redonda dentro del paréntesis, frente a «hourglass» en cursiva en el cuerpo y en redonda dentro de la Tabla 1.1.}\par
\rotulo{Problema}Los términos ingleses conservados por ser los de la literatura (shear-locking, hourglass, fill-in, stretch, taper, B-bar, stack tecnológico) alternan cursiva y redonda sin criterio a lo largo de unas quince apariciones, incluida la diferencia entre el cuerpo y las celdas de una misma tabla. Es cosmético, pero un tribunal que lee el Resumen y la Tabla 1.1 seguidos ve el mismo término tratado de dos maneras.\par
\rotulo{Corrección · unificar}Fijar un criterio único y aplicarlo de corrido: cursiva en la primera aparición de cada término inglés, con la traducción o glosa al lado, y redonda en las siguientes, incluidas las celdas de tablas y los pies de figura. Recorrer con una búsqueda las apariciones de shear-locking, hourglass, fill-in, stretch, taper, B-bar y stack para dejarlas homogéneas.\par
\rotulo{También en}\begin{itemize}[leftmargin=1.2em,itemsep=1pt,topsep=1pt]
\item {\small Tabla 1.1, f. 16: \cita{Hourglass y bloqueo por cortante, vía modos propios}}
\item {\small §1.3-1.13, f. 24-33: \cita{stretch, taper, fill-in, B-bar con tratamiento desigual}}
\end{itemize}
\end{ficha}
\begin{ficha}{baja}{\casilla\textbf{TIP-03}\quad\chip{baja}{baja}\quad Variantes léxicas y de puntuación (1.1, 3.6, Limitaciones, Anexos B y F) · folio 15 · PDF 30\hfill 25 min}
\rotulo{Dice}\cita{Seis variantes sueltas frente a la forma dominante: «no consta» en el texto contra «No consta» en las celdas (folio 15); «Postproceso» una vez contra 51 «post-proceso» (folio 73); «competencias» una vez contra 42 «destreza(s)» (folio 13); «caso canónico» una vez contra 14 «ejemplo canónico» (folio 108); «esfuerzo» en la cabecera de unidades del CSV del Anexo F.2 contra «tensión» en todo el documento; y un único inciso con tres guiones ASCII sin espaciar (folio 92).}\par
\rotulo{Problema}El documento mantiene una terminología estable y estas seis apariciones se salen de ella, incluida una —«esfuerzo» en la cabecera del CSV— en la que la salida del software usa una palabra distinta de la que el documento emplea para la misma magnitud, y otra —el inciso con guiones ASCII— que rompe el espaciado de la raya en el único lugar del capítulo donde ocurre.\par
\rotulo{Corrección · unificar}Escribir «No consta» también al citarlo en el texto; «Post-proceso» en la columna del folio 73; «destrezas» en el folio 13; «ejemplo canónico» en el folio 108; y sustituir «---e incluso … ---» por la raya con el mismo espaciado que los demás incisos del capítulo. En el Anexo F.2, emitir y documentar la cabecera como «\# Unidades: longitud=mm, fuerza=N, tension=MPa», para que la salida del software use la misma palabra que el documento.\par
\rotulo{También en}\begin{itemize}[leftmargin=1.2em,itemsep=1pt,topsep=1pt]
\item {\small §3.6, f. 73: \cita{Postproceso Contornos de tensión y desplazamiento, deformada, isolíneas}}
\item {\small Anexo B.6, f. 108: \cita{atajo Ctrl+E para el caso canónico}}
\item {\small Limitaciones, f. 92: \cita{---e incluso para mallas de decenas de miles de grados de libertad---}}
\end{itemize}
\end{ficha}
\begin{ficha}{baja}{\casilla\textbf{TIP-04}\quad\chip{baja}{baja}\quad Capítulo 1 (Secciones 1.3 a 1.13) · folio 24 · PDF 39\hfill 30 min}
\rotulo{Dice}\cita{«la convención transpuesta, también frecuente en la literatura, exigiría J\textasciicircum{}-T» (folio 24): incisos entre rayas a razón de tres o más por página, en varios casos anidados con paréntesis y con comillas de énfasis.}\par
\rotulo{Problema}La densidad de incisos entre rayas fragmenta la línea de lectura en el capítulo más técnico del documento, justo donde el tribunal no especialista necesita frases que avancen en línea recta. Es un defecto de ritmo, no de contenido.\par
\rotulo{Corrección · corregir}Al reescribir las oraciones largas del capítulo, convertir la mayoría de los incisos en oraciones independientes o en paréntesis breves, y reservar la raya para una sola aclaración por párrafo. Se resuelve en la misma pasada de reescritura que pidan los hallazgos de claridad, sin trabajo aparte.\par
\rotulo{También en}\begin{itemize}[leftmargin=1.2em,itemsep=1pt,topsep=1pt]
\item {\small §1.9, f. 28: \cita{Su incorporación se plantea como línea de extensión futura}}
\item {\small §1.12, f. 33: \cita{el módulo M0 emplea la forma por vértices de Verdict}}
\end{itemize}
\end{ficha}
\begin{ficha}{baja}{\casilla\textbf{TIP-05}\quad\chip{baja}{baja}\quad §2.2.2 — Tabla del stack tecnológico y cuerpo del capítulo · folio 50 · PDF 65\hfill 20 min}
\rotulo{Dice}\cita{«Se evitó deliberadamente toda dependencia de bibliotecas multimedia pesadas» (folio 50): la fila de manim incrusta un inciso aclaratorio de doce palabras dentro de la celda, la leyenda de la tabla no dice que todas las dependencias listadas son libres, y el cuerpo arrastra identificadores de código en versalita (node\_\allowbreak{}index\_\allowbreak{}map, dof\_\allowbreak{}x, dof\_\allowbreak{}y, to\_\allowbreak{}dict, restore\_\allowbreak{}from\_\allowbreak{}dict, JACOBIAN\_\allowbreak{}MIN\_\allowbreak{}DETERMINANT, MMD\_\allowbreak{}AT\_\allowbreak{}PLUS\_\allowbreak{}A).}\par
\rotulo{Problema}Detalles de composición acumulados en la sección del stack: una celda que hace de nota al pie, una leyenda que no enuncia el requisito que la tabla justifica y ocho identificadores de código que ningún lector de la tesis va a teclear y que en versalita interrumpen la prosa.\par
\rotulo{Corrección · corregir}Pasar el inciso de manim a nota al pie de la tabla; ampliar la leyenda a «Componentes del stack tecnológico de EduFEM. Todas las bibliotecas empleadas son de código abierto y licencia permisiva, conforme al requisito de independencia tecnológica»; y trasladar los identificadores de código a notas al pie o al anexo del manual, dejando en el cuerpo su descripción en español.\par
\end{ficha}
\subsection*{Los demás patrones, por eje (están en su capítulo o en la sección 3)}
{\small\begin{longtable}{@{}L{1.4cm}L{2.6cm}L{1.2cm}L{10.05cm}@{}}\toprule
\sffamily\bfseries ID & \sffamily\bfseries Eje & \sffamily\bfseries Veces & \sffamily\bfseries Patrón\\\midrule\endhead\bottomrule\endfoot
CAL-43 & Calibración & 24 & Patrón en toda la zona: la familia garantiz- aparece 11 veces y confirm- 13 en el documento (recuento por script sobre el texto completo); una parte sustancial califica lo que sólo se sostiene por diseño, por la existencia de un…\\\addlinespace[2pt]
CAL-50 & Calibración & 4 & «Garantiza», «valida» y «confirma» se emplean indistintamente para lo que se sigue de una demostración o de la construcción y para lo que una prueba puntual, con una tolerancia concreta, solo hace verosímil; así el verbo pierde…\\\addlinespace[2pt]
CAL-57 & Calibración & 8 & El valor pedagógico de decisiones de diseño y de resultados numéricos se afirma como hecho, y el software se vuelve sujeto de verbos de enseñanza («ejercita», «explican», «profundiza en la comprensión»), cuando quien ejercita o…\\\addlinespace[2pt]
CAL-62 & Calibración & 4 & El impersonal y la agencia atribuida a los propios instrumentos («la tabla estableció», «la matriz estableció», «los instrumentos señalan») ocultan que fue el autor quien construyó las tablas y emitió el juicio, justo el límite…\\\addlinespace[2pt]
CAL-70 & Calibración & 4 & En cuatro salvedades, por lo demás correctas, el artículo definido o el posesivo presuponen que el efecto existe y que sólo falta medirlo. La propia tesis tiene ya la fórmula justa en el cierre de 1.2: «el efecto que la hipótesis…\\\addlinespace[2pt]
TRZ-07 & Trazabilidad & 9 & De los diecisiete criterios de aceptación fijados en 2.1.6, cinco tienen cifra medida pero ningún veredicto —tasa mínima de la tensión recuperada, residuo de equilibrio menor que 1e-8, Cook con Q9 y N = 8 dentro del 1,5 \%,…\\\addlinespace[2pt]
TRZ-14 & Trazabilidad & 18 & Dieciocho de las veintinueve ecuaciones numeradas del Capítulo 1 —D en tensión y deformación plana, funciones de forma de Q4 y Q9, Jacobiano y su determinante e inversa, cinemática, matriz B, integral de rigidez, ensamblaje,…\\\addlinespace[2pt]
TRZ-16 & Trazabilidad & 6 & El indicador con que se contrasta la cláusula (a) llama «módulo interactivo» a lo que el resto del documento llama «módulo educativo» —33 apariciones frente a 6—, de modo que el indicador y la cosa medida no comparten nombre; el…\\\addlinespace[2pt]
DEF-13 & Riesgo de defensa & 6 & Pregunta de alto rendimiento para quien quiera atacar la V\&V: «la cota de 1,4, ¿la fijó antes o después de medir 1,54?». Los seis umbrales de aceptación de la subsección (±0,5 en las tasas; 1,4 y 1,9; 3 \%; 1 \%; 1,5 \%; 10\textasciicircum{}-8 en el…\\\addlinespace[2pt]
DEF-15 & Riesgo de defensa & 3 & Comprobación natural de una de las banderas del trabajo: «¿podemos reproducir sus resultados? ¿qué versión exacta del programa produjo estas cifras?». La reproducibilidad se apoya en «la versión del software que acompaña a este…\\\addlinespace[2pt]
DEF-16 & Riesgo de defensa & 4 & Cuatro subsecciones del capítulo de implementación (2.2.2, 2.2.3, 2.2.4 y 2.2.9) se agotan en el cómo técnico y terminan sin la frase que ate la decisión al argumento de la tesis, de modo que el tribunal no especialista no puede…\\\addlinespace[2pt]
RED-02 & Redundancia & 5 & El índice del documento se recita en prosa cinco veces: «Estructura del documento», el cierre de «Metodología de la investigación», la apertura de cada capítulo y una segunda vez al abrir §2.1. Ninguna repetición añade contenido;…\\\addlinespace[2pt]
RED-03 & Redundancia & 4 & El inventario de limitaciones se escribe entero tres veces —Introducción, §3.7.4 y Conclusiones— y una cuarta vez, en parte, al abrir las Recomendaciones, donde el estado actual se recapitula en páginas consecutivas con oraciones…\\\addlinespace[2pt]
RED-04 & Redundancia & 5 & El diagnóstico de la caja negra —textos abstractos por un lado, paquetes que ocultan el procedimiento por otro, software didáctico como puente— se desarrolla desde cero cinco veces. Las dos del capítulo 1 (§1.1 y §1.2, a cinco…\\\addlinespace[2pt]
RED-05 & Redundancia & 13 & La salvedad de que el efecto sobre el estudiante no se mide y queda para una prueba de campo se enuncia con oración completa una docena de veces. Es la garantía de calibración del documento y su ausencia sería un defecto mayor…\\\addlinespace[2pt]
RED-06 & Redundancia & 3 & Tres exposiciones del mismo respaldo bibliográfico. La diferencia no es solo de economía: únicamente §1.2 acota la evidencia (cualitativa, dependiente de la visualización y del apoyo instruccional, atribuida en parte a la autoría…\\\addlinespace[2pt]
RED-09 & Redundancia & 4 & El mismo hallazgo de Pérez-Santiago y Campos se enuncia entero cuatro veces, dos de ellas (§1.1 y §1.2) casi palabra por palabra y con la misma referencia a la misma página. Solo la de §1.2 cumple función argumental: es la que…\\\addlinespace[2pt]
RED-11 & Redundancia & 5 & La precaución de integrar las normas de error con un punto de Gauss adicional por dirección se enuncia cinco veces con su paréntesis y su justificación completa. Es un detalle de implementación que basta establecer una vez;…\\\addlinespace[2pt]
RED-13 & Redundancia & 4 & La separación entre identidad del nodo e índice de grado de libertad, que es una decisión de diseño y no un resultado, se explica cuatro veces: donde se decide, donde se prueba y dos veces al concluir.\\\addlinespace[2pt]
RED-14 & Redundancia & 4 & El estado actual del solucionador —ensamblaje por tripletes COO, conversión a CSR, factorización LU dispersa y ordenamiento de mínimo grado— se describe cuatro veces fuera de la sección que lo decide, incluida la recomendación,…\\\addlinespace[2pt]
RED-20 & Redundancia & 81 & El verbo «encarna» y los conectores «de forma» /\allowbreak{} «de manera» se repiten hasta convertirse en tic de redacción. No hay error de lengua, pero la monotonía léxica se nota en un texto de noventa y cinco folios y delata pasajes…\\\addlinespace[2pt]
ACC-05 & Accesibilidad & 3 & Tres términos metodológicos que el tribunal necesita para juzgar el trabajo —validez de contenido, tabla de especificaciones y variable interviniente— se usan en el Resumen y en la Introducción y solo se definen en la sección…\\\addlinespace[2pt]
ACC-07 & Accesibilidad & 33 & Patrón transversal: el término técnico se usa en su primera aparición sin una aposición de media línea que diga qué es y por qué importa, y su definición —cuando existe— llega decenas de folios después o solo en la nomenclatura.…\\\addlinespace[2pt]
ACC-09 & Accesibilidad & 20 & Patrón: una veintena de términos de programación y estructuras de datos —factorización LU dispersa, ordenamiento de mínimo grado, COLAMD, Cuthill-McKee inverso, fill-in, COO, CSR, índice espacial, firma topológica, idempotente,…\\\addlinespace[2pt]
ACC-11 & Accesibilidad & 4 & «Los cuatro principios de aprendizaje» se invocan cuatro veces entre el Resumen y las Conclusiones sin enumerarse nunca en esos capítulos, de modo que el lector que solo lea el Resumen y las Conclusiones —lectura habitual en un…\\\addlinespace[2pt]
CON-04 & Consistencia & 5 & Cinco colisiones o descuidos de notación frente a la Nomenclatura: Q designa la compacidad mientras Q4 y Q9 designan la familia de elementos; E designa el módulo de elasticidad en 1.4 y la matriz de extrapolación en 1.11,…\\\addlinespace[2pt]
CON-06 & Consistencia & 22 & El destinatario se llama «estudiante» en el problema, la hipótesis y las variables (62 apariciones), «alumno» en los fundamentos pedagógicos y en parte del Capítulo 2 (22 apariciones, seis de ellas en el folio 19) y «usuario» en…\\\addlinespace[2pt]
CON-07 & Consistencia & 4 & La Sección 1.13 declara una convención para «en el resto del documento» y el documento la incumple después en el término más delicado del trabajo: la misma raíz designa además la validez de contenido (folio 7), el validador de…\\\addlinespace[2pt]
CON-08 & Consistencia & 22 & Las tres unidades de análisis se nombran «casos de estudio» en el diseño metodológico (5 apariciones), «casos de referencia» en el Capítulo 3 y en los anexos (17), «problemas de referencia» en 1.13, 3.7.1 y los Aportes (3) y…\\\addlinespace[2pt]
CON-09 & Consistencia & 6 & EduFEM se nombra de seis maneras -herramienta (57), software (106), programa (51), aplicación (19), artefacto (15) y recurso (20)-, y las seis conviven dentro del mismo capítulo, lo que diluye el sujeto del que se predica cada…\\\addlinespace[2pt]
CON-12 & Consistencia & 6 & La sección de Cook separa las coordenadas con coma mientras el resto del capítulo usa punto y coma, de modo que en un texto con coma decimal «(48, 52)» se lee como un único número -cuarenta y ocho coma cincuenta y dos- en vez de…\\\addlinespace[2pt]
CON-17 & Consistencia & 4 & El concepto vertebrador del trabajo, «canal de cálculo» -27 apariciones y parte del objetivo general-, alterna con «canal del MEF» (folio 17), «flujo del Método de Elementos Finitos» (folio 109) y «proceso de cálculo» (folios 2 y…\\\addlinespace[2pt]
CON-18 & Consistencia & 13 & El mismo componente se llama «validador de salud» en el cuerpo (5 apariciones, folios 33, 44, 56, 73 y 78, entre ellas la celda de evidencia de la matriz de principios) y «comprobador de salud» en el Anexo B que lo documenta (8…\\\addlinespace[2pt]
CON-19 & Consistencia & 7 & El Anexo E -PDF 143-154, folios 128-139; no confundir con el Anexo D, PDF 136-142- rompe a la vez la terminología, la notación y la tipografía del documento: dice «esfuerzo» donde el resto dice «tensión» (7 de las 8 apariciones…\\\addlinespace[2pt]
APF-04 & Aparato formal & 26 & La lista de símbolos y siglas promete reunir la notación del documento y no cumple en ninguno de los dos sentidos. Por defecto omite más de veinte símbolos que aparecen en ecuaciones numeradas del Capítulo 1 y en las tablas de…\\\addlinespace[2pt]
APF-07 & Aparato formal & 12 & El documento remite a sí mismo sin número de sección, tabla ni página justo en los puntos donde el tribunal querría ir a comprobar: el objetivo específico que una sección afirma cumplir, la delimitación del alcance, el dato de…\\\addlinespace[2pt]
APF-09 & Aparato formal & 10 & Diez leyendas de figura y tabla son meros rótulos que no se pueden leer sin el cuerpo: no dicen qué representan las cajas ni en qué sentido corren las flechas, de qué módulo es la captura ni sobre qué modelo, qué magnitud va en…\\\addlinespace[2pt]
APF-14 & Aparato formal & 12 & La columna de evidencia mezcla dos formas de remisión -el signo de sección abreviado y la palabra completa- en doce celdas frente a ocho, mientras la Tabla 3.6 de la misma sección y toda la prosa del documento usan siempre…\\\addlinespace[2pt]
APF-15 & Aparato formal & 8 & Las capturas de los módulos M1 a M6 se presentan en bloque y quedan luego solas con su leyenda: ninguna se referencia individualmente ni desde la Tabla 2.5 ni desde la Subsección 2.2.7, que es donde el lector del cuerpo esperaría…\\\addlinespace[2pt]
APF-16 & Aparato formal & 5 & El anexo se presenta como memoria paso a paso trazable a la formulación teórica, pero el puente entre la sustitución numérica y la definición simbólica solo existe en dos de siete eslabones. El lector que quiera comprobar de…\\\addlinespace[2pt]
APF-18 & Aparato formal & 3 & Tres páginas quedan a medio llenar por colocación de flotantes o por cierre de sección. No hay error de contenido —son cierres legítimos—, pero en un documento de 168 páginas destinado a impresión desaprovechan espacio y, en el…\\\addlinespace[2pt]
REF-01 & Referencias & 6 & La cláusula que acota el alcance de la hipótesis —la afirmación más consultada del documento, y la que separa lo comprobado de lo pendiente— remite a «la literatura» en cinco puntos de máximo peso (Resumen, Introducción, 2.1,…\\\addlinespace[2pt]
REF-02 & Referencias & 22 & De las 174 llamadas de cita del cuerpo, 37 van sin localizador de página; 22 de ellas, repartidas en 12 puntos del texto, sostienen una definición o una atribución concreta y por tanto lo exigen según la regla que la propia tesis…\\\addlinespace[2pt]
REF-04 & Referencias & 5 & Cinco claves del .bib llevan un año distinto del que imprime la entrada. La clave no llega al lector, de modo que el PDF no muestra la discrepancia, pero sí señala el riesgo que sí es verificable: que los localizadores de página…\\\addlinespace[2pt]
\end{longtable}}
\clearpage
\section{Preguntas de defensa}
Un tribunal simulado de tres miembros de ingeniería civil —M1, docente de metodología; M2, estructuralista usuario de SAP2000, no especialista en MEF ni en programación; M3, docente de métodos numéricos con interés pedagógico— formuló las preguntas que haría tras leer el documento, y para cada una se comprobó dónde responde hoy el texto. «Qué falta» distingue lo que conviene escribir en la tesis de lo que basta con preparar de palabra. Los riesgos que exigen cambiar el texto están además en la tabla de hallazgos con prefijo DEF.
\needspace{5\baselineskip}\noindent\casilla\textbf{1. Su hipótesis afirma que la exposición del canal de cálculo mejorará el criterio de los estudiantes, y su objetivo general termina diciendo «para mejorar el criterio de los estudiantes de ingeniería civil». ¿A cuántos estudiantes midió?}\par
{\small\color{gris}\sffamily M1 metodologia\quad\chip{alta}{parcial}\quad riesgo critica}\par
{\small\rotulo{Dónde responde hoy}Introducción, Hipótesis y preguntas de investigación, folios 6-7; 2.1.1, folio 38 («la unidad de análisis de esta etapa es el artefacto»); 3.8, folio 84; Conclusiones, folio 89\par}
{\small\rotulo{Qué falta}La respuesta está construida y es sólida, pero el objetivo general (folio 5) y el OE5 (folio 6) siguen prometiendo la mejora y el «potencial para formar criterio»: reformularlos como propiedades del artefacto (DEF-02) y suavizar el «queda comprobada» del Resumen, 3.8 y Conclusiones (DEF-04). De palabra, tres frases: qué se comprobó, qué no y dónde queda diseñado lo que falta.\par}
\vspace{5pt}
\needspace{5\baselineskip}\noindent\casilla\textbf{2. Usted parte de que los estudiantes de esta carrera tienen «bajo criterio» para interpretar la respuesta estructural. ¿De dónde sale ese diagnóstico? ¿Lo midió aquí?}\par
{\small\color{gris}\sffamily M1 metodologia\quad\chip{critica}{sin respuesta}\quad riesgo critica}\par
{\small\rotulo{Dónde responde hoy}Introducción, Planteamiento del problema, folio 2: se apoya en el consenso de expertos y en la encuesta de Pérez-Santiago y Campos, ambos externos; no hay ningún dato de la UATF en todo el documento\par}
{\small\rotulo{Qué falta}El párrafo de DEF-01: declarar que el «bajo criterio» es una premisa tomada de la literatura y del diagnóstico docente, no un dato medido, y que la preprueba de la prueba de campo lo establecería en la población concreta.\par}
\vspace{5pt}
\needspace{5\baselineskip}\noindent\casilla\textbf{3. Usted cuenta que encontró un error en la matriz de extrapolación que su batería de convergencia no detectó, y luego concluye que el motor está «libre de errores de orden». ¿En qué quedamos?}\par
{\small\color{gris}\sffamily M3 numerico-pedagogico\quad\chip{alta}{parcial}\quad riesgo critica}\par
{\small\rotulo{Dónde responde hoy}3.7.1, folio 79 («no lo delataron las normas de desplazamiento sino la inspección de la memoria de cálculo»); Conclusiones, folio 88\par}
{\small\rotulo{Qué falta}Reformular el cierre de las Conclusiones (DEF-03). De palabra: la verificación prueba la ausencia de una clase de errores, y el trabajo demuestra madurez al mostrar el límite de su propia batería y la prueba que se añadió para cubrirlo.\par}
\vspace{5pt}
\needspace{5\baselineskip}\noindent\casilla\textbf{4. ¿Cuál es exactamente su variable dependiente y cómo puede medirla en un programa de computadora si la definió como una capacidad del estudiante?}\par
{\small\color{gris}\sffamily M1 metodologia\quad\chip{alta}{parcial}\quad riesgo alta}\par
{\small\rotulo{Dónde responde hoy}2.1.2 Variables e indicadores y Tabla 2.1, folios 39-40; medición de los datos, folio 75\par}
{\small\rotulo{Qué falta}La tabla lo resuelve, pero el texto no lo defiende: añadir el párrafo de DEF-11, con la distinción entre validez de contenido como condición necesaria y criterio logrado como magnitud que solo se mide en el estudiante.\par}
\vspace{5pt}
\needspace{5\baselineskip}\noindent\casilla\textbf{5. Su tabla de especificaciones da dieciocho de dieciocho. Si usted elige el universo, elige el instrumento y pone la marca, ¿qué prueba ese resultado?}\par
{\small\color{gris}\sffamily M3 numerico-pedagogico\quad\chip{alta}{parcial}\quad riesgo alta}\par
{\small\rotulo{Dónde responde hoy}3.7.2, folio 77 («la cobertura no es uniforme y conviene decirlo»); 2.1.6, folio 46, donde se fija el criterio «al menos un instrumento identificable»\par}
{\small\rotulo{Qué falta}Decir en el texto qué prueba y qué no prueba el 18/\allowbreak{}18 (DEF-20): es cobertura y trazabilidad, no suficiencia ni profundidad. De palabra, subrayar que el universo no lo fijó el autor sino el consenso publicado, y ahí está la parte que no es autoevaluación.\par}
\vspace{5pt}
\needspace{5\baselineskip}\noindent\casilla\textbf{6. ¿Quién validó el contenido de su software? ¿Hubo un panel de expertos, jueces, docentes de la materia?}\par
{\small\color{gris}\sffamily M1 metodologia\quad\chip{alta}{parcial}\quad riesgo alta}\par
{\small\rotulo{Dónde responde hoy}2.1.5, folios 45-46 («el criterio externo es el consenso publicado de expertos, que hace las veces de juicio de expertos sin panel propio»; «su limitación, que se declara, es que las construye el autor»); Limitaciones, folio 91\par}
{\small\rotulo{Qué falta}Está declarado pero no procedimentado: falta la regla de marcado de las celdas (DEF-15). De palabra, ofrecer un ejemplo de celda que quedaría vacía con esa regla y decir que el consenso publicado aporta el criterio externo que el autor no podía darse a sí mismo.\par}
\vspace{5pt}
\needspace{5\baselineskip}\noindent\casilla\textbf{7. Sus 67 expertos son mexicanos y de ingeniería mecánica, y el propio estudio dice que no son generalizables. ¿Por qué valen para ingeniería civil en Bolivia?}\par
{\small\color{gris}\sffamily M1 metodologia\quad\chip{alta}{parcial}\quad riesgo alta}\par
{\small\rotulo{Dónde responde hoy}2.1.5, folio 45 («se acota porque los ítems que se toman son contenido universal del método»); 2.1.4, folio 42\par}
{\small\rotulo{Qué falta}Sostener la universalidad con las dos vías de DEF-05 —los ítems tomados son operaciones del método respaldadas por los textos clásicos, y los ítems dependientes de contexto quedan fuera del universo— y llevar la salvedad al Resumen, a la Introducción y a las Limitaciones. Reconocer que lo intrasladable es la ponderación, y que la tabla no pondera.\par}
\vspace{5pt}
\needspace{5\baselineskip}\noindent\casilla\textbf{8. Su criterio de aceptación exige menos del 1 \% de error en la tensión normal, pero su propia tabla muestra 4,56 \% en la componente transversal frente a SAP2000 y 2,89 \% en el cortante. ¿No fijó el criterio sobre lo que sabía que iba a pasar?}\par
{\small\color{gris}\sffamily M3 numerico-pedagogico\quad\chip{alta}{parcial}\quad riesgo alta}\par
{\small\rotulo{Dónde responde hoy}2.1.6, folio 46, donde el criterio se fija solo sobre la tensión normal y la flecha; 3.6, folios 71-72, que reporta las componentes secundarias; 3.7.1, folio 80, que explica la diferencia de escala\par}
{\small\rotulo{Qué falta}Que el propio criterio diga por qué excluye las componentes secundarias (DEF-14). El argumento de escala existe, pero aparece después del dato y no antes, que es donde el tribunal lo necesita.\par}
\vspace{5pt}
\needspace{5\baselineskip}\noindent\casilla\textbf{9. La cota de 1,4 para la tasa del campo de tensiones, ¿la fijó antes o después de medir 1,54? ¿Y de dónde salen los demás umbrales?}\par
{\small\color{gris}\sffamily M3 numerico-pedagogico\quad\chip{alta}{parcial}\quad riesgo alta}\par
{\small\rotulo{Dónde responde hoy}2.1.6, folio 46; 3.2, folio 63\par}
{\small\rotulo{Qué falta}Explicar la procedencia de los seis umbrales o degradar la cota a observación sin criterio de aceptación (DEF-13). De palabra, tener a mano el registro en el guion de V\&V con su historial de versiones.\par}
\vspace{5pt}
\needspace{5\baselineskip}\noindent\casilla\textbf{10. Su tabla de síntesis dice «error máx.». ¿Máximo sobre qué? ¿Barrió todo el dominio o solo esos tres puntos?}\par
{\small\color{gris}\sffamily M3 numerico-pedagogico\quad\chip{critica}{sin respuesta}\quad riesgo alta}\par
{\small\rotulo{Dónde responde hoy}3.4, folio 68, tabla de puntos A, B y C; tabla de síntesis, folio 72\par}
{\small\rotulo{Qué falta}Declarar el criterio de elección de los tres puntos, reportar el error máximo sobre todos los nodos o decir que no se calculó, y cambiar el rótulo de la tabla de síntesis (DEF-18).\par}
\vspace{5pt}
\needspace{5\baselineskip}\noindent\casilla\textbf{11. ¿De dónde sale la partición del canal de cálculo en nueve etapas contra la que mide la cobertura de sus módulos?}\par
{\small\color{gris}\sffamily M3 numerico-pedagogico\quad\chip{critica}{sin respuesta}\quad riesgo alta}\par
{\small\rotulo{Dónde responde hoy}Cap. 3, cobertura del canal de cálculo, folio 75, y su eco en folio 78\par}
{\small\rotulo{Qué falta}La frase de DEF-19: que la partición reproduce la secuencia con que la literatura de referencia presenta la formulación isoparamétrica, con cita y localizador, y que con ella se fijó el criterio en 2.1.6 antes de construir los módulos.\par}
\vspace{5pt}
\needspace{5\baselineskip}\noindent\casilla\textbf{12. Yo calculo pórticos y losas. Si su programa no hace pórticos, ni losas, ni tres dimensiones, ¿para qué me sirve en ingeniería civil?}\par
{\small\color{gris}\sffamily M2 estructuralista\quad\chip{alta}{parcial}\quad riesgo alta}\par
{\small\rotulo{Dónde responde hoy}Introducción, Alcance y limitaciones, folios 8-9, con la justificación pedagógica del dominio plano; 3.6, folio 68, viga de hormigón; Limitaciones observadas, folio 85\par}
{\small\rotulo{Qué falta}La frase en positivo de DEF-07: vigas pared, ménsulas, regiones D, muros y, en deformación plana, presas, muros de contención y túneles. Y recordar que el propósito declarado es formativo, no productivo.\par}
\vspace{5pt}
\needspace{5\baselineskip}\noindent\casilla\textbf{13. ¿Por qué solo Q4 y Q9? El programa de 1998 que usted cita como antecedente ya tenía triángulos, Q4, Q8 y Q9.}\par
{\small\color{gris}\sffamily M2 estructuralista\quad\chip{alta}{parcial}\quad riesgo alta}\par
{\small\rotulo{Dónde responde hoy}Introducción, Alcance, folio 8; 1.1, folio 14, biblioteca de ED-Elas2D; Recomendaciones, folio 93\par}
{\small\rotulo{Qué falta}El párrafo de DEF-08: la pareja mínima contrastante, por qué no el Q8 serendípito y por qué los triángulos pertenecen al mallado automático que esta versión no incorpora.\par}
\vspace{5pt}
\needspace{5\baselineskip}\noindent\casilla\textbf{14. Usted afirma que «no existe» un entorno con estas características. ¿Qué buscó, dónde y con qué criterio?}\par
{\small\color{gris}\sffamily M3 numerico-pedagogico\quad\chip{alta}{parcial}\quad riesgo alta}\par
{\small\rotulo{Dónde responde hoy}1.1, folios 15 y 17; la nota de la Tabla 1.1 acota «no consta» como límite de la revisión, pero el cierre generaliza a «no existe»\par}
{\small\rotulo{Qué falta}El párrafo de método de la revisión y el acotamiento del cierre (DEF-09). De palabra, si aparece un contraejemplo, conceder de inmediato y reconducir al conjunto simultáneo de atributos, que es donde el vacío se sostiene.\par}
\vspace{5pt}
\needspace{5\baselineskip}\noindent\casilla\textbf{15. ¿Qué impide que el estudiante use EduFEM como otra caja negra, más lenta que SAP2000?}\par
{\small\color{gris}\sffamily M3 numerico-pedagogico\quad\chip{critica}{sin respuesta}\quad riesgo alta}\par
{\small\rotulo{Dónde responde hoy}En ninguna parte de forma directa; lo más próximo es la consigna del tratamiento de la prueba de campo, folio 94\par}
{\small\rotulo{Qué falta}El párrafo de limitaciones de DEF-22: la exposición es una posibilidad, no una obligación; el recorrido lo fija la consigna docente, y la memoria y los archivos de proyecto permiten verificar si se hizo.\par}
\vspace{5pt}
\needspace{5\baselineskip}\noindent\casilla\textbf{16. ¿Nadie más que usted ha usado el programa? ¿Lo probó algún estudiante o algún docente? ¿Quién revisó su código?}\par
{\small\color{gris}\sffamily M1 metodologia\quad\chip{critica}{sin respuesta}\quad riesgo alta}\par
{\small\rotulo{Dónde responde hoy}Aportes, folio 90, y Limitaciones, folio 91, que declaran que las tablas las construye el autor pero no dicen nada sobre uso o revisión por terceros\par}
{\small\rotulo{Qué falta}Las dos frases de limitaciones de DEF-21. Si hay tiempo antes de la defensa, hacer que dos o tres compañeros instalen el programa y resuelvan el ejemplo canónico, y reportarlo en una línea: es la evidencia más barata disponible.\par}
\vspace{5pt}
\needspace{5\baselineskip}\noindent\casilla\textbf{17. ¿Qué parte de este trabajo se hizo con asistencia de inteligencia artificial y cómo lo declara?}\par
{\small\color{gris}\sffamily M1 metodologia\quad\chip{critica}{sin respuesta}\quad riesgo alta}\par
{\small\rotulo{Dónde responde hoy}En ninguna parte del documento\par}
{\small\rotulo{Qué falta}La declaración de herramientas de DEF-06, en los términos que el reglamento de la carrera admita. De palabra: la responsabilidad intelectual del diseño, de la formulación y de la verificación es del autor, y la corrección del código no descansa en quién lo escribió sino en la batería de V\&V que cualquiera puede reejecutar.\par}
\vspace{5pt}
\needspace{5\baselineskip}\noindent\casilla\textbf{18. ¿Podemos reproducir sus resultados? ¿Qué versión exacta del programa produjo estas cifras?}\par
{\small\color{gris}\sffamily M3 numerico-pedagogico\quad\chip{alta}{parcial}\quad riesgo alta}\par
{\small\rotulo{Dónde responde hoy}2.1.5, folio 46, determinismo del motor y guiones reejecutables; nota al pie del folio 4 con la dirección del repositorio; Anexo A, manual de instalación\par}
{\small\rotulo{Qué falta}Fijar versión, etiqueta y confirmación del repositorio (DEF-16). Para la defensa: llevar el instalador y una corrida en vivo de uno de los guiones de V\&V, que responde la pregunta mejor que cualquier párrafo.\par}
\vspace{5pt}
\needspace{5\baselineskip}\noindent\casilla\textbf{19. ¿Cómo se usaría esto en la cátedra? ¿Tiene una guía o una secuencia de prácticas para el docente?}\par
{\small\color{gris}\sffamily M3 numerico-pedagogico\quad\chip{alta}{parcial}\quad riesgo alta}\par
{\small\rotulo{Dónde responde hoy}Recomendaciones, folio 94, dentro del tratamiento de la prueba de campo; Anexo B, manual de uso, folios 121 y siguientes\par}
{\small\rotulo{Qué falta}El ítem de guía de uso en la cátedra con tres prácticas (DEF-25). Hoy el manual enseña a operar el programa, no a enseñar con él.\par}
\vspace{5pt}
\needspace{5\baselineskip}\noindent\casilla\textbf{20. ¿En qué asignatura, de qué semestre y con cuántos estudiantes se introduce el MEF en esta carrera? Su población dice «las asignaturas» y más adelante «la asignatura».}\par
{\small\color{gris}\sffamily M1 metodologia\quad\chip{critica}{sin respuesta}\quad riesgo alta}\par
{\small\rotulo{Dónde responde hoy}2.1.4 Población, universo de contenido y casos de estudio, folio 41; Recomendaciones, prueba de campo, folio 94\par}
{\small\rotulo{Qué falta}Nombrar la asignatura con sigla, semestre y matrícula aproximada, en singular y de forma idéntica en 2.1.4 y en las recomendaciones (DEF-12).\par}
\vspace{5pt}
\needspace{5\baselineskip}\noindent\casilla\textbf{21. Su prueba de campo es de preprueba y posprueba con un solo grupo, sin hipótesis estadística ni criterio de éxito fijado de antemano. ¿Cómo atribuiría la ganancia a la exposición del canal, si es el mismo reparo que usted le hace a Lee?}\par
{\small\color{gris}\sffamily M1 metodologia\quad\chip{alta}{parcial}\quad riesgo alta}\par
{\small\rotulo{Dónde responde hoy}Recomendaciones, Prueba de campo, folio 94; 1.2, folio 18, donde se le reprocha a Lee la falta de grupo de control\par}
{\small\rotulo{Qué falta}Declarar preferente el diseño con grupo de comparación y reportar la limitación si solo hay un grupo (DEF-23), y completar el protocolo con hipótesis estadística, potencia, validación del instrumento, criterio de éxito a priori, ética y plazo (DEF-24).\par}
\vspace{5pt}
\needspace{5\baselineskip}\noindent\casilla\textbf{22. ¿Por qué debo creerle a su motor con solo tres casos de prueba?}\par
{\small\color{gris}\sffamily M3 numerico-pedagogico\quad\chip{verde}{respondida}\quad riesgo media}\par
{\small\rotulo{Dónde responde hoy}2.1.4, folios 42-43, selección intencional por valor probatorio y cobertura de todos los términos del modelo matemático; 3.2 a 3.6, folios 62-74; Limitaciones, folios 92-93\par}
{\small\rotulo{Qué falta}Nada esencial. Subrayar de palabra que el método de soluciones manufacturadas no es «un caso» sino veinte corridas en cuatro configuraciones con solución exacta, y que la verificación de código se juzga por cobertura de términos, no por número de ejemplos.\par}
\vspace{5pt}
\needspace{5\baselineskip}\noindent\casilla\textbf{23. Si ED-Elas2D ya hacía análisis de sólidos bidimensionales, en español y con pre, proceso y post, ¿qué aporta su trabajo?}\par
{\small\color{gris}\sffamily M2 estructuralista\quad\chip{verde}{respondida}\quad riesgo media}\par
{\small\rotulo{Dónde responde hoy}1.1, folios 14-15, con la diferencia declarada «no reside en la lengua sino en la concepción y la tecnología»; Tabla 1.1, folios 15-16; Aportes, folio 90\par}
{\small\rotulo{Qué falta}Nada. Llevar preparada la versión de una frase: exposición interactiva de cada matriz sobre el elemento del alumno, memoria de cálculo trazable y reproducible a mano, lectura de la respuesta con crudo y suavizado, reacciones y Mohr, y V\&V publicada con la herramienta.\par}
\vspace{5pt}
\needspace{5\baselineskip}\noindent\casilla\textbf{24. Usted dice «validación», pero no contrastó contra ningún ensayo físico. ¿Es correcto llamarlo validación?}\par
{\small\color{gris}\sffamily M3 numerico-pedagogico\quad\chip{alta}{parcial}\quad riesgo media}\par
{\small\rotulo{Dónde responde hoy}Cap. 3, entradilla, folio 61; Conclusiones, folio 89 («la validación en sentido estricto, contra mediciones físicas, queda fuera del alcance»)\par}
{\small\rotulo{Qué falta}Llevar la acotación al punto de uso, en la entradilla del capítulo (DEF-30). De palabra: verificación es resolver bien las ecuaciones, validación es resolver las ecuaciones correctas, y aquí la referencia es una solución analítica de la teoría de la elasticidad y un programa comercial, no un ensayo.\par}
\vspace{5pt}
\needspace{5\baselineskip}\noindent\casilla\textbf{25. SAP2000 resuelve ese modelo con cáscaras Mindlin-Reissner y usted con tensión plana. ¿Es legítimo comparar y llamar «error» a la diferencia? ¿Quién construyó ese modelo, con qué versión y con qué malla?}\par
{\small\color{gris}\sffamily M2 estructuralista\quad\chip{alta}{parcial}\quad riesgo media}\par
{\small\rotulo{Dónde responde hoy}3.6, folios 62-63 («pequeñas discrepancias entre ambos modelos son esperables y no deben interpretarse como mayor o menor exactitud»); 3.7.1, folio 80; Anexo E, folios 128-139\par}
{\small\rotulo{Qué falta}Comprobar que el Anexo E declare versión del programa, tipo y tamaño de elemento de cáscara, condiciones de apoyo y puntos de extracción; si falta alguno, añadirlo. Llevar el archivo del modelo a la defensa y usar el dato que juega en contra: en la componente cortante el modelo de cáscara queda más cerca del analítico que EduFEM.\par}
\vspace{5pt}
\needspace{5\baselineskip}\noindent\casilla\textbf{26. El campo de tensiones del Q4 no alcanza el orden cuadrático y usted lo atribuye a la capa de contorno. ¿Lo comprobó?}\par
{\small\color{gris}\sffamily M3 numerico-pedagogico\quad\chip{critica}{sin respuesta}\quad riesgo media}\par
{\small\rotulo{Dónde responde hoy}3.2, folio 65\par}
{\small\rotulo{Qué falta}Marcar el estatus de la explicación como plausible y no contrastada, y proponer la comprobación como trabajo futuro (DEF-17).\par}
\vspace{5pt}
\needspace{5\baselineskip}\noindent\casilla\textbf{27. El estudio que usted invoca como evidencia cuantitativa atribuye la mejora principalmente a que los alumnos formulan sus propios problemas, cosa que su software no hace. ¿Por qué deriva de ahí el diseño de sus módulos?}\par
{\small\color{gris}\sffamily M3 numerico-pedagogico\quad\chip{critica}{sin respuesta}\quad riesgo media}\par
{\small\rotulo{Dónde responde hoy}1.2 Fundamentos pedagógicos, tercer principio, folio 19\par}
{\small\rotulo{Qué falta}La frase de acotación de DEF-10: lo que se toma del estudio es la construcción por etapas de las herramientas de cálculo, y el punto de contacto con aquella condición es que la memoria se genera sobre el modelo que el propio alumno construye.\par}
\vspace{5pt}
\needspace{5\baselineskip}\noindent\casilla\textbf{28. ¿Por qué Python y no una hoja de cálculo, MATLAB u Octave, que ya tienen todo esto resuelto?}\par
{\small\color{gris}\sffamily M2 estructuralista\quad\chip{alta}{parcial}\quad riesgo media}\par
{\small\rotulo{Dónde responde hoy}Introducción, Justificación, plano metodológico, folio 4; Tabla 1.1, folio 15, que marca «Requiere MATLAB» como requisito del antecedente de Bishay; 2.2, pila tecnológica\par}
{\small\rotulo{Qué falta}Una frase explícita en la Justificación: gratuidad y ausencia de licencia para el alumno, ecosistema numérico y de interfaz en un solo lenguaje, empaquetado en un instalador autónomo sin requerir el intérprete; y la hoja de cálculo como límite, aunque sirva —y el trabajo lo aproveche— para reproducir a mano una etapa de la memoria.\par}
\vspace{5pt}
\needspace{5\baselineskip}\noindent\casilla\textbf{29. Si el Q4 bloquea, ¿qué impide que un alumno entregue un resultado equivocado sin enterarse? ¿El programa le avisa?}\par
{\small\color{gris}\sffamily M2 estructuralista\quad\chip{alta}{parcial}\quad riesgo media}\par
{\small\rotulo{Dónde responde hoy}Limitaciones observadas, folio 85, y Limitaciones, folio 91 (el bloqueo «exige al usuario refinar la malla o recurrir al Q9»); 3.5, membrana de Cook, folios 73-74\par}
{\small\rotulo{Qué falta}Decir qué hace el software al respecto —la conversión Q4-Q9 en un paso y la comparación sobre el mismo modelo son la respuesta de diseño— y, si no existe un aviso explícito, proponerlo entre las recomendaciones sobre el software.\par}
\vspace{5pt}
\needspace{5\baselineskip}\noindent\casilla\textbf{30. ¿Cómo eligió los cuatro principios de aprendizaje? ¿Por qué esos y no un marco completo de diseño instruccional?}\par
{\small\color{gris}\sffamily M1 metodologia\quad\chip{alta}{parcial}\quad riesgo media}\par
{\small\rotulo{Dónde responde hoy}1.2 Fundamentos pedagógicos, folios 18-20, donde cada principio se apoya en la literatura de enseñanza del MEF y en una fuente general; folio 20, que declara que no constituyen evidencia de eficacia de EduFEM\par}
{\small\rotulo{Qué falta}Una frase sobre el criterio de selección: son los principios que la literatura específica de enseñanza del MEF invoca para justificar el software didáctico, cada uno con respaldo en investigación general; no se adopta un marco completo de diseño instruccional porque el objeto evaluado es el contenido del artefacto y no una intervención educativa.\par}
\vspace{5pt}
\needspace{5\baselineskip}\noindent\casilla\textbf{31. ¿Qué pasa si la prueba de campo que usted propone refuta la parte diferida de la hipótesis? ¿Qué queda de este trabajo?}\par
{\small\color{gris}\sffamily M1 metodologia\quad\chip{critica}{sin respuesta}\quad riesgo media}\par
{\small\rotulo{Qué falta}No requiere texto nuevo. Respuesta oral: quedarían en pie el artefacto verificado, la cobertura del contenido y los dos instrumentos reutilizables; lo que caería sería la atribución del efecto, y un resultado negativo sería informativo porque las celdas de la tabla de especificaciones indican qué instrumento revisar.\par}
\vspace{5pt}
\needspace{5\baselineskip}\noindent\casilla\textbf{32. ¿Cuál es el aporte a la ingeniería civil y no solo a la didáctica?}\par
{\small\color{gris}\sffamily M2 estructuralista\quad\chip{alta}{parcial}\quad riesgo media}\par
{\small\rotulo{Dónde responde hoy}Aportes del trabajo, folio 90: herramienta libre verificada, motor completo Q4/\allowbreak{}Q9, memoria de cálculo y dos instrumentos metodológicos reutilizables\par}
{\small\rotulo{Qué falta}Nombrar el aporte en clave disciplinar, en el texto o de palabra: un medio para que el egresado sepa auditar los resultados del software que usará en su ejercicio profesional, y un ejemplo documentado de verificación y validación aplicado a un programa de cálculo estructural, procedimiento que la práctica civil rara vez ve escrito.\par}
\vspace{5pt}
\needspace{5\baselineskip}\noindent\casilla\textbf{33. ¿Cómo maneja las unidades? Usted valida en kgf y centímetros; ¿qué pasa si el alumno mezcla sistemas?}\par
{\small\color{gris}\sffamily M2 estructuralista\quad\chip{alta}{parcial}\quad riesgo media}\par
{\small\rotulo{Dónde responde hoy}3.6, folio 68, que declara el sistema técnico del caso; 2.2 y Anexo B, manual de uso\par}
{\small\rotulo{Qué falta}Comprobar qué dice el manual sobre la coherencia de unidades y, si el programa no la impone, decirlo con la misma franqueza con que se declaran las demás limitaciones.\par}
\vspace{5pt}
\needspace{5\baselineskip}\noindent\casilla\textbf{34. ¿Por qué la solución del sistema y la recuperación de tensiones no tienen módulo propio, si son las dos etapas donde el alumno más se pierde?}\par
{\small\color{gris}\sffamily M3 numerico-pedagogico\quad\chip{verde}{respondida}\quad riesgo media}\par
{\small\rotulo{Dónde responde hoy}Cap. 3, alcances, folio 76, y Limitaciones, folio 91: ambas se exponen en el post-proceso —sonda con modos crudo y suavizado, vista tridimensional, tabla de resultados— y en la memoria de cálculo\par}
{\small\rotulo{Qué falta}Nada. De palabra: son las dos etapas cuya salida el alumno manipula directamente en el post-proceso, de modo que su exposición es la interfaz misma más el desarrollo numérico de la memoria.\par}
\vspace{5pt}
\needspace{5\baselineskip}\noindent\casilla\textbf{35. Usted dice que las herramientas comerciales son costosas, pero SAP2000, ANSYS y Abaqus tienen versiones académicas o de estudiante gratuitas. ¿Sigue siendo un argumento?}\par
{\small\color{gris}\sffamily M2 estructuralista\quad\chip{alta}{parcial}\quad riesgo media}\par
{\small\rotulo{Dónde responde hoy}Introducción, Planteamiento del problema, folio 2 («las herramientas comerciales suelen ser costosas, no estar diseñadas con un propósito didáctico explícito y ofrecer una localización al español limitada»)\par}
{\small\rotulo{Qué falta}Conceder el punto y reordenar el argumento, en el texto o de palabra: el motivo principal no es el costo sino la opacidad y la ausencia de propósito didáctico; las versiones académicas son gratuitas pero igual de opacas, están limitadas en tamaño de modelo y atadas a registro, plataforma y vigencia anual.\par}
\vspace{5pt}
\needspace{5\baselineskip}\noindent\casilla\textbf{36. ¿Por qué la viga de Timoshenko se resolvió solo con Q9 y no también con Q4, que es el elemento que un estudiante usará por defecto?}\par
{\small\color{gris}\sffamily M3 numerico-pedagogico\quad\chip{verde}{respondida}\quad riesgo baja}\par
{\small\rotulo{Dónde responde hoy}2.1.4, folio 43, y 3.6, folio 68: el caso es «un contraste puntual de exactitud sobre una malla fina y no un estudio de convergencia»; la nota de la tabla lo repite\par}
{\small\rotulo{Qué falta}Nada en el texto. De palabra, ofrecer el dato si se tiene a mano o comprometerse a incorporar la corrida Q4 del mismo caso: el comportamiento del Q4 bajo flexión ya está documentado en la membrana de Cook.\par}
\vspace{5pt}
\clearpage
\section{Lista de control al terminar}
Comprobaciones para hacer sobre el documento ya corregido, con el buscador del editor o del lector de PDF. Ninguna sustituye a una relectura del Resumen, la Introducción, 3.8 y las Conclusiones de un tirón.
\begin{itemize}[label={\casilla},leftmargin=1.8em,itemsep=3pt]
\item Buscar «mejor», «formar criterio», «comprend», «aprend», «potencial», «permite que», «anclar», «intuitiv»: cada aparición que quede es una propiedad del artefacto, un resultado de la literatura acotado a otra herramienta, o la hipótesis pendiente con su remisión a la prueba de campo.
\item Buscar «garantiz», «demuestr», «confirm», «libre de», «riguros», «no existe», «ningún», «único», «cualquier»: cada aparición está acotada a los casos ensayados o al conjunto revisado.
\item El objetivo general se lee igual en la Introducción, en la Tabla 2.2 y en las Conclusiones; la pregunta del problema, igual en el folio 2, en 2.1.3 y en las Conclusiones.
\item Cada uno de los criterios de 2.1.6 aparece en la tabla de veredictos de 3.8 con su umbral, su cifra medida y «se cumple»; ningún criterio nuevo aparece solo en la hipótesis o solo en 3.8.
\item Las cifras del Resumen, del Capítulo 3, de 3.8 y de las Conclusiones coinciden y llevan el mismo calificador: 0,04 \% y 0,26 \% (analítica), 0,21 \% en σx y 0,56 \% en desplazamientos (SAP2000), 2,9 \% y 4,56 \% en las componentes secundarias, declaradas fuera del criterio.
\item Un nombre por concepto: variable dependiente, estudiante, EduFEM o «el software», casos de referencia, canal de cálculo, módulo educativo, validador (o comprobador) de salud, post-proceso; «validación» en la acepción que fija 1.13.
\item La asignatura que define la población tiene nombre, sigla y semestre, y no queda ningún comentario DATO PENDIENTE en el fuente.
\item Todas las figuras y tablas se citan en el texto antes de aparecer (B.7–B.12 incluidas); los Anexos C y F se citan desde el cuerpo con un propósito; ninguna ecuación numerada queda sin que nadie la cite, o pierde el número.
\item Las 22 citas que sostienen una definición, cifra, umbral o atribución concreta llevan localizador; «[3, p. 1159; 3, p. 1171]\allowbreak{}» pasa a «[3, pp. 1159, 1171]\allowbreak{}»; las cinco claves del .bib tienen el año de su entrada.
\item Recompilar: 0 referencias indefinidas, 0 desbordes; índices regenerados; los folios 85 y 95 no quedan con tres líneas; el Anexo E tiene un párrafo de presentación que dice qué es y quién lo elaboró.
\end{itemize}
\end{document}